\chapter{财产犯罪}

% Source PDF physical page 252; see audit/source_coverage.json.
\begin{EditionOpeningPanel}
1.财产犯罪是分则最重要的板块。盗窃罪、抢劫罪、诈骗罪是重中之重。

2.罪名体系图（四阶图）

\InsertDiagram{figures/ch19_property.tex}
\end{EditionOpeningPanel}

\section{基本原理}

[问题]狗蛋偷了小芳的电瓶车。次日，小芳又从狗蛋家偷回来。小芳是否构成盗窃罪？（偷回电瓶车案）

\subsection*{一、保护法益}

就有形财物而言，财产犯罪的保护法益：合法占有，所有权，非法占有。

模型1：甲盗窃了乙的摩托车。次日，第三人丙从甲家盗窃了该摩托车。丙构成盗窃罪。

第一，丙没有侵犯甲的所有权，因为甲对摩托车没有所有权。甲对车仅存在非法占有事实。第二，非法占有事实若未达到平稳状态，则对其可以实施正当防卫，予以挽回。例如，甲将摩托车刚偷到手，便被群众发现，群众可以帮助主人抢回来。抢回来的行为属于正当防卫。第三，非法占有事实若达到平稳状态，则刑法予以保护，不允许这种类型的“黑吃黑”。若第三人丙盗窃该车，则构成盗窃罪。这是因为，刑法要保护稳定的财产秩序。财产权利要正常流转，需要财产秩序稳定。如果天下大乱，则财产权利无从谈起。

模型2：甲盗窃了乙的摩托车。次日，主人乙从甲处盗窃回来。

乙不构成盗窃罪。理由：对于甲的平稳的非法占有事实，刑法只是有限保护，也即第

% Source PDF physical page 253; see audit/source_coverage.json.
三人不能侵犯。非法占有事实可以对抗一般人，但不能对抗所有权人。所有权人恢复权利的行为没有非法占有他人财物的目的，不具有违法性。\footnote{少数说：乙构成盗窃罪。理由：所有权人应当通过法律程序救济权利，不能私力救济，刑法要保护稳定的财产秩序。}因此，本节开头的“偷回电瓶车案”中，小芳不构成盗窃罪。

模型3：甲借到乙的摩托车。未到期，主人乙从甲处盗窃回来。

乙构成盗窃罪。理由：合法占有能够对抗所有权人。所有权人偷回财物，虽然没有非法所有目的，但是有非法占有目的。成立盗窃罪，只要求非法占有目的，不要求非法所有目的。（2003年第10题、2008年第2题）\footnote{少数说：乙不构成盗窃罪。理由：所有权人偷回自己的财物，没有非法所有目的。成立盗窃罪，要求有非法所有目的。}

\InsertDiagram{figures/manual_table_p253_1.tex}

保护等级：合法占有\ensuremath{>}所有权\ensuremath{>}非法占有\ensuremath{>}一般人。

\subsection*{二、行为对象}

财产犯罪的行为对象主要是财物。

1.包括无体物，如电、通信网络信号。

2.包括虚拟财产，如Q币、游戏点卡。

3.包括违禁品。例如，甲看到乙非法持有毒品或假币，抢夺之，构成抢夺罪。（2010年第59题）

4.包括债权凭证，如存折、银行卡。

5.包括不动产，如房屋、房屋产权。

6.不包括人的身体。但是，如果从身体分离出来的器官、血液，属于财物。

\subsection*{三、非法占有目的}

取得型财产犯罪与毁弃型财产犯罪的区别在于：有无非法占有目的。如果有，则是取得型财产犯罪，如盗窃罪；如果没有，则是故意毁坏财物罪。

非法占有目的由排除意思和利用意思构成。排除意思，是指永久性排除占有人占有，将财物转为自己占有的意思；利用意思，是指对财物进行利用的意思。两种意思缺一不可。

1.排除意思

缺少排除意思，也即具有返还意思，不构成盗窃罪。例如，甲盗用乙的自行车去买酱油，然后放回原地。甲的盗用行为不构成盗窃罪。

% Source PDF physical page 254; see audit/source_coverage.json.
2.利用意思

（1）缺少利用意思，就可能构成故意毁坏财物罪。例1，甲偷出他人电动自行车，不是想用，而是毁弃，构成故意毁坏财物罪。例2（2009年第17题），甲殴打乙，乙拿出手机要报警，甲为防止报警，抢夺了手机，殴打后离去，走了5公里后，将手机扔进水沟里。甲有排除意思，但没有利用意思，甲不构成抢夺罪，构成故意毁坏财物罪。

（2）利用意思，不要求完全遵从财物的正常用途。

例1，某猥琐男基于怪癖盗窃女士内衣，也是利用，定盗窃罪。（2011年第61题）

例2，甲偷出图书馆的珍藏本，然后用来给自已垫桌腿，也是利用，定盗窃罪。

例3，乙盗取他人的钢材设备作为废品卖给废品收购站，也是利用，定盗窃罪。（2011年第61题）

例4，丙为了取暖，盗窃他人家具用来烤火，也是利用，定盗窃罪。（2011年第61题）

[巩固练习］题1，甲从超市偷出商品，然后又伪装退货，以此获取货款。这种偷出商品的行为，第一，虽然有返还意思，但这种返还意思不算数，因为这种返还是以获取货款为条件的。因此，甲仍有排除意思。第二，甲偷商品虽然不是为了自用，但仍具有利用意思。将商品作为等价物换取货款，就是一种利用方式。所以，甲具有非法占有目的，构成盗窃罪。甲冒领货款还构成诈骗罪。对盗窃罪与诈骗罪不需要并罚，按照重罪吸收轻罪处理。

题2，甲女与丈夫在打离婚官司，为了证明丈夫有外遇，偷走丈夫送给第三者的戒指，打算在诉讼后予以归还。甲女只有利用意思，没有排除意思，不构成盗窃罪。

题3，甲将乙杀害后，为了防止警方追查，将乙的行李及轿车推到山下。甲只是为了不让他人利用而将财物抛弃，不属于利用意思，而属于不让他人利用的意思。（2011年第61题）

题4，甲将乙杀害后，盗走乙的车，将户体运到外地，然后将车抛弃。甲既有利用意思，也有排除意思，构成盗窃罪。

题5，挪用公款罪缺少哪种意思？贪污罪有非法占有目的。挪用公款罪没有非法占有目的，具体而言，“挪”表示缺少排除意思，“用”表示具有利用意思。

3.为谁非法占有？

非法占有目的，既包括为行为人本人占有的目的，也包括为第三人占有的目的。

例1，甲爬到乙家的苹果树上，偷摘苹果，扔到隔壁丙家（自已暗恋对象）后院，让丙女拥有。甲属于盗窃行为。

例2，火车上，甲悄悄拿出熟睡的乘客乙的钱包，直接放进女友的包里，让女友拥有。甲构成盗窃罪。

例3，官员甲利用职权将公家的财物送给朋友乙。贪污罪要求有非法占有目的。甲构成贪污罪。

\section{侵占罪}

[问题】在婚礼现场，新娘委托伴娘保管一下戒指，等司仪喊“交换戒指”，由伴娘拿上台。伴娘答应，却悄悄拿走离开现场。伴娘是否构成侵占罪？（伴娘保管案）

第270条将代为保管的他人财物非法占为己有，数额较大，拒不退还的，处二年以下有期徒刑、拘役或者罚金；数额巨大或者有其他严重情节的，处二年以上五年以下有期

% Source PDF physical page 255; see audit/source_coverage.json.
徒刑，并处罚金。

将他人的遗忘物或者埋藏物非法占为己有，数额较大，拒不交出的，依照前款的规定处罚。

本条罪，告诉的才处理。

本罪没有侵害他人的占有权，而只侵害了他人的所有权，因为行为人事先已经占有他人财物。例如，乙将汽车交给甲保管，约定一周后来取。一周后乙来取，甲拒不返还。甲构成侵占罪。由于本罪是告诉才处理的犯罪，因此真正定罪的被告人很少。

\InsertDiagram{figures/manual_p255_1.tex}

\subsection*{一、行为对象}

本罪的行为对象：他人所有、自己占有的财物。对一件物品，所有权和占有权可以相分离。所有权人对物品虽有所有权，但由于某些原因，占有权可能又归其他人拥有。这样的财物主要包括两大类：一是代为保管物，二是遗忘物、埋藏物。这主要是根据主人失去占有的原因来划分。前者是主人基于自愿而将财物转移占有给行为人。后者是主人非基于自愿而脱离占有。

\subsubsection*{（一）代为保管物}

所有权人自愿将财物委托给行为人占有。例如，代为保管、暂时出借、用于担保等。因此，对本罪条文中的“代为保管的财物”应作扩大解释，是指委托保管物，包括出借物、担保物等。

1.本罪的“代为保管的财物”必须是行为人占有的财物。行为人的有些保管、借用行为并没有形成对财物的占有，不会构成侵占罪。

例1，在本节开头“伴娘保管案”中，伴娘只是在现场临时保管，戒指仍在新娘实际控制范围内，仍由新娘占有，伴娘的保管没有形成对戒指的占有，只是新娘占有戒指的辅助者。伴娘将新娘占有的财物转移为自己占有，构成盗窃罪，而非侵占罪。

例2，甲乙在餐厅吃饭，乙借用甲的手机打个电话，打完后悄悄拿走手机溜掉。乙的借用没有形成对手机的占有。乙拿走手机属于将甲占有的财物转移为自己占有，构成盗窃罪。

[巩固练习]题1（2011年第62题），甲委托乙帮自己照看房子。乙将房子里的电脑擅自变卖。乙对电脑没有形成占有，构成盗窃罪。

题2，甲将车交给洗车店的乙，约定次日来取车。乙在洗车时发现车内有张彩票，拿走兑奖，中奖100万元。甲将车委托给乙洗车，乙占有车。但甲没有委托乙占有车内物品，乙没占有车内的物品（包括彩票）。彩票是有可能价值数额较大甚至巨大的财物。乙将甲占有的彩票，转移为自己占有，构成盗窃罪。\footnote{盗窃罪的既遂数额不能以100万元认定，因为乙在盗窃时，不知道该彩票能够中奖、价值100万元。盗窃罪的既遂数额可认定为数额较大。乙中奖100万元后，甲要求乙归还100万元，乙不归还，构成侵占罪。}（2023年试题）

% Source PDF physical page 256; see audit/source_coverage.json.
2.基于不法原因而委托管理的财物

（1）委托给付的财物。例如，甲欲向国家工作人员行贿，将财物委托给乙，让乙转交给国家工作人员，乙却将财物据为已有，乙是否构成侵占罪？（2017年第18题）

肯定说（少数说）的理由：甲有返还请求权，乙拒不返还，构成侵占罪。

否定说（多数说）的理由：①对于贿赂款，国家应予以没收，甲丧失所有权，没有返还请求权，因此，乙拒不返还，不构成侵占罪。②如果认定为侵占罪，会得出民法上甲没有返还请求权，但刑法上认为其有返还请求权的结论，刑法和民法对相同问题会得出不同结论，法秩序的统一性会受到破坏。\footnote{有人认为，根据《民法典》第462条，无权占有也能产生返还请求权。然而，第462条指的是，对财物无权占有时，财物被侵夺了，无权占有者对财物有返还请求权。例如，甲无权占有一辆摩托车，该车被乙侵夺，甲对乙有返还请求权。如果甲基于自愿且基于不法原因，将财物转移给乙占有后，则甲不能根据第462条主张返还请求权。}

（2）委托保管的财物。例如，甲将其盗窃的财物委托给乙，告知是偷来的，让乙窝藏（或代为销售）。乙窝藏后，并且拒不返还。乙构成掩饰、隐瞒犯罪所得罪之外，是否还构成侵占罪？

肯定说（少数说）的理由：甲有返还请求权，乙拒不返还，构成侵占罪。

否定说（多数说）的理由：甲是基于不法原因将赃物交给乙保管，没有返还请求权。因此，乙拒不返还，不构成侵占罪。这表明，这种类型的“黑吃黑”不构成犯罪。

[注意】不是基于不法原因委托保管的财物，而是基于合法原因委托保管的财物，是侵占罪的保护对象。例如，甲盗窃了一辆摩托车，对修理师傅乙称：“请将我的摩托车修理好，我明天来取。”乙察觉出该车是甲盗窃来的，第二天拒不归还。由于甲是基于合法原因（修车）而委托乙保管，因此甲享有返还请求权，乙拒不返还，构成侵占罪。这表明，这种类型的“黑吃黑”构成犯罪。

\subsubsection*{（二）遗忘物、埋藏物}

1.遗忘物

例如，甲将钱包遗忘在街边，乙捡到，甲索要，乙拒不返还。当乙捡到时，此时钱包属于甲所有、乙占有的财物。此时乙并不构成侵占罪。当乙拒不返还，表明变成自己所有，则构成侵占罪。

[注意］“遗忘物”应扩大解释为“非自愿而脱离占有的财物”。例如，甲不在家，甲家阳台的衣服掉落到楼下邻居乙家。甲向乙索要，乙拒不返还。当乙捡到衣服时，此时衣服属于甲所有、乙占有的财物。当乙拒不返还，表明变成自已所有，则构成侵占罪。

2.埋藏物

例如，狗蛋继承了爷爷一座房产，爷爷去世时没有告诉狗蛋，院子里埋了一坛金子。狗蛋将该房产卖给小芳，小芳发现这坛金子。狗蛋知道后，向小芳索要，小芳拒不返还。当小芳发现这坛金子时，此时金子属于狗蛋所有、小芳占有的财物。当小芳拒不返还，表明变成自己所有，则构成侵占罪。

\subsection*{二、行为方式}

行为方式：行为人将占有变成自己所有（据为己有）。变成自己所有的方式可分为作

% Source PDF physical page 257; see audit/source_coverage.json.
为方式和不作为方式。作为方式，如无权处分，例如，变卖掉、消费掉、抵押掉、赠与他人。不作为方式，如拒不返还。\footnote{“拒不返还”包括两种情形：一是主人索要，行为人拒不返还；二是主人尚未索要，但是能够推定，即使主人索要，行为人也会拒不返还。}

例1，甲委托乙保管自已的小狗，乙保管了两天就将小狗给卖掉了。乙的行为属于将自己占有变成自已所有，构成侵占罪。又如，乙将小狗给下锅炖了、吃了，也属于将自己占有变成自已所有，也构成侵占罪。乙不构成故意毁坏财物罪，因为有利用意思。再如，乙没卖也没吃小狗，但就是不还给甲，这也属于将自己占有变成自已所有，也构成侵占罪。

例2，甲委托乙保管现金1万元，保管5天。第3天，乙用这1万元给自己买了台电脑。此时很难说明，乙将甲的现金据为己有，因为现金是种类物，放在一起无法进行区分，产生混同效果，只要第6天乙如数归还，就说明乙没有将1万元变成自已所有。如果第6天拒不返还，就表明乙将1万元变成自己所有。此时，拒不返还就具有了“变成自己所有”的性质。

[巩固练习]题1（错误找钱），甲给客户乙多找了钱，乙拿到手后发现多找了1万元。由于资金是种类物，有混同效果，乙花掉这笔钱，不构成侵占罪。但乙拒不返还，构成侵占罪。

[提示]对于资金，民法有“占有即所有”的规则。但是，刑法不能承认该规则。这是因为，侵占罪的成立，以承认占有与所有相分离为前提。若承认“占有即所有”规则，则意味着上文例2、题1中的乙对资金已经有了所有权，那么就不存在侵占罪要求的“变占有为所有”，就都不构成侵占罪了。刑法在此采取实质标准，也即资金的实际所有权人是谁，不能认为“资金只要到了谁手里，谁就有所有权”。

题2（2021年主观题），甲敲诈勒索乙，让乙将现金放进街边垃圾桶。乙被迫照办。甲告知丙真相，让丙帮助取回现金。丙取回现金，交给甲，甲丙平分了现金。假设，乙将现金放进垃圾桶时，甲便构成敲诈勒索罪既遂，那么丙的行为是否构成侵占罪？构成。丙从垃圾桶拿到现金，便现实占有了现金，然后交给甲进行平分，属于将乙所有、自己占有的现金变成自己和甲所有，构成侵占罪。丙同时构成掩饰、隐瞒犯罪所得罪，与侵占罪想象竞合，择一重罪论处。

\section{盗窃罪}

[问题]甲骑自行车在前面，让三轮车夫乙载着自己的货物跟着自己走。走到十字路口，乙看到人流挺多，便掉转车头，甲发现后呼喊，但因为人群间隔，眼睁睁看着乙拉走自己的货物。乙构成盗窃罪还是侵占罪？（三轮车夫案）

盗窃罪（第264条），是指带着非法占有目的，通过平和手段，将他人占有的财物转移为自己占有。转移占有由两个环节组成：一是破坏他人对财物的占有。有该行为，就成立盗窃罪。二是建立自己对财物的占有。有该行为，盗窃罪就既遂了。转移占有的手段必须是平和手段。所谓平和手段，是指手段不能对人身具有暴力、胁迫的性质。

% Source PDF physical page 258; see audit/source_coverage.json.
\InsertDiagram{figures/manual_p258_1.tex}

\subsection*{一、实行行为的第一项条件“转移占有”}

考试主要通过盗窃罪与侵占罪的区分来考查盗窃罪的“转移占有”。

\InsertDiagram{figures/manual_table_p258_2.tex}

\CJKunderwave{关键区分}：谁在占有财物？如果是主人在占有，行为人破坏主人的占有就是盗窃；如果行为人事先在占有，行为人就是侵占。如何判断谁在占有？刑法上的占有分为事实上的占有和观念上的占有。

\subsubsection*{（一）事实上的占有}

这是指财物在主人实际控制范围之内。

1.\CJKunderwave{空间距离问题}。主人占有财物，不要求随身携带，只要财物在主人实际控制范围内即可。主人占有财物，可以与财物保持一定距离，也即主人占有财物，既可以是紧密占有，也可以是松弛占有。（2008年第16题，2010年第62题，2011年第62题，2016年第59题)

例1，在本节开头“三轮车夫案”中，注意不要先判断行为人乙是否占有财物。若如此判断，容易陷入错误思维：乙在占有财物，变成自己所有，构成侵占罪，属于侵占代为保管物。\CJKunderwave{正确的判断顺序是}：先判断主人是否在占有，然后判断行为人是否在占有。\CJKunderwave{一旦}\CJKunderwave{确认主人在占有，那么即使行为人在实际的持有，也视为对主人占有的一种辅助手段，只}\CJKunderwave{是占有的辅助者}。本案中，正确的思维是：主人在占有财物，因为财物在主人实际控制范围内，乙只是占有的辅助者。乙将主人占有的财物转移为自已占有，构成盗窃罪。

[总结模型］主人占有自已的财物，但距离财物有一定距离，行为人又持有该财物，然后将该财物转移为自己占有。行为人构成盗窃罪。躲坑提示：行为人不构成侵占罪，不属于侵占代为保管物。注意：不能看行为人是否“保管”财物，要看行为人是否占有财物。

\InsertDiagram{figures/manual_p258_3.tex}

例2，旅客放在火车行李架上的行李，属于旅客在占有。

例3，秘书与领导同行，帮领导拎包，包仍属于领导在占有，秘书拎包只是对领导占有的一种辅助手段。（2016年第59题）

例4，商场里，顾客在试穿衣服，衣服仍属于销售者在占有。

例5，大学食堂，学生将书包放在座位上占座，然后去窗口买菜。书包仍属于主人在占有。（2011年第62题，2014年第60题）

% Source PDF physical page 259; see audit/source_coverage.json.
[注意】刑法上的占有与民法上的占有不同，前者侧重于事实状态，后者侧重于权利状态。刑法不承认“间接占有”概念，民法承认“间接占有”概念。

2.\CJKunderwave{短暂遗忘问题}。财物处于主人实际控制范围内，主人的短暂遗忘不改变主人的占有事实。躲坑提示：行为人不构成侵占遗忘物型侵占罪。例如，乙离开公园长椅后，将提包遗忘在长椅上，甲见此情景，在乙走了十米左右，马上走过去拿走提包。由于乙尚未走远，提包尚在乙的实际控制范围内，所以乙仍在占有提包。即使乙短暂遗忘，也不能改变乙对提包的占有。甲构成盗窃罪，而非侵占遗忘物型侵占罪。

3.\CJKunderwave{占有的转化问题}。特定场所的财物，主人丧失占有，便转化为场所管理者占有。此时第三人将财物变成自己占有，构成盗窃罪。

例如，乘客甲将钱包遗忘在出租车上，下了出租车，走了。乘客乙上车后，发现钱包，下车时将钱包拿走。甲虽然对钱包失去占有，但是钱包转化为司机占有（不以司机发现钱包为条件）。乙拿走钱包，属于将司机占有的财物转移为自己占有，构成盗窃罪。在此，所有与占有相分离。盗窃罪只要求转移占有，不要求转移所有权。

\CJKunderwave{占有转化的判断标准}：场所管理人对场所里的财物的管控力大小。如果管控力大，则能转化为管理人占有。判断管控力大小，主要考察人员流动性、管理人的注意力范围、场所的封闭性等因素，进行综合判断。

[巩固练习］题1，遗忘在飞机上、高铁上的财物，可以转化为乘务人员占有（多数说）。但遗忘在普通火车或公交车上的财物，就不能进行这种转化，只能视为无人占有的遗忘物。（2011年第62题，2012年第18题）

题2，客人甲离开宾馆房间时，将钱包忘在房间，负责打扫的服务员看到后据为已有。虽然甲失去对钱包的占有，但是转化为宾馆管理者在占有，服务员的行为仍属于将他人占有的财物变成自己占有，构成盗窃罪。（2010年第62题）

题3，顾客将钱包遗忘在大型超市、热闹的较大空间的餐厅，不能转化为场所管理人占有。

题4，储户遗忘在储蓄所柜台上的现金，转化为储蓄所管理者占有。乘客遗忘在地铁、机场安检台上的行李，转化为现场管理者占有。

[引申］盗窃罪只要求转移占有，不要求转移所有权，因此，盗窃罪的被害人是财物的占有人，不要求是财物的所有权人。成立盗窃罪，要求违背被害人意愿而非法转移占有。这里的被害人是财物的占有人，不要求是财物的所有权人。

例如，甲将自已的自行车（价值5000元）出借给乙，并约定“乙如果遗失则按3倍价格赔偿”。甲很期待车被偷。几日后丙从乙处偷走该车。甲得知车被偷后，非常开心。虽然丙没有违背所有权人甲的意愿，但是违背了占有人乙的意愿，因此构成盗窃罪。（2022年试题）

4.\CJKunderwave{上下占有问题}。范例：小卖部老板甲让外甥乙帮忙一起卖货，乙将店里财物拿回家。如果认为乙占有财物，那么乙属于侵占。

\CJKunderwave{结论}：上级所有的财物，下级有无占有，关键看上级对下级的授权大小。如果上级对下级比较信赖，充分授权，下级享有独立处分权，则表明下级在占有财物。上例中，甲让乙一起卖货，表明甲也在店里，乙只是辅助者，乙没有占有财物，构成盗窃。如果甲将一家分店交给乙经营，乙将店里财物拿回家，乙属于侵占，由于利用了职务便利，构成职务侵占罪。

% Source PDF physical page 260; see audit/source_coverage.json.
例如，老板甲让司机乙运输货物到外地，让秘书丙押运。在运输途中，乙没有占有财物，只是占有的辅助者。如果拿走财物，属于盗窃。如果甲信任乙，让乙一个人长途运输货物，则乙在占有货物，属于侵占。

5.\CJKunderwave{共同占有问题}。例如，甲乙各出资一半，共同购买一台电脑。甲趁乙不在，将电脑变卖。电脑是甲乙共同占有的财物。甲侵犯了乙的占有，构成盗窃罪。

6.\CJKunderwave{封缄物的占有}。范例：甲将一个行李箱交给乙，让乙拿回家保管。乙在保管时，打开行李箱，取走里面的珍贵珠宝。对乙该如何处理？\CJKunderwave{结论}：主人甲对内容物（珠宝）仍占有，受托人乙对封物的整体在占有。所以，乙打开封物，取走内容物，转移为自己占有，构成盗窃罪。乙将封物整体据为已有，不打开，也不归还，构成侵占罪。如果不打开，将封物整体卖给丙，则构成盗窃罪，因为该行为中包含了“将甲占有的内容物转移给第三人占有”。

[巩固练习]题1，快递员将负责运送的包裹打开，拿走内容物，构成盗窃罪。内容物由主人占有，快递员及其公司没有占有。

题2（2024年试题），乙将电脑拿到修理店，约定次日来取。甲修电脑时，发现主机内藏着100克金条，拿出来放入口袋。次日，乙想起金条，发现金条没了，要求甲归还。甲归还。甲构成盗窃罪既遂，而非中止。

\subsubsection*{（二）观念上的占有}

这是指虽然财物处于主人实力控制范围外，但是根据社会一般观念，可以推知有人在占有该财物。

例1，甲将自行车停在家门前，去外地出差。身在外地，但对自行车仍在占有。有人偷走自行车，定盗窃罪。

例2，主人为了躲避地震，将财物搬放在门前大街上，然后去防震棚。财物仍属于主人在占有。

例3，信鸽爱好者放飞的信鸽，仍属于主人占有。

[巩固练习］题1，货主的轮船遇海难，沉海底。轮船内的货物仍属于货主占有。货主的卡车遇事故翻车，货物撒落路边。货物仍属于货主占有。

题2（2019年试题），甲公司将共享单车投放在街边，该单车在无人使用时，属于甲公司占有，即使车锁坏了或没有上锁，也属于甲公司占有；\footnote{锁是用来防小偷的，不是主人占有自己财物的必要条件。例如，主人携带的行李箱忘记上锁，不能因此改变主人对行李箱的占有；主人将没有上锁的自行车放在自家楼下，主人依然在占有自己的自行车。世界上本来没有锁，偷的人多了，才有了锁。}该单车由乘客使用时，由乘客占有。情形一，乙将该单车的锁拆掉，将单车放在自家院子，属于将甲公司占有的财物转移为自己占有，成立盗窃。情形二，乙正常使用完共享单车后，将车停在楼下，方便自已下次使用，有时也被其他乘客提前使用。乙属于正常使用该单车，不成立盗窃。情形三，乙将市区的共享单车偷偷搬到偏远农村，供村民扫码使用、正常使用。乙没有将该车转移为自己占有，不成立盗窃。

\subsubsection*{（三）占有意思}

无论是事实上的占有，还是观念上的占有，占有人对财物都要有占有意思。如果行为

% Source PDF physical page 261; see audit/source_coverage.json.
人只是客观上占有财物，但主观上没有占有意思，则不属于在占有财物。

例如，甲走路时脚底粘上一张百元钞票，但甲不知情，带着钞票向前走。由于甲对钞票没有占有意思，不算在占有钞票。

1.当占有事实有些松弛时，占有意思对占有事实可以起到补充作用。

例如，甲在五楼阳台上不慎将手机掉到地面，行人乙看到后便要捡走，甲呼喊住手，乙不理睬捡起就跑。甲虽然对手机不是紧密占有，但甲具有明确的、强烈的占有意思，仍视为甲在占有手机。乙构成盗窃罪，而非侵占罪。该案例说明盗窃罪不要求秘密性。这是貌似侵占、实为盗窃的情形。

2.占有意思不要求特别声明，可以推定存在。

例1，主人外出，快递员将包裹放在主人家门口。主人虽然不知道门口有包裹，但在占有包裹。

例2，继承祖上房屋，不知墙体内有金条，装修工发现后据为已有，构成盗窃。

例3，甲以伤害故意将乙打晕，临时发现乙身上钱包，然后拿走。虽然乙昏迷，但仍在占有钱包。甲构成故意伤害罪和盗窃罪。

3.死者占有问题。

（1）当场杀，杀人者拿。例如，甲基于报仇目的杀死乙。第一种情形，甲当场发现乙身上的财物，当场拿走。多数说认为，甲构成盗窃罪，因为甲故意制造死亡结果，并很快利用该死亡状态，此时死者对财物的占有状态仍在短时间延续，没有消失。\footnote{根据2005年6月8日最高人民法院《关于审理抢劫、抢夺刑事案件适用法律若干问题的意见》第8条，实施故意杀人罪后，临时起意拿走他人财物的，定盗窃罪。}少数说认为，甲构成侵占罪，理由是，死人对财物没有占有意思，因此失去占有。第二种情形，甲离开现场，过后又返回现场，才发现死者身边的财物，拿走财物。唯一结论：甲构成侵占罪，因为死者对财物的占有状态已经消失（\CJKunderwave{总结}：当场杀，当场拿，多数说定盗窃；当场杀，过后拿，定侵占）。

（2）第三人拿。例如，甲杀死乙后离开现场，无关的丙路过，发现死者身上的财物并拿走。对于路过的第三人而言，无须承认死者的占有。第三人任何时候路过拿走死者财物，均构成侵占罪（\CJKunderwave{总结}：第三人，任何时候拿，都定侵占）。

（3）家里拿。例如，甲基于报仇目的在乙住宅内杀死乙。无论甲当场发现财物，拿走财物，还是过后发现财物，拿走财物，均构成盗窃罪。第三人任何时候入户拿走财物，均构成盗窃罪。理由是，财物在乙家中，应当承认死者对财物的占有（\CJKunderwave{总结}：家里杀，家里拿，都定盗窃)。

\subsection*{二、实行行为的第二项条件“平和手段”（秘密性问题）}

转移占有的手段必须是平和手段，不能是暴力手段。

问题是，在平和手段的基础上，是否还要求手段具有秘密性？

\subsubsection*{（一）观点展示}

个别公开说认为，盗窃罪的绝大多数情形是秘密的，个别情形是公开的。全部秘密说

% Source PDF physical page 262; see audit/source_coverage.json.
认为，盗窃罪的所有情形都是秘密的。\footnote{张明楷教授主张，盗窃罪不要求秘密性，可以公开盗窃。参见张明楷：《盗窃与抢夺的界限》，载《法学家》2006年第2期。}

个别公开说认为，\CJKunderwave{盗窃罪、抢夺罪、抢劫罪的区别在于手段的暴力程度：（1）盗窃}\CJKunderwave{罪，是对物平和手段，对人没有危险。（2）抢夺罪，是对物暴力（如此才能体现“抢”}\CJKunderwave{的含义），夺取财物，对人有危险。（3）抢劫罪，是对人暴力，压制人的反抗}。三罪是性质相同、程度不同的关系，是A与A+B的位阶关系（包容评价关系）。抢劫罪可以包容评价为抢夺罪、盗窃罪。

全部秘密说认为，盗窃罪与抢夺罪、抢劫罪的区分在于，盗窃罪是秘密的，抢夺罪、抢劫罪是公开的。盗窃罪与后两罪是A与非A的对立排斥关系，抢劫罪、抢夺罪无法包容评价为盗窃罪。全部秘密说认为，\CJKunderwave{盗窃罪的秘密性，是指犯罪人自认为自己的手段是秘}\CJKunderwave{密的}。

例1（\CJKunderwave{抢劫乌龙案}），乙站在街边抽烟，旁边停放一辆摩托车，不是乙的，是丙的，丙去办事，乙不认识丙。甲看到该摩托车，以为是乙的车，想公开抢劫该车，便上前将乙打倒，导致乙轻伤，然后骑走该车。甲对乙构成故意伤害罪，对丙构成盗窃罪，并罚。第一，甲不构成抢劫罪，因为抢劫罪的暴力对象要求是财物的主人、占有人等，不能是无关的第三人。第二，甲不构成抢夺罪，因为甲没有对物暴力，对车主的人身没有危险，不属于夺取车主的财物。第三，按照全部秘密说，秘密性是指犯罪人自认为自己的手段是秘密的，那么甲自认为自已是公开手段，甲便不构成盗窃罪。甲对丙的车便不构成犯罪，由此导致处罚漏洞。按照个别公开说，抢劫罪的故意可以包容评价为盗窃罪的故意，所以甲有盗窃罪的故意，有盗窃罪的行为，构成盗窃罪。

例2（保姆搬东西案），主人中风，失去说话能力，瘫痪在床，睡着了。保姆将主人家的财物往外搬，搬的过程中，主人醒了，只能干瞪眼看着。保姆看到主人醒来，没搭理，继续搬。主人又睡着了，保姆仍继续搬。依据全部秘密说，主人醒来之前，保姆构成盗窃罪；主人醒来之后，保姆构成抢夺罪；主人又睡着了，保姆又构成盗窃罪。这会导致，主人的眼晴一闭一睁，保姆的同一个行为，一会儿构成盗窃罪，一会儿构成抢夺罪。根据个别公开说，保姆构成盗窃罪。（2021年试题）\footnote{源自实务案例，参见朱卿新：《利用被害人无反抗能力取走财物的行为应如何定性》。案情是，李某乘同村村民陈某家仅有身患重病不能行走的陈某雄一人在家之机，在见到陈某雄坐在屋厅的轮椅上时，先后两次径直从大门进人屋内，使用一把螺丝刀撬开衣柜、写字台抽屉等处，盗窃现金人民币3624元。文章认为，盗窃罪中的窃取，不是以实施隐秘方式实施为必要条件。信宜市人民法院于2014年3月3日以李某犯盗窃罪判处有期徒刑8个月，并处罚金500元。访问网址为htps：//www.maoming.jcy.gov.cn/lyj/yasf/202102/t20210223\_3137632.shtml}

例3（楼上掉东西案），甲在三楼阳台，不慎将钱包掉到一楼地面。路人乙要捡走，甲不让乙捡，乙不予理睬，仍捡走。根据全部秘密说，乙构成抢夺罪。然而，乙是捡走，不是抢夺走。根据个别公开说，乙构成盗窃罪，属于极少数的公开盗窃。

\subsubsection*{（二）真题答案支持个别公开说}

由于全部秘密说和个别公开说争议日益激烈，自2026年以后，法考大概率不会问“唯一正确答案”。如果问“唯一正确答案”，则应选择个别公开说。这是因为，法考官方真题答案采取个别公开说。在“九大本”法考教材与官方答案之间，官方答案的效力更高，是终极

% Source PDF physical page 263; see audit/source_coverage.json.
选择依据。

1.真题1（\CJKunderwave{2017年第60题}）

下列说法是否正确：“甲欲进王某家盗窃，正撬门时，路人李某经过。甲误以为李某是王某，会阻止自已盗窃，将李某打昏，再从王某家窃走财物。甲不构成抢劫既遂。”

第一，甲误以为李某是主人王某，为了继续取财，对其实施暴力，表明甲的意图从盗窃升级为公然抢劫，意图抢劫主人的财物。

第二，甲不构成抢劫罪，因为抢劫罪的暴力对象要求是财物的主人等，李某是路人，不是主人，不是抢劫罪的合格对象。

第三，甲不构成抢夺罪，因为甲没有对物暴力，对主人的人身没有危险，不属于夺取主人的财物。

第四，甲对李某构成故意伤害罪（假如李某有轻伤），对王某应当构成盗窃罪，两罪并罚。甲构成盗窃罪的理由是：（1）客观上，甲实施了盗窃王某财物的行为。（2）主观上，甲具有抢劫的故意。抢劫的故意可以包容评价为盗窃的故意，因此甲具有盗窃的故意。所以甲构成盗窃罪。

第五，\CJKunderwave{全部秘密说认为，秘密性是指盗窃者自认为自已是秘密的。根据这种看法，甲}\CJKunderwave{自以为自己在对主人实施暴力，在公然抢劫，甲便不具有秘密盗窃的故意，那么甲就不满}\CJKunderwave{足“自认为自己是秘密的”这一要求，因此不构成盗窃罪}。这会导致对甲的取财行为无法定罪，出现处罚漏洞。

本题官方答案认为，甲对王某的财物不构成抢劫罪既遂，那么只能构成盗窃罪既遂。这足以证明，甲成立盗窃罪，不要求“自认为自已是秘密的”。\footnote{参见张明楷：《刑法学》（第五版），第972页脚注。}全部秘密说也赞成，甲构成盗窃罪既遂。但赞成该结论，就必须放弃“盗窃罪要求秘密性，也即自认为自已是秘密的”，否则会自相矛盾。

2.真题2（2011年第55题）

下述说法是否正确：“乙为吴某入户盗窃望风，但吴某入户后实施抢劫行为。乙是盗窃罪的共犯”。官方答案认为，该说法正确。

官方答案解析认为，“乙是盗窃罪的帮助犯，吴某是盗窃罪的实行犯。由于吴某又构成抢劫罪，最终对吴某以抢劫罪论处”。其意思是，根据共犯从属性，要认定帮助者乙构成盗窃罪的帮助犯，必须认定一个前提：实行犯吴某构成盗窃罪的实行犯。吴某构成抢劫罪，抢劫罪能够包容评价为盗窃罪，所以吴某构成盗窃罪。这个盗窃罪显然是公开进行的。

如果认为盗窃罪必须是秘密的，那么吴某便不构成盗窃罪，由此帮助者乙便无法构成盗窃罪的共犯。

有些同学认为，既然乙主观上有帮助盗窃的故意，客观上有帮助行为，所以当然构成盗窃罪的帮助犯，不需要考察吴某是否构成盗窃罪。这种看法忽略了共犯从属性的要求。如果实行犯不构成盗窃罪（既遂），则帮助犯不可能构成盗窃罪的帮助犯（既遂）。（2026年“九大本”法考教材刑法卷第59页提醒了这一点）

% Source PDF physical page 264; see audit/source_coverage.json.
3.真题3（2013年第60题）

甲潜入他人房间欲盗窃，忽见床上坐起一老，哀求其不要拿她的东西。甲不理睬而继续翻找，拿走一条银项链（价值400元）。关于本案的分析，下列哪些选项是正确的？

A.甲并未采取足以压制老姬反抗的方法取得财物，不构成抢劫罪

B.如认为区分盗窃罪与抢夺罪的关键在于是秘密取得财物还是公然取得财物，则甲的行为属于抢夺行为；如甲作案时携带了凶器，则对甲应以抢劫罪论处

C.如采取B选项的观点，因甲作案时未携带凶器，也未秘密窃取财物，又不符合抢夺罪“数额较大”的要件，无法以侵犯财产罪追究甲的刑事责任

D.如认为盗窃行为并不限于秘密窃取，则甲的行为属于入户盗窃，可按盗窃罪追究甲的刑事责任

官方答案：ABCD。虽然考查的是观点展示，但是官方的立场倾向很明显。AB项想揭示，如果将公开盗窃认定为抢夺，则携带凶器公开盗窃就属于携带凶器抢夺，就需要定处罚很重的抢劫罪，由此导致罪刑不相适应。CD项想揭示，如果要求盗窃罪具有秘密性，就导致甲不构成任何财产犯罪，由此出现处罚漏洞。这显然不合理。

4.真题4（2014年第19题）

乙购物后，将购物小票随手扔在超市门口。甲捡到小票，立即拦住乙说：“你怎么把我购买的东西拿走？”乙莫名其妙，甲便向乙出示小票，两人发生争执。适逢交警丙路过，乙请丙判断是非，丙让乙将商品还给甲，有口难辩的乙只好照办。关于本案的分析（不考虑数额），下列哪一选项是错误的？

A.如认为交警丙没有处分权限，则甲的行为不成立诈骗罪

B.如认为盗窃必须表现为秘密窃取，则甲的行为不成立盗窃罪

C.如认为抢夺必须表现为乘人不备公然夺取，则甲的行为不成立抢夺罪

D.甲虽未实施恐吓行为，但如乙心生恐惧而交出商品的，甲的行为构成敲诈勒索罪

官方答案：D。本题考查的是三角诈骗与盗窃罪的间接正犯的区别，区别在于受骗人有无处分受害人财物的权限。如果有，则为三角诈骗；如果没有，则为盗窃罪的间接正犯（对此诈骗罪中有详细讲解）。本题考查的是观点展示。但法考界官方有明显倾向：\CJKunderwave{交警不}\CJKunderwave{是法官、不是片区民警，对这件事没有管辖权，没有权力处分乙的财物，所以甲不构成三}\CJKunderwave{角诈骗，而构成盗窃罪的间接正犯}。如果要求盗窃罪具有秘密性，则甲不构成盗窃罪。甲也不构成三角诈骗。对甲便无法定罪，由此出现处罚漏洞。

5.真题5（2021年试题）

小学班主任在微信家长群中通知大家交学生的生活费。家长们纷纷用发红包的方式交纳。某个奇葩家长甲竟然在群里抢收了这些红包，数额较大。问：甲是否构成盗窃罪？这是真实案例，司法机关对其以盗窃罪论处。甲意识到自已是在微信群这种公众群里抢红包，自认为自己是公然进行的。成立盗窃罪，如果要求盗窃者必须自认为自己是秘密的，

% Source PDF physical page 265; see audit/source_coverage.json.
则对甲无法定盗窃罪，会出现处罚漏洞。\footnote{有观点认为，甲在抢红包时，发红包者可能没看手机界面，所以甲的行为仍然是秘密性的。然而，完全有可能，发红包者在看着手机界面时，甲抢走红包，此时甲的行为仍然是公开的。而且，无论哪种情形，甲在主观上都认为自已是在公众群里抢红包，不是秘密进行的，而是公开进行的。在微信群里抢红包，之所以用“抢”字，就是要表达手段的公然性。如果甲非法登录乙的支付宝账户，划走资金，则甲主观上认为自己在秘密操作。甲在公众群里操作，甲主观上肯定认为自己的手段是公然的。所以，按照全部秘密说，对甲不能定盗窃罪。这会出现处罚漏洞。本题参见张明楷编著：《刑法的私垫（之四）》，北京大学出版社2025年版，第555-556页。作者认为甲构成公开盗窃。}

6.“自认为自己是秘密的”的问题所在

第一，这种\CJKunderwave{秘密性的含义违反了主客观相一致原则的一个子原则“客观决定主观原}\CJKunderwave{则”}。该原则是指，客观构成要件决定了主观认识内容。凡不是客观构成要件的要素，就不是主观认识内容。该原则也被称为“客观构成要件对故意内容的规制功能”。全部秘密说也承认，成立盗窃罪，客观上可以是公开的，也即被财物主人、周围人看见，都可以定盗窃罪。既然客观上不要求秘密，那么主观上不应要求行为人“自认为是秘密的”。

第二，“自认为自己是秘密的”这种观点会成为犯罪人随意狡辩的理由，因为“自认为自己是秘密的”是犯罪人的心理活动，无法对此予以证明。

例1，上述“保姆搬东西案”中，保姆发现主人看见她搬东西，知道自已是公开进行的。假如她认为盗窃罪可能处罚轻，就向法官辩解说：“我自认为我是秘密的。”如此，法官就必须判她构成盗窃罪。她如果认为抢夺罪可能处罚轻，就辩解说：“我自认为我是公开的。”法官就必须判她构成抢夺罪。若如此，则法官判她构成什么罪，完全取决于犯罪人如何描述自己的心理活动。这显然是不合理的。

例2，甲为了非法占有，偷偷给乙下迷药，乙昏迷后，甲拿走财物。根据全部秘密说，甲自认为自己是秘密的，因此对甲应定盗窃罪。但是，甲显然更应构成抢劫罪，属于昏醉抢劫。

第三，\CJKunderwave{逻辑推理问题}。全部秘密说误将一个罪的常见情形当作一个罪的必要条件。秘密性只是盗窃罪的常见情形。从逻辑上讲，不能将一个罪的常见情形上升为一个罪的必要条件。虽然百分之八十的盗窃案是秘密进行的，但不能因此将秘密性上升为盗窃罪的必要条件。黑格尔提醒：不能误将熟知当真知。

\subsubsection*{（三）个别公开说的回应}

1.全部秘密说认为，德日刑法的盗窃罪可以公开进行，是因为德日刑法没有抢夺罪；我国刑法有抢夺罪，所以必须要求盗窃罪是秘密的。“个别公开说”的回应是，德日刑法的盗窃罪可以公开进行的观点，与其立法上有无抢夺罪没有关系。\footnote{对此可参考德日权威学者的著作。例如，山口厚、大谷实、雅科布斯、金德霍伊泽尔等教授的教材。}我国台湾地区刑法有抢夺罪，但多数说认为盗窃罪不要求秘密性。\footnote{对此可参考中国台湾权威学者的著作。例如，许泽天：《刑法分则（财产法益篇）》，中国台湾新学林出版股份有限公司2022年版，第28页。张丽卿：《窃盗与抢夺的界限》，载蔡墩铭主编：《刑法争议问题研究》，中国台湾五南图书出版公司1999年版，第504页。曾淑瑜：《刑法分则实例研习一个人法益之保护》，中国台湾三民书局2004年版，第198页。}

2.盗窃罪不要求秘密性，这一点没有违背法律规定。刑法第264条（盗窃罪的条文）中，没有“秘密”二字。

% Source PDF physical page 266; see audit/source_coverage.json.
3.现行司法解释没有规定，盗窃罪必须秘密。\footnote{（1）1998年3月17日最高人民法院《关于审理盗窃案件具体应用法律若干问题的解释》规定：“以非法占有为目的秘密窃取财物，构成盗窃罪”。然而，这句话不等于“盗窃罪是指秘密窃取”。2013年4月4日两高《关于办理盗窃刑事案件适用法律若干问题的解释》规定：废止1998年3月17日的司法解释，并且在新的司法解释中删除“秘密窃取”的表述。（2）最高法院第7批指导性案例“臧进泉案”中，审理案件的基层法官认为盗窃罪要求秘密性。而最高法院对本案提炼的“裁判要点”中没有要求盗窃罪必须秘密。该“裁判要点”才是最高人民法院具有指导性的结论。}

4.有同学问，许多教材和实务判决主张“全部秘密说”，怎么办？如本书前言所述，根据《行政许可法》第54条规定，在法考领域，法律不允许有官方指定教材。对于争议问题，最高等级的选择标准是司法部公布的真题答案，也即2017年第60题及其答案。\footnote{需要说明的是，作者在法考界讲课，从不输出自己的观点，只输出命题官方的观点。网上有人指责作者输出自己的个性观点。这是混淆视听。就考试而言，讲解命题官方观点，竞被指责，让人很无语，也很无奈（满满的求生欲)。}

\subsection*{三、盗窃罪的特殊类型}

1.多次盗窃

根据2013年的相关司法解释，\footnote{2013年4月4日最高人民法院、最高人民检察院《关于办理盗窃刑事案件适用法律若干问题的解释》。}2年内盗窃3次以上的，应当认定为“多次盗窃”。

（1）计算次数的标准。基于同一个盗窃犯意，在同一地点连续针对多个被害人实施盗窃，属于一次盗窃。例如，甲来到公交车上盗窃，连续盗窃了三位乘客的钱包，然后下车，属于一次盗窃。

（2）不要求每次盗窃行为都成立盗窃罪（也即不要求每次盗窃数额较大财物），更不要求都成立盗窃罪既遂。（2026年试题）

（3）已经受过治安处罚的盗窃行为，可以计人“多次盗窃”的次数，但治安拘留期限应折抵刑期。（2026年试题）

2.入户盗窃

（1）“户”是指相对隔离的家庭生活场所。

（2）入户时应具有非法目的，但不限于盗窃目的。如果是合法入户，在户内临时起意盗窃，不属于入户盗窃。例如，为了诈骗主人入户，被主人识破骗局，便盗窃主人财物，也属于入户盗窃。

3.携带凶器盗窃

（1）凶器包括性质上的凶器和用法上的凶器。

（2）携带，不要求显示，更不要求使用，如果对人使用，就属于抢劫了。

（3）携带凶器的状态，要求具有对人使用的意图，譬如防止主人抓捕。例如，甲趁着夜色，拎着镰刀，偷割了邻居的一亩小麦，不属于携带凶器盗窃。

4.扒窃

这是指在公开场合盗窃他人随身财物。

（1）要求在公开场合，也即不特定人出入的场合，包括公共交通工具上。

（2）盗窃他人随身携带的财物。随身，不要求贴身。例如，盗窃他人身上口袋的财

% Source PDF physical page 267; see audit/source_coverage.json.
物，盗窃他人放在座位旁边提包里的财物，盗窃乘客放在行李架上的行李箱中的财物，都属于扒窃。（2011年第16题）

（3）扒窃不要求秘密性，公开扒窃也可成立盗窃罪。

\OriginalPanel{典型真题}

甲的下列哪些行为属于盗窃（不考虑数额）？（2014年·卷二·60题）\footnote{[答案]ABCD。}

A.某大学的学生进食堂吃饭时习惯于用手机、钱包等物占座后，再去购买饭菜。甲将学生乙用于占座的钱包拿走

B.乙进入面馆，将手机放在大厅6号桌的空位上，表示占座，然后到靠近窗户的地方看看有没有更合适的座位。在7号桌吃面的甲将手机拿走

C.乙将手提箱忘在出租车的后备厢。后甲搭乘该出租车时，将自己的手提箱也放进后备厢，并在下车时将乙的手提箱一并拿走

D.乙全家外出打工，委托邻居甲照看房屋。有人来村里购树，甲将乙家山头上的树谎称为自家的树，卖给购树人，得款3万元

\section{抢夺罪}

[问题】狗蛋牵一条藏樊，抢夺小芳的提包。对狗蛋能否适用“携带凶器抢夺，定抢劫罪”的规定？（牵条藏案）

抢夺罪（第267条），是指以非法占有为目的，当场直接夺取他人紧密占有的财物的行为。

\subsection*{一、行为结构}

\InsertDiagram{figures/manual_p267_1.tex}

1.抢夺罪的本质是对物实施暴力。“抢”说明手段要有暴力，不是平和手段。“夺”说明是夺取财物，对财物实施暴力。

首先，不能是对人实施暴力，否则就是抢劫罪。

其次，对物实施暴力要求对人具有人身危险性，否则就是平和手段，构成盗窃罪。手段何时才具有人身危险性？一般而言，当被害人紧密占有财物时，夺取财物的手段就具有人身危险性。（2010年第59题）

例1，甲拎着包逛街，乙从后面夺下包迅速逃离。因为甲对财物紧密占有，乙的行为对甲具有人身危险性，构成抢夺罪。

例2，甲拎着包逛街，不慎摔倒，包摔出5米外。乙路过迅速捡起逃离。因为甲对财物没有紧密占有，乙的行为对甲没有人身危险性，构成盗窃罪。（2016年第18题）

[巩固练习]题1（2018年试题），乙骑摩托车载着甲，行驶到泥泞路面。甲带着非法占有目的，欺骗乙：“我帮你骑过泥泞路段，在前方平坦处等你。”乙答应，甲骑着摩托

% Source PDF physical page 268; see audit/source_coverage.json.
车，驶过泥泞路段，在前方平坦处向乙拜拜！第一，由于被害人乙没有实施处分占有，而是犯罪人甲实施了转移占有，所以甲不构成诈骗罪。第二，乙对财物从紧密占有变成松弛占有。甲在前方平坦处骑走车，不属于对车暴力、夺取车，手段对乙也没有危险，属于平和手段。因此，甲不构成抢夺罪，而构成盗窃罪，属于公开盗窃。再次重申，答案是定盗窃罪，而非抢夺罪。

题2，甲来到摩托车销售店，对店员乙称：“我能否试骑一下？”乙称可以，让甲在店门前空地试骑。乙站在一边观看。甲试骑时，骑到空地另一边，然后直接骑走，向乙拜拜！第一，乙没有处分财物，所以甲不构成诈骗罪。第二，乙对财物属于松弛占有，甲直接骑走的行为，不属于对车暴力、夺取车，手段对乙没有危险，属于平和手段。因此，甲不构成抢夺罪，而构成盗窃罪，属于公开盗窃。

题3，甲拿一把剔骨刀，看到乙背着挎包坐在公园长椅上睡着了，悄悄轻轻割断乙的挎包背带，拿走挎包。甲虽然拿着剔骨刀，但是悄悄轻轻割断背带的行为对乙的人身没有危险，所以不构成抢夺罪，而构成盗窃罪。（2023年试题）

2.盗窃罪、抢夺罪、抢劫罪的区分。

区分标准：手段的暴力程度不同。三罪不是对立排斥关系，而是位阶关系，可以包容评价。\footnote{参见张明楷：《盗窃与抢夺的界限》，载《法学家》2006年第2期。}

\InsertDiagram{figures/manual_p268_1.tex}

[提示1]传统观点认为，盗窃罪与两抢罪的区分在于秘密与公开，两抢罪必须是公开进行的，所谓“明抢暗偷”。这种看法有误。

第一，公开进行只是抢夺罪的常见情形，不是其必要条件。例如，夜晚，甲坐在公园长椅上睡着了，乙悄悄来到甲身后，迅速扯走甲脖子上挂的手机，导致甲的脖子勒出红印。几秒后，甲迷迷糊糊醒来，发现手机没了。乙已经消失在黑暗中。乙的行为具有秘密性，但仍构成抢夺罪。

第二，公开进行只是抢劫罪的常见情形，不是其必要条件。例1，甲欲非法占有乙的财物，暗中向乙的饮料中下迷药，乙昏迷后，甲拿走财物。甲的手段是秘密的，但仍构成抢劫罪，属于昏醉抢劫。例2，深夜，甲欲非法占有乙的财物，悄悄站在乙身后，用铁榔头猛击乙的后脑勺，乙当即昏倒，甲拿走财物。甲的手段是秘密的，但仍构成抢劫罪。

可见，用公开与秘密并不能真正区分盗窃罪与两抢罪，只能区分一些常见情形。然而，考试不喜欢考常见情形，而喜欢考不常见的、非典型的情形。

[提示2]“乘人不备”问题。传统观点认为，乘人不备是抢夺罪的成立条件。这种看法

% Source PDF physical page 269; see audit/source_coverage.json.
有误。乘人不备只是抢夺罪的常见情形，不是抢夺罪的必要条件。例如，在候车室，妇女乙坐在椅子上，甲欲非法占有乙的提包，乙很警惕地抱紧提包，但甲仍走上前一把从乙怀中夺走提包。虽然乙时刻准备着，甲没有乘人不备，但甲仍构成抢夺罪。（2006年第60题）

\subsection*{二、携带凶器抢夺}

第267条第2款携带凶器抢夺的，依照本法第二百六十三条（抢劫罪）的规定定罪处罚。

本款是法律拟制，不是注意规定。本款是把不符合抢劫罪的行为拟制为抢劫罪。如果没有该款规定，携带凶器抢夺，仍应定抢夺罪。例如，刑法没有规定“携带凶器盗窃的，定抢劫罪”，对此只能仍定盗窃罪。

1.凶器

（1）性质上的凶器，主要是法律规定禁止个人携带的违禁品。例如，枪支、爆炸物、管制刀具。

（2）用法上的凶器，是指本来用途不是凶器，但是可以用于杀伤人的物品。例如，菜刀、砖块、钢管、硫酸、藏葵。一件物品能否认定为用法上的凶器，需综合考察：

①该物品的杀伤力大小。例如，菜刀杀伤力大，可以作为用法上的凶器。

②一般人对该物品的畏惧感大小。例如，汽车可以撞死人，但是开车抢夺不属于携带凶器抢夺。因为大街上汽车川流不息，人们对汽车的畏惧感不高（2010年第59题）。又如，在本节开头“牵条藏葵案”中，藏葵有杀伤力，一般人对其有畏惧感，狗蛋属于携带凶器抢夺，定抢劫罪。

2.携带

（1）不要求本人随身携带凶器，可以让第三人携带。例如，行为人可以让手下携带凶器，自己实施抢夺。

（2）不要求显示凶器，可以藏在包里；也不要求向被害人暗示有凶器。如果显示凶器，进行胁迫，直接定抢劫罪。也不要求使用凶器。如果使用凶器，对人暴力，直接定抢劫罪。

（3）要求有随时使用的可能性。例如，甲将一把手枪放在带有密码锁的大行李箱中，提着大行李箱出了酒店，偶然看到一位拎包妇女，便上前抢夺提包。因为甲不可能立刻拿出手枪，所以不构成携带凶器抢夺。

（4）要求有对人使用的意图。例如，甲下班去超市买菜刀回家做饭，买好菜刀放进提包，回家路上偶然看到身旁一位拎包大妈，便突然抢夺。因为甲没有使用菜刀的意图，所以不构成携带凶器抢夺。判断有无对人使用的意图，可以考察该物品平日被携带的可能性大小。例如，甲拎一块砖头走在大街上伺机抢夺。因为一般人不会拎着砖头逛街，因此可以推断甲有对人使用砖头的意图。

\subsection*{三、司法解释的两点规定\protect\footnotemark}\footnotetext{2013年11月18日最高人民法院、最高人民检察院《关于办理抢夺刑事案件适用法律若干问题的解释》。}

1.飞车抢夺问题

司法解释第6条规定：“驾驶机动车、非机动车夺取他人财物，具有下列情形之一的，应当以抢劫罪定罪处罚：（1）夺取他人财物时因被害人不放手而强行夺取的；（2）驾驶

% Source PDF physical page 270; see audit/source_coverage.json.
车辆逼挤、撞击或者强行逼倒他人夺取财物的；（3）明知会致人伤亡仍然强行夺取并放任造成财物持有人轻伤以上后果的。”

上述规定比较繁琐，记忆起来比较困难。其实上述规定都属于注意规定，这些情形本身就符合抢劫罪的构成要件，所以用抢劫罪的构成要件来判断上述情形即可，不需要特别记忆。（2003年第35题，2010年第59题，2025年试题）

2.结果加重犯

上述司法解释将抢夺行为同时过失致人重伤、死亡规定成结果加重犯“抢夺罪致人重伤、死亡”，也即定抢夺罪，加重处罚。例如，甲骑摩托车，趁行人乙不备，一把扯走乙的提包，乙倒地，头部碰到石头，导致死亡。甲构成抢夺罪致人死亡，加重处罚。\footnote{提示：此前的司法解释规定，抢夺致人死亡，同时触犯抢夺罪与过失致人死亡罪，择一重罪论处。注意司法解释的变化。}

\section{抢劫罪}

[问题］乙在公园长椅上睡着，甲偷乙的财物，不远处有个小孩在观看甲的举动。甲担心小孩喊人，来到小孩面前，呵斥道：“如果出声，就割你耳朵。”小孩不敢出声。甲窃得财物。甲构成抢劫罪还是盗窃罪？（威胁割耳朵案）

\subsection*{一、构成要件}

抢劫罪（第263条），是指以非法占有为目的，以暴力、胁迫或者其他方法，强取他人财物的行为。本罪的保护法益有人身法益和财产法益。人身法益是手段型法益，财产法益是目的型法益。本罪的性质由财产法益决定，也即本罪是财产犯罪。\footnote{有司法解释规定：抢劫罪的行为人导致被害人轻伤的，便构成抢劫罪既遂。其理由是，抢劫罪的保护法益有人身法益和财产法益，侵害其中任何一个就构成既遂。然而，这种看法将人身法益和财产法益平等看待。实际上，财产法益才是决定性法益。只有侵害了财产法益，取得财物，实现非法占有目的，抢劫罪才能既遂。抢劫罪既遂要有犯罪所得。例如，甲抢劫乙，将乙打成轻伤，压制了反抗，然后取得乙的财物。如果认为打成轻伤就既遂，那么此时该罪就终局了。一个犯罪只能有一次终局形态，那么此后的取财行为就不能被视为抢劫罪既遂了，该财物就不能被视为犯罪既遂所得了，那算什么？所以，该司法解释规定只能被视为特殊规定，不能推广适用，在承继共犯的场合不需要考虑该规定。法考考查该规定的概率较小。}

\InsertDiagram{figures/manual_p270_1.tex}

\subsubsection*{（一）行为对象}

抢劫罪的对象包括财物的占有者、占有的辅助者、财物的保护者（具有保护意思和保护行动的人），不包括无关的第三人。

例1，甲、乙晚上来到网吧，由乙偷睡着的上网者的财物，甲望风。网吧老板看到，上前阻止，甲用刀威胁老板，老板不敢上前，乙顺利偷到财物。甲、乙构成抢劫罪。

例2，甲看到商店门口停了一辆电动车，乙站在旁边，以为乙是车主，使用暴力将乙推倒在地，致乙轻伤，骑走电动车。实际上，车主是丙，乙是无关的人，临时站在那里。

% Source PDF physical page 271; see audit/source_coverage.json.
虽然甲主观上有抢劫的故意，但应先看客观要件，乙不是合格的抢劫罪的对象。甲对丙构成盗窃罪，对乙构成故意伤害罪，并罚。

[练习]本节开头“威胁割耳朵案”中，小孩没有表现出保护意思和保护行动，不是保护者。甲不构成抢劫罪，而构成盗窃罪。

\subsubsection*{（二）成立条件（第一步，行为方式是强制手段）}

1.暴力。暴力的程度要求足以压制对方反抗。

2.胁迫。这是指，以恶害相通告，使对方产生恐惧心理。恐惧心理的程度要求是足以使对方不敢反抗。

（1）由于要求达到这种程度，所以要求以暴力相胁迫，不包括以非暴力的恶害相通告。例如，以揭发隐私胁迫交付财物，不构成抢劫罪，而构成敲诈勒索罪，因为这种胁迫难以达到不敢反抗的程度。以暴力相胁迫的模式应是：不给钱，就\CJKunderwave{当场（当即）}实现暴力恶害。如此，才能压制对方反抗。

（2）以暴力相胁迫，并不要求行为人真正具有加害能力和加害意思。只要行为人的威胁内容使被害人以为行为人会实现威胁内容即可。例如，甲将一团卫生纸装入背包，进入银行，威胁工作人员：“我包里装着炸弹，给我10万元现金，否则我就引爆！”尽管甲客观上不可能实现恶害内容，依然属于以暴力相胁迫，构成抢劫罪。

3.其他方法。这是指其他足以压制人反抗的方法。

（1）昏醉型抢劫（使对方不知反抗）。以非法占有为目的，用药物、酒精麻醉被害人，然后窃取财物，构成抢劫罪，而非盗窃罪。\CJKunderwave{注意}：被害人自己陷入昏醉状态，行为人单纯利用被害人这种状态取走财物，仅成立盗窃罪。（2005年第60题）

(2）拘禁型抢劫（使对方不能反抗）。例如，甲入室抢劫，将主人乙骗到阳台上然后反锁阳台小隔门，从屋里拿走财物。因为甲造成乙不能反抗的状态，构成抢劫罪。（2005年第60题）

4.“\CJKunderwave{行为与目的同时存在}”原则

成立抢劫罪，必须具备第一步。第一步是抢劫罪的实行行为（强制手段）。如何判断一个暴力行为、胁迫行为、其他方法是抢劫罪的实行行为？根据行为与故意或目的同时存在原则，抢劫罪的强制手段是指带着劫财的目的而实施的强制手段。

例1，甲灌晕了丙，然后离去。乙看到昏醉的丙，拿走其财物。乙不构成抢劫罪，因为乙没有实施抢劫罪的强制手段，乙只是单纯利用甲制造的丙无法反抗的状态。乙构成盗窃。

例2，甲将丙牢牢绑在树上，准备掏丙的钱。乙赶到，将甲赶跑，自己在丙身上掏走钱（掏钱动作对丙没有人身伤害）。乙不构成抢劫，因为乙没有实施抢劫罪的强制手段，乙只是单纯利用甲制造的丙无法反抗的状态。乙构成盗窃，属于盗窃（公开型）。\footnote{乙为何不构成甲的抢劫罪的承继共犯？因为成立共犯，要求带着意思联络，一起制造违法事实。乙将甲赶跑，没有与甲一起取财的意思。乙想独占。}

[巩固练习］题1，甲出于伤害目的将乙打晕在地，然后看到乙的钱财，拿走。甲不构成抢劫罪，因为甲没有实施抢劫罪的强制手段。虽然甲实施了暴力行为，但没有带着取财目的实施暴力行为，不属于抢劫罪的强制手段。甲只是单纯利用自己的故意伤害罪所制

% Source PDF physical page 272; see audit/source_coverage.json.
造的乙无法反抗的状态，拿走乙的钱财，构成盗窃。\footnote{有个司法解释规定：先强奸或伤害，被害人未失去知觉，利用被害人不能反抗、不敢反抗的处境，临时起意使用暴力或胁迫劫取财物，构成强奸罪（或故意伤害罪）和抢劫罪，数罪并罚；被害人失去知觉或者没有发觉，临时起意拿走财物，构成强奸罪（或故意伤害罪）和盗窃罪，数罪并罚。注意，上述行为人要构成抢劫，必须在强奸或故意伤害后实施了独立的抢劫罪的强制手段。如果只是单纯利用之前的强奸或故意伤害造成的被害人无法反抗的状态，拿走财物，只构成盗窃。在这里，被害人有无失去知觉不是关键因素。例如，甲为强奸妇女乙，将乙打倒在地，乙重伤，无法反抗，但仍清醒。甲看到乙掉落在身旁的钱包，捡起离去。甲的拿财行为构成盗窃（公开型）。}（2007年第7题，2021年试题）

题2，甲带着教训的目的，将乙绑在树上，然后发现乙身上的钱包，此时产生非法占有目的，伸手拿钱包。由于乙被绑着，无法反抗，钱包被甲拿去。新观点认为（多数说），甲拿钱包构成抢劫罪。理由是，甲实施的拘禁行为是继续犯，在拘禁期间甲产生非法占有目的，此时拘禁行为仍在持续，该持续的拘禁行为便属于抢劫罪的实行行为。该案与上述

题1的区别在于，题1中，甲的伤害行为已经结束，没有持续。结论：利用自己已经结束的伤害行为，拿钱，不构成抢劫罪；利用自已持续的拘禁行为，拿钱，构成抢劫罪。

\subsubsection*{（三）既遂条件（第三步与第四步具有因果性）}

成立抢劫罪既遂，要求第三步与第四步具有因果关系。行为人虽然取得财物，但是被害人不是因为无法反抗而放弃财物，行为人不构成抢劫罪既遂。

例如，甲暴力抢劫乙的财物，乙挣脱后逃跑，甲追赶，乙逃跑时钱包不慎掉下，乙不知情并跑远。甲赶到捡走钱包，迅速离去。甲拿走财物不是因为乙无法反抗而到手的，而是因为乙遗落了财物。由于缺乏第三步与第四步的因果联系，甲的暴力行为构成抢劫罪未遂，甲的取财行为构成侵占罪既遂。两个行为构成两罪，应并罚。

[巩固练习]题1（2013年第56题），甲欲采取昏醉抢劫的方式非法占有乙的财物，在乙的办公室给乙的水杯投了迷药，然后离去。1小时后，甲估计乙已经喝下昏迷，来到乙办公室，没有看到乙，原来乙喝下后就立即出门了。甲拿走财物。由于乙不在现场，甲拿走财物时，不需要应对乙的反抗，也即甲放迷药的行为对甲的取财行为没有起到应有作用（压制乙的反抗、排除乙的阻力）。因此，甲的下迷药行为构成抢劫罪（未遂），甲的取财行为构成盗窃罪（既遂）。两个行为构成两罪，应并罚。（2024年试题）

题2，甲欲入户抢劫，入户后看到乙在卧室睡觉，将其卧室门反锁，拿走客厅财物。事后乙才苏醒。甲实施了足以压制反抗的行为，成立抢劫罪。但甲拿走财物不是基于乙被剥夺现实自由、无法反抗。甲构成抢劫罪（未遂）和盗窃罪既遂，并罚。如果乙在甲取财途中醒来但无法反抗，甲拿走财物，则构成抢劫罪既遂。\footnote{提示：成立拘禁型抢劫罪，不要求拘禁手段达到非法拘禁罪的程度，也即只要求足以压制对方反抗，不要求实际剥夺人身自由。成立拘禁型抢劫罪既遂，要求拘禁手段实际压制反抗。}

\subsection*{二、事后转化抢劫}

第269条犯\CJKunderwave{盗窃、诈骗、抢夺罪}，为\CJKunderwave{窝藏赃物、抗拒抓捕或者毁灭罪证}而\CJKunderwave{当场使用}\CJKunderwave{暴力}或者以暴力相威胁的，依照本法第二百六十三条的规定定罪处罚。

该条是法律拟制，也即原本应定盗窃罪（诈骗罪、抢夺罪）+故意伤害罪（以暴力致人轻伤为前提），但法条拟制为抢劫罪。其成立条件有三个：第一，三个轻罪；第二，三个目的；第三，实行行为。

% Source PDF physical page 273; see audit/source_coverage.json.
\subsubsection*{（一）三个轻罪：犯盗窃、诈骗、抢夺罪}

1.\CJKunderwave{三罪范围}。传统看法是，三个罪是指三种行为，而非罪名。这种看法不妥当。例如，狗蛋在乱坟岗盗窃小芳的户体，想奸户，被人发现，为抗拒抓捕实施暴力。虽然狗蛋实施了盗窃行为，但不转化为抢劫罪。

正确标准：一个行为能否转化抢劫，需要符合两点：第一，该行为具有财产犯罪属性。这是因为，要转化成抢劫罪，而抢劫罪属于财产犯罪。\CJKunderwave{二者在属性上具有相通性}。

例如，上述狗蛋的盗窃户体行为不具有财产犯罪属性，因此无法转化为属于财产犯罪的抢劫罪。第二，该行为符合盗窃罪、诈骗罪、抢夺罪这三罪的构成要件。

例1，盗伐林木的行为符合盗窃罪的构成要件，可以转化为抢劫罪。

例2，信用卡诈骗罪也符合诈骗罪的构成要件，也可以转化为抢劫罪。

例3，使用假币欺骗他人钱财的行为，也符合诈骗罪的构成要件，也可以转化为抢劫罪。

例4，盗窃户体罪、招摇撞骗罪、抢夺国有档案罪等，由于不具有保护财产法益的性质，不具有财产犯罪的属性，所以不能转化为抢劫罪。

2.\CJKunderwave{三罪的时间阶段}。要求着手实行，但不要求既遂。也即，三罪在预备阶段不转化为抢劫。例如，狗蛋欲入户盗窃小芳家，正在翻院墙时，被保安发现，为抗拒抓捕，将保安打成重伤。狗蛋不转化抢劫，定盗窃罪（犯罪预备）和故意伤害罪，并罚。\footnote{（1）司法解释规定：\CJKunderwave{三罪针对的财物}，一般情况下要求是“数额较大的财物”，但如果使用暴力致人轻微伤或更严重伤害，则不要求是“数额较大的财物”。例如，狗蛋盗窃小芳的自行车（300元），为抗拒抓捕实施暴力，导致小芳轻微伤，狗蛋构成转化型抢劫罪。

（2）司法解释规定：“已满14周岁不满16周岁的人盗窃、诈骗、抢夺他人财物，为窝藏赃物、抗拒抓捕或者毁灭罪证，当场使用暴力，故意伤害致人重伤或者死亡，或者故意杀人的，应当以故意伤害罪或者故意杀人罪论处。”原本，这种情形可以转化抢劫，按照抢劫罪追究刑事责任。司法解释在此作了特殊规定，做题应按司法解释作答。}

\subsubsection*{（二）三个目的}

使用暴力或以暴力相威胁，要求带着三个目的：“为窝藏赃物、抗拒抓捕或者毁灭罪证”。例如（2008年第15题），甲抢夺到妇女乙的项链，以为是不值钱的项链，打了乙两耳光，骂道：“出来混，也不知道戴条好的！”甲不是出于三个目的，不转化抢劫。

[总结]转化型抢劫与正常抢劫的区分

\CJKunderwave{区分标准}：\CJKunderwave{使用暴力的目的不同}。\footnote{根据行为与目的同时存在原则。一个暴力行为究竟是哪种抢劫的实行行为，关键看使用暴力的目的。}正常抢劫，行为人是为了继续取得财物。转化型抢劫，行为人是为了窝藏赃物、抗拒抓捕、毁灭罪证。判断步骤如下：

\InsertDiagram{figures/manual_p273_1.tex}

% Source PDF physical page 274; see audit/source_coverage.json.
例1，甲偷到钱包，主人乙抓捕甲，为了窝藏赃物、抗拒抓捕而对乙使用暴力，定转化型抢劫。

例2，甲正在偷钱包，主人乙护住钱包，甲为了继续取财，对乙使用暴力，定正常抢劫。

例3，甲正在偷钱包，主人乙抓捕甲，甲为了抗拒抓捕，对乙使用暴力，定转化型抢劫。

例4，甲正在偷钱包，主人乙抓捕甲，甲为了抗拒抓捕，也为了继续取财，对乙使用暴力，打倒乙，拿走钱包，定正常抢劫。

\subsubsection*{（三）实行行为：当场使用暴力或以暴力相威胁}

1.\CJKunderwave{当场}。当场是个时间概念。这是要求\CJKunderwave{三个轻罪和使用暴力之间具有时间上的紧密}\CJKunderwave{联系，也即时间上相隔不是很久}。\CJKunderwave{注意}：判断重点不是空间上相隔多远，而是时间上相隔多久。不要误将“当场”理解成“现场”。“现场”是个空间概念。考试经常给个空间距离，这是想让考生据此来判断时间相隔多久。（2003年第35题，2014年第8题，2015年第17题)

\InsertDiagram{figures/manual_table_p274_1.tex}

例如，甲在行驶的大巴车上盗窃了乙的财物。一小时后，大巴车进入高速公路的服务区，乙发现财物被甲盗窃，要抓捕甲，甲为了抗拒抓捕将乙打成重伤。甲的盗窃行为与暴力行为相隔1小时，相隔太久，不具有时间上的当场性，因此不转化抢劫。甲构成盗窃罪和故意伤害罪，并罚。（2022年试题）

2.\CJKunderwave{使用暴力或以暴力相威胁}

（1）“使用暴力”，是指故意使用暴力。

例1，乙盗窃后逃跑，主人甲追赶，乙为了摆脱追赶，跨越栏杆，导致栏杆倒塌，刚好砸倒甲，导致甲轻伤。乙不是故意对甲使用暴力，不构成转化抢劫罪。

例2，乙飞车抢夺后逃跑，主人甲追赶，乙的车轧到小石子，小石子弹起来砸中甲的太阳穴，导致甲受伤。乙不是故意对甲使用暴力，不构成转化抢劫罪。

（2）\CJKunderwave{“以暴力相威胁”}，是指行为人以施加暴力相通告：你（抓捕者）若再抓捕我，就对你（抓捕者）施加暴力，以此使抓捕者产生恐惧心理。

例1，小偷偷到财物后，主人追，小偷用刀刺破自己手臂，并对主人说：“你再追，我就死在你面前。”主人见小偷鲜血直流，一下愣住了。小偷迅速逃离现场。小偷以对自已施加恶害相威胁，不是这里的以暴力相威胁。该盗窃不转化为抢劫。（2008年第62题）

% Source PDF physical page 275; see audit/source_coverage.json.
例2，甲盗窃乙的林木后逃跑，乙追赶，甲威胁乙：“再追，我就告发你的违法事情。”乙放弃追赶。甲不是以暴力相威胁，而是以非暴力手段相威胁，故不转化抢劫。（2020年主观题）

（3）\CJKunderwave{暴力、威胁的对象}。对象不限于被害人，包括其他抓捕人，但不包括动物（狼狗）。（2008年第62题，2017年第60题）

第一，对象错误。例1，甲盗窃后逃跑，误以为路人乙是主人，对乙实施暴力，实际上现场没有抓捕者。多数说认为，甲不转化为抢劫罪，因为“为抗拒抓捕而实施暴力”是个主客观相一致的要素，要求存在现实的抓捕者。甲构成盗窃罪和故意伤害罪，并罚。少数说认为，甲能构成转化抢劫，理由是，“为抗拒抓捕”只是一种主观心理，不要求存在现实的抓捕者。（2021年主观题）

例2，甲趁乙睡觉，偷盗乙的财物，转身离开时看到乙翻身，以为乙要起身抓自已，便将乙打晕。实际上乙只是睡梦中的翻身。多数说认为，甲不构成抢劫，而构成盗窃罪和故意伤害罪，并罚。（2013年主观题）

第二，打击错误。例如（2017年第60题），甲抢夺到财物后逃跑，主人乙追赶，甲为抗拒抓捕，向乙开枪，不慎打中附近行人丙，致丙死亡。首先，甲向乙的开枪行为构成转化抢劫罪。其次，甲过失导致丙死亡，存在打击错误。具体符合说和法定符合说在预定目标上的结论是一致的，区别在于，对实害对象的结论不一致。具体符合说，尊重案件事实，所以，甲构成抢劫罪（过失）致人死亡。法定符合说，特事特办，将过失拟制为故意，所以，甲构成抢劫罪（故意）致人死亡。

（4）\CJKunderwave{暴力、威胁的程度}，要求达到足以压制对方反抗的程度。这是因为，要转化成抢劫罪，而抢劫罪的暴力、威胁要求足以压制对方反抗，故这里的暴力、威胁也要达到这个程度。\CJKunderwave{二者在程度上具有相通性}。例如（真题），小偷偷到财物后逃跑，主人追。小偷向主人泼开水。暴力程度高，转化为抢劫。小偷将报纸揉成一团，砸向主人。暴力程度低，不转化抢劫。\footnote{司法解释规定：以摆脱方式抗拒抓捕，暴力程度较小，未造成轻伤的，不转化抢劫。例如（2003年第35题），小偷取得财物后翻墙逃跑，主人抓住小偷的脚，小偷在墙上顺势端主人一脚。这是摆脱行为，暴力程度不高，不转化为抢劫。如果这一脚导致主人轻伤，则转化抢劫。参见2016年1月6日最高人民法院《关于审理抢劫刑事案件适用法律若干问题的指导意见》。}（2005年第60题，2008年第62题）

3.既遂标准

关于转化型抢劫罪的既遂标准，多数说认为，以前三罪为准，前三罪既遂了，转化型抢劫罪便既遂；前三罪未遂，则转化型抢劫罪未遂。少数说认为，只有使用暴力后最终取得财物，才构成转化型抢劫罪既遂。例如，甲在乙家盗窃到珠宝，出门后便被小区保安发现，保安要抓捕甲，甲为了窝藏赃物、抗拒抓捕，对保安实施暴力，但被保安控制住，并扣押了珠宝。按照多数说，由于盗窃罪既遂，因此甲构成转化型抢劫罪既遂。（2022年试题）

\subsubsection*{（四）转化型抢劫罪的共犯问题}

1.共犯与身份问题

事后转化抢劫=前提事实（实施前罪）+实行行为（事后使用暴力）。“前罪实施者”类似于一种身份，只有具有这种身份才能构成转化抢劫的实行犯。因此，转化抢劫类似于一种真正身份犯。可用该原理来指导转化抢劫的共犯问题。（2017年第87题）

% Source PDF physical page 276; see audit/source_coverage.json.
一人单独盗窃，另一人事后参与。例如，甲单独盗窃并既遂，逃离现场时，乙参与进来，产生不同情形：

情形一：乙教唆甲对抓捕者实施暴力，甲照办。甲构成转化抢劫，乙构成转化抢劫的教唆犯。

情形二：甲教唆乙帮自己断后，乙知道甲实施了盗窃，为了帮助甲，对抓捕者使用暴力。甲不能构成转化抢劫的教唆犯，因为没有实行犯，乙没有身份（不是前罪实施者），不能担任转化抢劫的实行犯角色。为此，只能让甲成为转化抢劫的间接正犯，乙构成转化抢劫的帮助犯。该案类似于总论“间接正犯”一节的“虐囚案”。（2022年试题）

情形三：甲欺骗乙：“这个人要伤害我，你帮我防卫他”，乙不知情，对抓捕者实施暴力。乙不构成转化抢劫，属于假想防卫。甲构成转化抢劫的间接正犯，利用了不知情的乙对抓捕者实施暴力。

情形四：乙为帮助甲，对抓捕者使用暴力，甲对此不知情。由于甲不知情，甲不构成转化抢劫，只定盗窃罪。乙不构成转化抢劫，因为乙没有身份（不是前罪实施者），不能单独构成转化抢劫，要构成转化抢劫的帮助犯，必须以甲构成转化抢劫（正犯）为前提，但甲不知情，不构成转化抢劫。乙也不构成盗窃罪的片面帮助犯或承继的共犯，因为甲的盗窃罪已经结束。乙构成故意伤害罪，同时构成窝藏罪，因为帮助犯罪分子逃匿。两罪想象竞合，择一重罪论处。（2025年试题）

2.共犯与实行过限

二人共同盗窃，一人实行，另一人帮助。一人使用暴力构成转化抢劫，另一人是否也跟着构成转化抢劫，判断标准：知情+参与。例如，甲乙共同盗窃，甲入室，乙在外望风，产生不同情形：

情形一：甲在盗窃时为抗拒抓捕而使用暴力，乙对此不知情。甲构成转化抢劫，乙以盗窃罪论处。如果乙对此知情并继续望风，则构成甲的转化抢劫的共犯。

情形二：甲正在盗窃，乙在望风时，保安抓捕乙，乙为了抗拒抓捕而使用暴力，甲对此不知情。\CJKunderwave{乙虽然是帮助犯，但也可转化为抢劫罪，前罪实施者不限于实行犯，也包括帮}\CJKunderwave{助犯}。乙为了抗拒抓捕自己，使用暴力，构成转化抢劫。甲仍定盗窃罪。

情形三：甲盗窃既遂后逃跑，主人抓捕甲，乙对主人使用暴力，甲对此不知情。乙与甲是盗窃罪的共同犯罪，乙具有身份（前罪实施者），能够转化抢劫。甲仍定盗窃罪。（2020年试题）

\subsubsection*{（五）转化型抢劫罪与“携带凶器抢夺型”抢劫罪的关系}

例如，甲携带凶器抢夺乙的提包，抢夺到手后逃跑，乙追赶，甲为了抗拒抓捕，使用凶器将乙捅成重伤。

第一种处理方案（多数说），首先，甲携带凶器抢夺，构成抢劫罪，属于法律拟制来的抢劫罪。然后，甲又构成故意伤害罪。对抢劫罪和故意伤害罪并罚。\footnote{参见2005年6月8日最高人民法院《关于审理抢劫、抢夺刑事案件适用法律若干问题的意见》第4条规定。}

第二种处理方案（少数说），甲犯抢夺罪，为了抗拒抓捕而当场使用暴力，构成抢劫罪，属于法律拟制来的抢劫罪。该抢劫罪的暴力行为致人重伤，因此构成抢劫罪（故意）致人重伤。（2025年主观题）

% Source PDF physical page 277; see audit/source_coverage.json.
\subsection*{三、法定刑升格条件\protect\footnotemark}\footnotetext{以下内容吸收了司法解释的要点。主要有三个司法解释：（1）2016年1月6日最高人民法院《关于审理抢劫刑事案件适用法律若干问题的指导意见》。（2）2005年6月8日最高人民法院《关于审理抢劫、抢夺刑事案件适用法律若干问题的意见》。（3）2000年11月28日最高人民法院《关于审理抢劫案件具体应用法律若干问题的解释》。}

犯抢劫罪，具备以下法定刑升格条件，加重处罚（十年有期徒刑至死刑）。

\subsubsection*{（一）入户抢劫}

1．“户”的范围。“户”有两大特征：一是场所特征，也即与外界相对隔离；二是功能特征，也即用于家庭生活。这里的家庭不要求是夫妻型家庭，单身居住的出租房，也是单身者的家。（1）除了房屋，封闭的院落，也属于户的范围。（2）户，不要求空间很大。司法解释认为，牧民生活的帐篷属于户。（3）户，不要求固定不动。司法解释认为，渔民生活的渔船属于户。同理，车主长途旅行的房车也属于户。（2024年试题）。（4）刚装修好的无人居住的新房不算“户”。（5）对于\CJKunderwave{部分时间从事经营、部分时间用于生活起居的}\CJKunderwave{场所}，在非营业时间时，属于“户”。（2013年第8题，2025年试题）。（6）没有家庭生活属性的集体宿舍、旅店房间、临时工棚，不属于“户”。\footnote{参见2005年6月8日最高人民法院《关于审理抢劫、抢夺刑事案件适用法律若干问题的意见》。}

2.“入户”的目的。（1）必须带着实施人身或财产犯罪的非法目的入户。（2）合法入户，在户内抢劫，不属于“入户抢劫”。应将“\CJKunderwave{入户抢劫}”与“\CJKunderwave{在户内抢劫}”区别开来。

例1，甲在自己居住的户内抢劫同居室友，不属于入户抢劫。

例2，甲教唆外面的乙抢劫甲的同居室友丙。乙构成入户抢劫。甲不构成入户抢劫。如果丙是寄宿在甲家，则乙不构成入户抢劫，因为乙进入户内只需要经甲的同意，不需要经丙的同意。

［巩固练习]题1，甲在乙家做客，临时起意，将乙迷晕，窃取财物。甲构成昏醉抢劫，但不构成入户抢劫。（2024年试题）

题2，卖淫女甲趁家人不在家，招来客乙来到家里，性交易完成后，乙看到甲挺有钱，便抢劫了甲。乙不构成入户抢劫。

3.抢劫行为。必须在户内实施暴力、胁迫等强制手段。

例1，为了抢劫，将主人骗至户外，然后抢劫，不属于入户抢劫。

例2，在户外将主人打晕在地，然后入户拿走财物，不属于入户抢劫。

［巩固练习］题1，为了抢劫而入户，使用暴力将主人赶出户外，然后拿走户内财物，属于入户抢劫。

题2，甲在户外持刀抢劫乙，乙称自己没带钱，要回家拿。于是甲拿着刀，押着乙，进入乙家，取得财物。甲构成入户抢劫。（2025年试题）

4.主观认识。必须认识到是他人的“户”，在此需要主客观相一致。例如，甲以为抢劫的是仓库，实际上是他人住宅，不构成入户抢劫。

5.转化抢劫。入户盗窃、诈骗、抢夺时，为了窝藏赃物、抗拒抓捕、毁灭罪证，在户内使用暴力或以暴力相威胁，转化为抢劫，同时也认定为入户抢劫。在户外使用暴力，转化为抢劫，不认定为入户抢劫（2025年试题）。同理，入户盗窃时，被主人发现，为了继续取得财物，使用暴力直接升级为抢劫，属于入户抢劫。（2017年第60题）

% Source PDF physical page 278; see audit/source_coverage.json.
6.共犯。共犯人是否构成“入户抢劫”，需要根据参与时间来判断。

例1，甲乙共谋入户抢劫，甲负责进入户内抢劫，乙负责望风。乙构成“入户抢劫”的帮助犯，对乙也适用“入户抢劫”的法定刑。

例2，甲入户后，抢劫过程中，电话叫乙为自己的抢劫望风。乙照办。乙构成甲的抢劫罪的承继的帮助犯。由于乙是在甲“入户”后才帮助的，因此不构成“入户抢劫”的帮助犯。

例3，甲入户后，实施抢劫，将主人丙打晕，想搬走保险柜，太沉，电话叫乙来帮忙，乙进来后一起搬走保险柜。乙构成甲的抢劫罪的承继的正犯。虽然乙进入户内，但“入户抢劫”要求在户内实施暴力行为，乙在户内没有实施暴力行为，因此不构成“入户抢劫”。（2022年试题）

\subsubsection*{（二）在公共交通工具上抢劫}

1．“公共交通工具”的范围。（1）不含小型出租车。（2）接送职工的单位班车、接送师生的校车等大、中型交通工具，视为“公共交通工具”。

2.行为方式。（1）公共交通工具须处在运营状态。（2）既包括身体进入交通工具内抢劫，也包括拦截后，身体不进入交通工具进行抢劫。（3）抢劫对象包括乘客、司机及其他特定人员。

3.转化抢劫。在公共交通工具上盗窃、诈骗、抢夺，为窝藏赃物、抗拒抓捕、毁灭罪证，在公共交通工具上当场使用暴力或以暴力相威胁，转化为抢劫罪，同时认定为在公共交通工具上抢劫。

注意一：转化抢劫要演变为“在公共交通上抢劫”，要求暴力发生在公共交通工具内。普通抢劫要演变为“在公共交通工具上抢劫”，不要求这一点，也即，拦截后，身体不进入交通工具进行抢劫，也属于“在公共交通工具上抢劫”。（2022年试题）

注意二：在户内、公共交通工具上，携带凶器抢夺，虽然能定抢劫罪，但不属于“人户抢劫”“在公共交通工具上抢劫”，因为在户内、公共交通工具上没有对人使用暴力，只是对物暴力。

\subsubsection*{（三）抢劫银行或者其他金融机构}

1.抢劫对象须是经营资金。仅仅抢劫银行的办公用品不属于“抢劫银行”；抢劫银行大厅储户身上的现金，也不属于“抢劫银行”。

2.抢劫正在使用中的银行或者其他金融机构的运钞车的，视为“抢劫银行或者其他金融机构”。

\subsubsection*{（四）多次抢劫}

“多次抢劫”是指抢劫3次或以上，并且每次均构成抢劫罪。基于同一个抢劫犯意，在同一地点连续针对多个被害人实施抢劫，属于一次抢劫。例如，甲来到公交车上抢劫，连续抢劫了三位乘客，然后下车，属于一次抢劫。

\subsubsection*{（五）抢劫致人重伤或死亡（结果加重犯）}

1.主观心理，包括抢劫罪过失致人重伤、死亡和抢劫罪故意致人重伤、死亡。例1，甲抢劫乙时，不慎将乙打死。甲构成抢劫罪过失致人死亡。例2，甲抢劫乙，为了压制乙的反抗，杀死乙，然后取财。甲构成抢劫罪故意致人死亡。这便是常说的谋财害命。

% Source PDF physical page 279; see audit/source_coverage.json.
2.实行行为。抢劫罪致人重伤、死亡，是指抢劫罪的实行行为致人重伤、死亡。认定抢劫罪的实行行为，可根据“行为与故意或目的同时存在原则”判断。抢劫罪的实行行为，是指带着非法占有他人财物的目的，实施的暴力、胁迫等行为。不能评价为抢劫罪的实行行为的其他行为导致重伤或死亡结果，不属于抢劫罪致人重伤或死亡。对其他行为导致重伤、死亡需单独定罪。（2016年第16题）

例1，甲欲抢劫乙，为了制服乙，使用暴力将乙打成重伤或打死，属于抢劫致人重伤或死亡。

例2，甲抢劫既遂后逃跑，主人追赶，为了摆脱，将主人打成重伤，不属于抢劫致人重伤，而构成抢劫罪和故意伤害罪，并罚。若为了灭口，杀死主人，构成抢劫罪和故意杀人罪，并罚。\footnote{司法解释规定：“行为人为了劫取财物而预谋故意杀人，或者在劫取财物过程中，为制服被害人反抗而故意杀人的，以抢劫罪（致人死亡）定罪处罚。行为人实施抢劫后，为灭口而故意杀人的，以抢劫罪和故意杀人罪定罪，实行数罪并罚。”该司法解释符合上述原理，属于注意规定，不是特殊规定，因此不需要死记硬背。}（2012年第59题）

3.加重结果。抢劫罪致人重伤、死亡，是指抢劫罪的实行行为致人重伤、死亡。“致人重伤、死亡”的“人”的范围包括：

（1）最初的被害人（财物的主人）。例如，甲抢劫乙，将乙打死。

（2）前来阻止抢劫的人。例如（2009年第58题），甲抢劫乙，丙前来阻止，甲重伤丙，进而顺利抢劫了乙。甲构成抢劫罪致人重伤。

（3）打击错误中的其他人。这是指，为了抢劫，向财物主人开枪，过失致无关第三人重伤、死亡。例如，甲抢劫乙，向乙开枪，不慎打死行人丙，甲构成抢劫罪致人死亡。具体符合说认为构成抢劫罪（过失）致人死亡，法定符合说认为构成抢劫罪（故意）致人死亡。（2022年试题）

[提示]上述范围适用于事后转化抢劫的抢劫罪致人重伤、死亡。例如，甲盗窃后逃跑，主人乙追赶，甲为抗拒抓捕，向乙开枪，不慎打中行人丙，致丙死亡。甲构成抢劫罪致人死亡。

\InsertDiagram{figures/manual_table_p279_1.tex}

［总结]原则上，人身犯罪中“致人重伤、死亡”的“人”仅限于最初的被害人，财产犯罪中“致人重伤、死亡”的“人”还包括前来阻止的第三人。背后原理是，人身犯罪保护人身法益，人身法益具有人身专属性，具有不可替代性；财产犯罪保护财产法益，财产法益不具有人身专属性，具有可替代性。

% Source PDF physical page 280; see audit/source_coverage.json.
4.因果关系。抢劫行为与加重结果之间应具有刑法上的因果关系。考试常考存在介入因素的案件。对此，可根据“两步走”标准来判断。

例如，甲为抢劫乙，将乙打成重伤昏迷，甲拿走财物离去。乙被冻死。甲构成抢劫罪致人死亡。（2009年第58题）

[巩固练习］题1，甲持刀抢劫乙，乙为逃命，被迫闯红灯，被车撞死。乙的逃跑行为导致的危险是甲引起的，不异常。甲属于抢劫罪（过失）致人死亡。

题2（2019年试题），甲抢劫到乙的包，甲逃离，乙为了追赶而闯红灯，被车撞死。注意本案例与上述案例不同。乙不是逃跑，而是追赶，为了追回自已的包而冒险闯红灯，比较异常。死亡结果不能归属于甲。甲不构成抢劫罪致人死亡。\footnote{实际处罚方案的选择问题。先回顾下想象竞合犯与结果加重犯的转化过程。例如，鲁智深拳打镇关西，一个行为同时触犯故意伤害罪和过失致人死亡罪，想象竞合，择一重罪论处，定故意伤害罪。但是，立法者认为该处罚结论过轻，进行特殊规定，也即将过失致人死亡作为故意伤害罪的加重处罚结果，生成结果加重犯“故意伤害罪致人死亡”，该处罚绝对比单纯的故意伤害罪重。

同理，甲为了劫财，杀死乙，取走财物。甲的一个暴力行为同时触犯抢劫罪和故意杀人罪，想象竞合，择一重罪论处，定故意杀人罪。但是，立法者认为该处罚结论较轻，进行特殊规定，将故意杀人作为抢劫罪的加重处罚结果，生成结果加重犯“抢劫罪（故意）致人死亡”。

问题在于，立法者给该结果加重犯配置的刑罚并不必然比故意杀人罪更重，也即立法者的配套措施没做好。由此导致，在实际处罚方案的选择上需要根据具体案件情节具体判断。如果抢劫罪（故意）致人死亡的处罚比故意杀人罪重，就按照抢劫罪（故意）致人死亡处理。如果抢劫罪（故意）致人死亡的处罚比故意杀人罪轻，就按照故意杀人罪处理，也即抢劫罪和故意杀人罪想象竞合，择一重罪论处，定故意杀人罪，此时不需要转化为结果加重犯。}

\subsubsection*{（六）冒充军警人员抢劫}

1.“冒充军警人员”，要求足以使一般人相信是军警。对此，要考察行为人是否穿着军警制服、携带枪支、出示军警证件等。对于行为人仅穿着类似军警的服装或仅以言语宣称系军警人员但未携带枪支、也未出示军警证件而实施抢劫的，一般而言，不足以使一般人相信是军警。

2.“冒充军警人员抢劫”包括交警冒充防暴警察、军人冒充警察的情形。军警人员利用自身的真实身份实施抢劫的，不认定为“冒充军警人员抢劫”（司法解释）。

\subsubsection*{（七）持枪抢劫}

1.枪须是真枪，不包括仿真枪，但不要求实弹，可以是空枪。

2.枪须向被害人使用或显示。

例如，甲携带枪支抢夺，定抢劫罪，但不属于持枪抢劫。这是因为，携带凶器抢夺，定抢劫罪。在此，枪支已经被评价为凶器，此后就不能再评价枪支了，否则违反了禁止重复评价原则。

[巩固练习］题1，甲犯盗窃罪，为了抗拒抓捕，当场使用枪支，定抢劫罪，属于持枪抢劫。这是因为，盗窃罪要转化为抢劫，只要求当场使用暴力，并不要求当场使用凶器；枪支在此尚未被评价，所以可以评价给抢劫。

题2，甲为了抢劫，使用枪支逼迫乙喝下毒药，乙中毒昏迷，甲拿走财物。甲构成持枪抢劫。

\subsubsection*{（八）抢劫军用物资或者抢险、救灾、救济物资}

[提示]认识错误问题。例如，甲欲抢劫救灾物资，抢到手发现是军用物资。这种对

% Source PDF physical page 281; see audit/source_coverage.json.
象认识错误属于同一构成要件内的错误，不影响该法定刑升格条件的适用。

\subsection*{四、抢劫罪与绑架罪}

1.两罪区分

两罪的客观行为均可以表现为劫持被害人。\CJKunderwave{二者的核心区分在于主观目的不同}：抢劫罪，行为人是带着向A要钱的目的而劫持了A；绑架罪，行为人是带着向B要钱的目的而劫持了A。

例1（2006年第14题），甲使用暴力将乙扣押在某废弃的建筑物内，强行从乙身上搜出现金3000元和1张只有少量金额的信用卡，甲逼迫乙向该信用卡中打入人民币10万元。乙便给其妻子打电话，谎称自己开车撞伤他人，让其立即向自己的信用卡打入10万元救治伤员并赔偿。乙妻信以为真，便向乙的信用卡中打入10万元，被甲取走，甲在得款后将乙释放。甲劫持乙的目的是向乙要钱，因此甲构成抢劫罪。

例2（2010年第16题），甲持刀将乙逼入山中，让乙通知其母送钱赎人。乙担心其母心脏病发作，遂谎称开车撞人，需付5万元治疗费，其母信以为真。甲劫持乙的目的是向乙的母亲要钱，因此甲构成绑架罪。

[注意]不具有通用性的区分标准。

（1）用是否形成现实的三角关系来区分，抢劫罪只是两角关系，绑架罪形成现实三角关系。

有观点认为，要成立绑架罪，必须形成三角关系，也即绑匪向第三人勒索了钱财。实际上，成立绑架罪，并不要求实施行为二（向第三人勒索财物）。行为二只是常见情形，并非必要条件。通过上述例1、例2发现，“第三人打钱”这个环节不重要。例2中，当甲带着向乙的母亲要钱的目的而劫持了乙，甲就构成绑架罪，并且既遂。不能将一个罪的常见情形作为区分标准。若用“第三人打钱”这个环节区分绑架罪和抢劫罪，则例1也构成绑架罪。这显然是错的。命题组就是有意用“第三人打钱”这个环节来干扰考生。

（2）用“空间上的当场性”来区分，抢劫罪具有当场性，绑架罪不具有当场性。

实际上，“空间上不在当场”只是绑架罪的常见情形（先将人质劫持到某个秘密地点，然后通知第三人限期打钱），但并非必要条件。有些绑架罪也具有空间上的当场性。

例如，警方追捕逃犯甲时，甲用枪将一名行人劫持在身前，命令警方放下枪并退下，否则就杀死人质。双方陷入对峙局面。甲当场劫持人质，当场提出不法要求，当场要求警方满足要求，甲构成绑架罪。何况，抢劫罪不要求空间上的当场性。所以，不能将一个罪的常见情形作为区分标准。

2.两罪关系

\CJKunderwave{结论}：抢劫罪与绑架罪不是A与-A的对立排斥关系，而是A与B的中立关系。因此，二者可以想象竞合。这是因为，绑架罪是人身犯罪，抢劫罪是财产犯罪，虽然人身犯罪与财产犯罪有所不同，但不是A与-A的不同（对立排斥关系），而是A与B的不同（中立关系）。人身犯罪和财产犯罪完全可以想象竞合。

例如，乙带着女儿丙（4岁）逛街。甲一把将丙劫持过来，用刀架在丙肩膀，威胁乙“不给1万元，就要孩子的命！”乙被迫给钱。对同一个案件事实可以从不同角度评价。当评价重点放在乙身上，甲对乙构成胁迫型抢劫罪，也即“若不给钱，就对你或你的亲人施加恶害”。甲完全符合抢劫罪的构成要件。当评价重点放在丙身上，甲对丙构成绑架罪，也即甲劫持丙，将丙作为人质，侵害丙的人身安全。甲完全符合绑架罪的构成要件。甲同时构成抢劫罪和绑架罪，想象竞合，择一重罪论处。有些同学，将绑架罪与抢劫罪总是理

% Source PDF physical page 282; see audit/source_coverage.json.
解成对立排斥关系，非此即彼关系，不懂得对同一个案件事实可以从不同角度、不同侧面去提炼。

司法解释并未排斥以上观点。该观点也是法考官方立场。真题证据，例如（2017年第15题），下列说法是否正确：“乙闯入银行营业厅挟持客户王某，以杀害王某相要挟，迫使银行职员交给自已20万元。乙不构成抢劫罪，仅成立绑架罪”。官方答案及解析：该说法错误，绑架罪和抢劫罪可以想象竞合，乙同时触犯绑架罪和抢劫罪，想象竞合，择一重罪论处。

3.罪数问题。（1）先抢劫，后绑架，数罪并罚。例如，甲使用暴力压制乙的反抗，抢劫乙，只抢到100元，不满意，便继续控制乙，给乙的妻子打电话，勒索财物。甲构成抢劫罪和绑架罪，并罚。（2）先绑架，后抢劫，司法解释规定：重罪吸收轻罪。\footnote{2005年6月8日最高人民法院《关于审理抢劫、抢夺刑事案件适用法律若干问题的意见》。}例如，甲绑架了乙，向乙的妻子打电话要钱。甲在等待期间，看到乙身上有钱，抢走该钱。甲构成绑架罪和抢劫罪，重罪吸收轻罪。

\OriginalPanel{典型真题}

关于抢劫罪的认定，下列哪些选项是正确的？（2017年·卷二·60题）\footnote{[答案]ABD。}

A.甲欲进王某家盗窃，正撬门时，路人李某经过。甲误以为李某是王某，会阻止自己盗窃，将李某打昏，再从王某家窃走财物。甲不构成抢劫既遂

B.乙潜入周某家盗窃，正欲离开时，周某回家，进屋将乙堵在卧室内。乙掏出凶器对周某进行恐吓，迫使周某让其携带财物离开。乙构成入户抢劫

C.丙窃取刘某汽车时被发现，驾刘某的汽车逃跑，刘某乘出租车追赶。途遇路人陈某过马路，丙也未减速，将陈某撞成重伤。丙构成抢劫致人重伤

D.丁抢夺张某财物后逃跑，为阻止张某追赶，出于杀害故意向张某开枪射击。子弹未击中张某，但击中路人汪某，致其死亡。丁构成抢劫致人死亡

\section{敲诈勒索罪}

[问题］夜晚，小芳独自坐在校园长椅上欣赏月光，书包在身旁。狗蛋为了非法占有该书包，戴着鬼面具跳出来吓唬小芳，小芳很恐惧，跑开后躲在附近查看究竟。狗蛋拿走小芳的包。狗蛋是否构成敲诈勒索罪？（鬼面具案）

\subsection*{一、构成要件}

\InsertDiagram{figures/manual_p282_1.tex}

% Source PDF physical page 283; see audit/source_coverage.json.
1.成立条件（第一步）

第一步是本罪的实行行为，也即恐吓行为（胁迫、威胁、敲诈行为）。行为模型：以恶害相通告，使对方产生恐惧心理。

（1）“以恶害相通告”。例如，本节开头“鬼面具案”中，狗蛋没有以恶害相通告，不属于恐吓行为，不构成敲诈勒索罪，构成盗窃罪。

[注意】恶害的内容既可以是违法内容，也可以是合法内容。例如，甲得知乙有受贿事实，向乙勒索5万元，否则就向司法机关告发。尽管向司法机关告发是合法的，但甲不能将此作为敲诈钱财的手段。甲构成敲诈勒索罪。

（2）“使对方产生恐惧心理”。例如（2006年第15题），甲捡到乙的手机及身份证，给乙打电话，索要3000元，并称若不付钱就不还手机及身份证。乙迫于无奈付给甲3000元现金，赎回手机及身份证。由于“不给手机和身份证”不是重大恶害，不会让乙产生恐惧心理，因此甲不构成敲诈勒索罪。

[注意］恶害内容要能使对方产生恐惧心理，必须以对方相信该恶害内容为前提。例如，甲对乙说：“你如果不给我1万元，我会让天上的雷来劈你！”乙对此不会相信，只会觉得可笑。甲的行为就不属于恐吓行为。

2.既遂条件（第三步+第四步）

第三步是“对方基于恐惧心理而交付财物”，第四步是“行为人取得财物”。行为人的恐吓行为与取得财物之间要具有因果关系，否则成立敲诈勒索罪未遂。

例1，甲恐吓乙交付1万元，乙实际上是当地真正的黑社会老大，没有产生恐惧心理，反而很欣赏甲的勇气，便给甲1万元。甲成立敲诈勒索罪未遂。

例2（真题），甲敲诈勒索乙的钱财，让乙将钱放到指定地点，否则将乙的丑闻曝光。乙报警后，警察让乙按照甲的指示去放钱，以此诱捕甲。乙配合警方，将钱放到指定地点，甲刚拿到钱后便被埋伏的警察抓捕。甲是敲诈勒索罪的未遂，而非既遂。即使甲没有当场被抓，甲也是敲诈勒索罪的未遂，因为乙不是基于恐惧心理而交付财物。

\subsection*{二、敲诈勒索罪与抢劫罪的区分}

\InsertDiagram{figures/manual_p283_1.tex}

% Source PDF physical page 284; see audit/source_coverage.json.
例如，黑社会分子甲来到乙的饭馆收保护费。

情形1，甲扇乙一耳光，警告：“给钱！”乙害怕而给钱。甲构成敲诈勒索罪，因为甲的暴力（扇耳光）不足以压制乙的反抗。乙是有得选，但最好给。这表明，\CJKunderwave{第一，敲诈勒}\CJKunderwave{索罪的手段可以是暴力。第二，敲诈勒索罪可以当场取财}。因此，是否使用暴力和是否当场取财，不是抢劫罪与敲诈勒索罪的区分标准。

情形2，甲威胁：“后天不给钱，后天就放你的血！”乙害怕，后天给了钱，或者当下就给了钱。甲构成敲诈勒索罪，因为乙属于有得选，最好给。

情形3，甲威胁：“现在不给钱，现在就放你的血！”甲构成抢劫罪，因为乙属于没得选，只好给。（2007年第63题）

情形4，甲威胁：“现在不给钱，后天就放你的血！”乙害怕，给了钱。甲构成敲诈勒索罪，因为甲是威胁日后施加暴力，由此乙有两天的考虑时间和选择自由。乙给钱，属于有得选，最好给。这表明，\CJKunderwave{抢劫罪的胁迫手段，要求：不给钱，就当场施加暴力}。这是抢劫罪胁迫手段的第一个当场性。

情形5，甲威胁：“如果不答应后天给钱，现在就放你的血！”乙答应，后天给了钱。甲构成敲诈勒索罪，因为甲让乙后天给钱，由此乙有两天的考虑时间和选择自由。乙给钱，属于有得选，最好给。这表明，\CJKunderwave{抢劫罪的胁迫手段，要求：当场给钱，当场取财}。这是抢劫罪胁迫手段的第二个当场性。

[总结］抢劫罪有三个当场性。当场，是个时间概念，不是空间概念，不是指同一现场。

（1）胁迫手段中有两个当场性。

第一，威胁当场施加暴力，也即若不给钱，就当即施加暴力。

第二，要求当场给钱，当场取财。这里的当场，是指胁迫行为与取财行为具有时间上的紧密联系。例如，甲胁迫乙：“现在不给钱，就杀了你！”乙身上没钱，打电话给朋友丙，让丙替自已还个民事债务，债主会找丙取钱。丙答应照办。甲派丁冒充债主前往取钱，丁从丙处取到钱，甲结束胁迫。虽然胁迫行为与取财行为不在同一现场，存在空间距离，但是胁迫行为一直在持续，与取财行为在时间上具有紧密关联性，属于当场取财，因此甲构成抢劫罪。这表明，抢劫罪不要求胁迫行为与取财行为发生在同一现场。提示，甲不构成绑架罪，因为甲的目的是向乙要钱，而非向第三人要钱。

（2）暴力手段中有一个当场性，也即当场取财。这里的当场，是指暴力行为与取财行为要具有时间上的紧密联系。

例1，甲带着非法占有目的，使用暴力劫持乙，向乙要钱，乙称身上没钱，家里有钱。甲便押着乙，来到十公里外乙的家中取钱。虽然暴力行为与取财行为不在同一现场，存在空间距离，但是，甲的暴力行为一直在持续，与取财行为在时间上仍具有紧密联系，属于当场取财，因此甲构成抢劫罪。这表明，抢劫罪不要求暴力行为与取财行为发生在同一现场。提示，甲不构成绑架罪，因为甲的目的是向乙要钱，而非向第三人要钱。

例2，甲为了非法占有乙的财物，在上海杀死乙，坐火车到乙的北京家里取财。甲不构成抢劫罪，因为甲的暴力行为和取财行为没有时间上的紧密联系，没有当场性。甲构成故意杀人罪和盗窃罪，并罚。

\subsection*{三、敲诈勒索罪与行使权利的区分}

例如，甲的摩托车被乙偷去，甲恐吓乙：“3天之内若不将摩托车交回，就小心你的狗

% Source PDF physical page 285; see audit/source_coverage.json.
命！”乙被迫将摩托车交还给甲。甲在行使权利，没有非法占有目的，不构成敲诈勒索罪。

判断重点应放在行为是否符合敲诈勒索罪的构成要件。只要符合，就成立敲诈勒索罪。该罪成立条件：带着非法占有目的，实施恐吓行为。

第一步，主观上是否有非法占有目的。如果行为人具有财产请求权，在行使权利，则表明没有非法占有他人财物的目的，不构成敲诈勒索罪。行为人要有请求权，需要有请求权基础（事实依据和法律依据）。

第二步，客观上是否实施了恐吓行为。这是敲诈勒索罪的实行行为，也即以恶害相通告，使对方产生恐惧心理，以此威胁对方给钱。

\InsertDiagram{figures/manual_p285_1.tex}

[巩固练习］题1，某W姓明星的情妇甲向该明星声称：“你若不赔偿我青春损失费3000万元，我就将咱们的那些丑事视频曝光。”首先，青春损失费没有法律依据，表明甲不是在行使权利，甲具有非法占有目的。其次，甲的行为属于恐吓行为。甲构成敲诈勒索罪。

题2，王某得知局长刘某有事实，便要求刘某给自己1万元，否则将此事向公安机关告发。王某根本就没有权利要求1万元，具有非法占有目的，而且行为属于恐吓行为，构成敲诈勒索罪。这表明，\CJKunderwave{公民虽然有检举揭发他人违法犯罪的权利，但是不能将此}\CJKunderwave{作为勒索财物的手段}。

[注意]非法占有目的的认定。请求权有事实依据和法律依据，但请求数额有无法律依据，双方存在争议。多数说认为，这种场合，认定行为人没有非法占有目的。例如（2006年第15题），甲到乙的餐馆吃饭，在食物中发现一只苍蝇，遂以向消费者协会投诉为由进行威胁，索要精神损失费3000元。乙迫于无奈付给甲3000元。甲的请求权有事实依据和法律依据，但是，请求数额（3000元）有无法律依据，法律上是否支持，有争议。在这种场合，甲没有非法占有目的，因此不构成敲诈勒索罪。

\subsection*{四、敲诈勒索罪与诈骗罪的关系}

1.相同点：被害人都是基于有瑕疵的意思而处分财物。

2.区分：意思表示有瑕疵的原因不同。诈骗罪是被害人基于认识错误而处分财物，敲诈勒索罪是被害人基于恐惧心理交付财物。前者是被骗，后者是被迫。

3.认识错误与恐惧心理的竞合问题。

诈骗罪中行为人虚构的事实可以有暴力内容，敲诈勒索罪中行为人通告的恶害可以是虚假内容。当这些事实既是虚假的，又是让人恐惧的，行为就既符合诈骗罪又符合敲诈勒索罪，属于想象竞合。（2007年第63题）

例1，甲对女明星乙说：“我偷拍了你的裸照，不给钱就网上曝光！”乙以为甲手中的照片是自已的裸照，实际上是甲伪造的。乙便答应照办。乙一方面产生认识错误，另一方面产生恐惧心理，甲既触犯诈骗罪，又触犯敲诈勒索罪，想象竞合。

% Source PDF physical page 286; see audit/source_coverage.json.
处理办法：当诈骗罪与敲诈勒索罪产生竞合后，一般情况下，择一重罪论处。

例2，甲男和乙女合作，甲男提供当地众多官员的电话号码，乙女逐个打电话碰运气，诸如“张处长，我是会所的小王，请你打钱，否则我把我俩的那种视频曝光”。张处长记性不好，“做贼心虚”，受骗，害怕，于是打钱。甲、乙同时触犯诈骗罪和敲诈勒索罪，想象竞合，择一重罪论处。

4.敲诈勒索与编造恐吓信息型诈骗的区分。

\CJKunderwave{敲诈勒索罪的行为模式是}：我恐吓你，你若不给钱，我就对你施加恶害，由此你对我产生恐惧心理。如果行为人只是转告虚假恐吓信息，不构成敲诈勒索罪，而构成诈骗罪。判断的焦点是：被害人对谁产生恐惧心理。如果对行为人产生恐惧心理，则行为人构成敲诈勒索罪。如果对“第三人”产生恐惧心理，则行为人不构成敲诈勒索罪。

例如（2005年第19题），甲对乙谎称：“我打听到，丙这个黑社会老大要教训你，你给我钱，我帮你摆平。”乙信以为真，因害怕丙的伤害而给了甲一笔钱。甲属于转告虚假恐吓信息，乙不是对甲产生恐惧心理，而是对“黑社会老大”产生恐惧心理，所以，甲不构成敲诈勒索罪，只构成诈骗罪。

[巩固练习]题1，甲对岳父谎称：“我的儿子是你的宝贝外孙，被黑社会老大王某绑架了，人家说不给10万元，就撕票！我没钱，你赶紧弄钱给我，我好赎人。”岳父信以为真，担心外孙安危，便给了甲10万元。甲属于转告虚假恐吓信息，岳父不是对甲产生恐惧心理，而是对“黑社会老大王某”产生恐惧心理，因此，甲不构成敲诈勒索罪，只构成诈骗罪。

题2，莱温斯基欺骗克林顿：“我闺蜜甲知道了我们的事，让我问你要封口费，否则曝光。”克林顿很恐惧，便给了钱。莱温斯基构成诈骗罪，不构成敲诈勒索罪。

[提示】“碰瓷”问题。道路上“碰瓷”，涉嫌敲诈勒索罪、诈骗罪、抢劫罪、交通肇事罪等。是否构成某罪，以是否符合该罪的构成要件为准。\footnote{2020年9月22日最高人民法院、最高人民检察院、公安部《关于依法办理“碰瓷”违法犯罪案件的指导意见》。}

\OriginalPanel{典型真题}

吴某在路上将骑摩托车的王某拦下，威胁王某，称次日向自己转账1万元，否则见一次打一次。王某吓得落荒而逃。吴某追赶王某，想继续威胁，在追出200米后，追丢了王某，于是返回现场将王某的摩托车骑走。对此，下列说法正确的是？（2025年回忆版金题）\footnote{[答案］D。（1）吴某要求王某次日给钱，不是要求当场给钱，给王某保留了一定的选择自由，所以吴某不构成抢劫罪，而构成敲诈勒索罪。（2）吴某骑走摩托车，构成盗窃罪，因为王某虽然离开现场，但是王某不是遗忘了车，王某对车仍有占有意思，王某对车由紧密占有变成松弛占有，王某的短暂离开不影响王某对车的占有。吴某通过平和手段将王某占有的车转移为自己占有，构成盗窃罪。吴某的手段不是对物暴力，所以不构成抢夺罪。（3）吴某不构成敲诈勒索罪既遂，而是未遂，因为王某并没有基于恐惧心理而交付财物，王某离开现场时对车没有交付给吴某的意思。}

A.吴某以暴力相威胁，压制王某的反抗，构成抢劫罪

B.吴某以敲诈的方式取得王某的摩托车，构成敲诈勒索罪既遂

C.吴某利用王某离开现场的机会，抢夺王某的摩托车，构成抢夺罪

D.吴某骑走王某摩托车的行为，属于秘密窃取他人财物，构成盗窃罪

% Source PDF physical page 287; see audit/source_coverage.json.
\section{诈骗罪}

[问题］婚介所老板狗剩欺骗顾客狗蛋：“交1万元会费，给你介绍貌美如花的女子相亲。”狗蛋交了钱，认为相亲对象颜值很低，控告狗剩欺骗。狗剩是否构成诈骗罪？（貌美如花案)

\InsertDiagram{figures/manual_p287_1.tex}

\subsection*{一、构成要件}

\subsubsection*{（一）欺骗行为}

1.欺骗内容。对事物的评价分为事实判断和价值评价。事实判断，能做真假判断。价值评价，不能做真假判断。例如，“真善美”中，“真假”判断是事实判断。“善恶”“美丑”是价值评价。诈骗罪的欺骗是指虚构事实或隐瞒真相。这表明，只能就事实判断进行欺骗。就价值评价进行欺骗，原则上不构成诈骗罪。

例如，本节开头“貌美如花案”中，狗剩不构成诈骗罪中的欺骗，因为美丑是价值评价。假如，狗剩欺骗狗蛋：“交1万元会费，给你介绍单身女子相亲。”狗蛋发现来相亲的小芳已婚，或者发现来相亲的小红是个男的，狗剩构成诈骗，因为这些事项属于事实判断。

2.欺骗方式。虚构事实是作为的欺骗。隐瞒真相是不作为的欺骗。在此关键看行为人\CJKunderwave{有无告知真相的义务}。告知义务的来源与作为义务相同，包括法律、职务、契约、先前行为等产生的告知义务。（2007年第62题）

例1，行为人隐瞒既往严重病史，与保险公司签订生命保险合同，属于欺骗行为。

例2，在金钱借贷关系中，担保物有瑕疵，如果出借人知道就不会出借，借用人隐瞒，属于欺骗行为。

例3，甲在古玩市场看到一件宝物，价值连城，而摊主不识货，廉价出售，甲隐瞒宝物价值而成交。因为甲没有告知真相的义务，交易中的单纯沉默不构成诈骗，属于正常的淘宝行为。又如（真题），乙请鉴赏家甲鉴定一件藏品，甲明知该藏品价值连城，但故意贬其价值，并要求以低价购买，乙答应低价出卖。甲有告知真相义务，故意隐瞒，并提出购买，构成诈骗。

3.欺骗类型。一是使对方从无到有产生认识错误；二是利用对方已有的认识错误。后者包括两种情形：

（1）有告知真相义务。这种情形下，只要不告知，不作为，单纯维持利用已有认识错误，便构成诈骗。当然，如果通过积极作为方式强化利用，则更构成诈骗。

例1，顾客以为摆在柜台上的一部无法使用的笔记本电脑是合格产品，要求购买。销售员不告知并出售，构成诈骗。

例2（2011年第86、87题），甲将一把壶的壶底落款“民國叁年”磨去，放在自己的古玩店里出卖。某日，钱某看到这只壶，误以为是明代文物。甲见钱某询问，谎称此壶确

% Source PDF physical page 288; see audit/source_coverage.json.
为明代古董，钱某信以为真，按明代文物交款买走。又一日，顾客李某看上一幅标价很高的品，以为名家亲笔，但又心存怀疑。甲遂拿出虚假证据，证明该画为名家亲笔。李某以高价买走品。甲作为店家有告知真相义务。甲对钱某不告知，不作为，维持利用钱某已有的认识错误，构成诈骗。甲对李某不告知，并且用积极作为方式强化李某已有的认识错误，更构成诈骗。

（2）没有告知真相义务。这种情形下，单纯维持利用，不构成诈骗；但是，通过积极作为方式强化已有认识错误，则构成诈骗。

例1（默然接受多找钱案），银行柜台员乙给顾客甲找钱时，弄错数字，甲看到这一幕，没有指出来，默不作声地接受。弄对账目数字是柜台员的义务。法律上，顾客没有核对义务和告知真相的义务。甲只是单纯维持利用，没有通过积极作为强化，不构成诈骗。当然，事后乙向甲索要，甲拒不退还，构成侵占罪。

例2（2008年第14题），甲在某银行的存折上有4万元存款。某日，甲将存款全部取出，但由于银行职员乙工作失误，未将存折底卡销毁。半年后，甲又去该银行办理存储业务，乙对甲说：“你的4万元存款已到期。”甲听后，灵机一动，对乙谎称存折丢失，请办理挂失及取款。乙为甲办理了挂失手续，甲取走4万元。甲谎称存折丢失，让办理挂失，并取走钱，属于通过积极作为方式强化乙的认识错误，构成诈骗。

4.欺骗程度：足以使一般人相信。如此才有侵犯他人财产的可能性、危险性。

例如，甲兜售假冒的劳力士手表，对顾客谎称：“这是劳力士，200元卖给你。”一般人不会相信，不构成诈骗。

[总结】罪名中，涉及到虚假的要素，都要求足以使一般人相信。

例1，诽谤罪要求捏造虚假事实。甲诽谤乙“脑子进水”，不构成诽谤罪。

例2，甲伪造面额为70元的假币，不构成伪造货币罪。

例3，甲没有穿制服、没有任何证件，仅凭一张嘴，对乙谎称：“我是警察，不给钱，就砍人了！”甲构成抢劫罪，但不构成“冒充军警人员抢劫”。

\subsubsection*{（二）认识错误}

1.认识错误必须是具有意思自治能力的人产生的，也即受骗者只能是有意思自治能力的人。欺骗幼儿、精神病患者的财物，成立盗窃罪，因为这些人没有意思自治能力，不能为自己的行为负责，其所产生的认识错误不是诈骗罪认可的认识错误。

例如（2006年第17题）：“甲潜入乙家，搬走乙家1台价值2000元的彩电，走到门口，被乙5岁的女儿丙看到。丙问甲为什么搬我家的彩电，甲谎称是其父亲让他来搬的。丙信以为真，让甲将彩电搬走。甲的行为属于诈骗。”官方答案：这种说法错误，甲构成盗窃罪，属于公开盗窃。如果要求盗窃罪必须是秘密的，由于甲与丙公开对话，完全公开，不构成盗窃罪。甲没有对物实施暴力，也不构成抢夺罪。如此对甲只能作无罪处理。

2.机器“被骗”问题。机器不能成为诈骗罪的对象。理由：机器（如ATM机、自动售货机）没有意思自治能力和自我意识，机器的运行不属于刑法上的“行为”。\footnote{少数说：机器能够被骗。}

例1，甲将小铁片塞进自动售货机，机器当作硬币，吐出商品。甲不构成诈骗，而是盗窃。

% Source PDF physical page 289; see audit/source_coverage.json.
例2（2003年第47题），甲利用ATM的出错，不断按键，机器不断吐出现金，甲拿走。甲构成盗窃罪。

[巩固练习］题1（2006年第60题），“甲将一张作废的IC卡插入银行的自动取款机试探，碰巧自动取款机显示能够取出现金，于是甲取出5000元。甲将IC卡冒充借记卡的欺骗行为在本案中起到了主要作用，因而构成诈骗罪”。官方答案是，该说法错误。

题2（2022年试题），甲在无人售货的超市将高价价签换成低价价签，然后在自助付款机上扫描低价价签，后拿走该商品。由于甲在付款过程中，不存在“欺骗受害人，受害人基于认识错误而处分财物”，因此，甲不构成诈骗罪。甲的手段是平和手段，没有对物暴力，因此构成盗窃罪。实际上，在这种无人售货超市的自助付款区上方都有摄像头监控（监控室的人在监控），甲的行为带有“堂而皇之”的特征，但仍构成盗窃罪。

\subsubsection*{（三）处分财物}

成立诈骗罪，要求被害人基于认识错误而处分财物。处分财物是指被害人将自己占有的财物处分给对方占有。

1.处分行为的表现方式

第一，积极举动，也被称为作为方式。例如，甲诈骗乙的钱，乙受骗，将钱交付给甲。这种交付行为就是一种积极举动。

第二，消极静止，也被称为不作为方式。这是指被害人不行使权利。\footnote{这里的“不作为”与不作为犯中的“不作为”含义不同。后者是指负有作为义务而不履行作为义务。前者的“不作为”是指不行使权利。}例如，甲将车伪装成军车，企图骗免过路费。收费员以为该车是军车，便免收过路费。收费员的处分行为的表现方式是不作为，亦即不行使权利。甲构成诈骗罪。

第三，默许，这是指被害人默许、容忍行为人转移财物。例如，乙将院子里的石头卖给丙，丙称次日让工人来运走石头。此事被甲得知，次日甲冒充工人搬运石头。乙误以为甲是丙雇请的工人，便默许甲运走石头。乙的处分行为的表现方式是容忍、默许。甲构成诈骗罪。\footnote{参见张明楷：《诈骗犯罪论》，法律出版社2021年版，第216页。}（2024年主观题）

2.处分行为的效果（直接性要件）

处分行为要达到一种效果，也即直接处分了自己对财物的占有，使自己对财物失去占有，行为人不需要再实施转移占有的行为。例如，甲欺骗乙称：“我会魔法，会将10元钱变成100元。”乙信以为真，给甲10张10元让他变。甲将钱放进碗里，在变来变去时，趁乙不注意，将钱拿走，换成白纸，并嘱托乙10分钟后打开，遂离去。虽然乙基于认识错误将钱“交给”甲，但是，乙的“交给”行为并没有处分占有，乙对钱只是由紧密占有变成松弛占有。真正实施转移占有是由甲实施的。因此甲不构成诈骗罪，而构成盗窃罪。

\subsubsection*{（四）遭受损失}

成立诈骗罪，要求被害人有遭受财产损失的危险性、可能性。成立诈骗罪既遂，要求行为人取得财物，被害人实际遭受损失。被害人有无财产损失的判断标准。

% Source PDF physical page 290; see audit/source_coverage.json.
个别财产说（多数说）：应算个别的账，不应算总账。也即，应考虑被害人处分财物的目的是否实现（目的落空理论），“这笔钱是否花冤枉了”。

整体财产说（少数说）：应算总账，不应算个别的账。如果算总账，不吃亏，就没有财产损失。

例如，甲盗窃了一辆汽车，谎称是自己的车，以二手车的价格卖给乙，乙不知情而购买。

个别财产说（多数说）：乙花钱是想买合法的车，却买到赃车，花钱的目的没有实现，花冤枉钱了，有财产损失。即使车价很便宜，也存在财产损失。甲构成盗窃罪和诈骗罪，并罚。

整体财产说（少数说）：看乙花的钱与车的实际价值是否相当。如果相当，例如车价比较便宜，乙就没有财产损失。甲对乙不构成诈骗罪。

[巩固练习]题1，甲来书店找法考辅导用书，书店老板将公务员考试用书包装好，冒充法考辅导用书卖给甲。甲花300元购买。该公务员考试用书价值300元。个别财产说（多数说）：甲有财产损失。整体财产说：甲没有财产损失。

题2，募捐问题。在此不存在相互给付，而是单方给付。对此，不存在个别财产说与整体财产说之争。根据目的落空理论判断即可。

（1）甲谎称为四川地震灾区募捐，实际上将募捐来的款项据为己有，构成诈骗罪。

（2）甲组织为四川地震灾区捐款。乙是四川人，为了给家乡救灾，捐了1万元。实际上甲将获得的捐款捐给了贵州某养老院。乙的捐款目的失败了，存在财产损失。甲构成诈骗罪。

（3）甲组织为四川地震灾区捐款。乙是四川人，准备捐1000元。甲欺骗乙：“你们单位同事都捐1万元，你好意思捐1000元？”乙觉得有点不好意思，便捐了1万元。甲将捐款交给四川灾区。甲只是让乙的爱心膨胀了一下，乙的捐款目的没有失败，没有财产损失。甲不构成诈骗罪。（2015年第63题）

\subsection*{二、诈骗罪与盗窃罪的区分I（两角关系的情形)}

\CJKunderwave{关键区分}：被害人是否基于认识错误处分财物。盗窃罪缺少诈骗罪的第三步“基于认识错误而处分财物”。这一步是个主客观相一致的概念，包括两个要件：一是客观上有处分行为，二是主观上有处分意识。二者缺一不可。

做题步骤：先审查客观处分行为，后审查主观处分意识。

\InsertDiagram{figures/manual_p290_1.tex}

\subsubsection*{（一）客观处分行为}

\CJKunderwave{处分行为，是指处分占有，也即被害人将自己占有的财物处分给对方占有}。诈骗罪与盗窃罪都是转移占有的财产犯罪。二者区别在于：谁实施了转移占有。诈骗罪，由被害人实施转移占有，也即将自己的财物转移给犯罪人占有。盗窃罪，由犯罪人实施转移占有，

% Source PDF physical page 291; see audit/source_coverage.json.
也即将被害人的财物转移为自己占有。

例如（调虎离山案），甲在商场假装买衣服，欺骗售货小姐乙：“你帮我拿个凳子过来。”乙去拿凳子，甲趁机拿走衣服。被害人乙没有实施处分占有，犯罪人甲实施了转移占有。甲不构成诈骗，而构成盗窃。

[巩固练习]题1，甲与乙一起乘火车旅行。火车在某车站仅停2分钟，但甲欺骗乙说“本站停车12分钟”，乙信以为真，下车购物。乙刚下车，火车便发车了。甲立即将乙的财物转移至另一车厢，然后在下一站下车后携物潜逃。被害人乙没有实施处分占有，犯罪人甲实施了转移占有。甲构成盗窃。（2008年四川卷第15题）

题2，甲来到乙家，对乙谎称：“你家小孩在楼下被人打，还不赶紧去看。”乙赶忙下楼。甲趁机进屋拿走财物。被害人乙没有实施处分占有，犯罪人甲实施了转移占有。甲构成盗窃。

\subsubsection*{（二）主观处分意识}

这个知识点比较疑难，又比较重要，希望勇敢的你耐心学习，多看两遍。

对方客观上有处分行为，但主观上没有处分意识，行为人不构成诈骗罪，而构成盗窃罪。所谓处分意识，是指被害人意识到，自己将自己占有的财物转移给犯罪人占有，也即处分占有的意识。被害人要有处分意识，要求事先意识到自己在占有财物，要求意识到财物的\CJKunderwave{现实存在}。判断的关键点：\CJKunderwave{被害人是否意识到财物的现实存在}。

范例1（调包案），甲在超市，偷偷将一箱饼干中的饼干取出，将一部照相机放进去，重新包装好，拿到收银台。收银员误以为里面是饼干，收了一箱饼干的钱，将箱子交给甲。第一步，客观上，收银员将箱子交给甲，处分了照相机，将照相机转移给甲占有。也即收银员对照相机有处分行为。第二步，主观上，收银员没有意识到照相机的现实存在，因此没有处分照相机的意识。也即收银员对照相机没有处分意识。因此，甲不构成诈骗罪，而构成盗窃罪。\CJKunderwave{总结}：收银员对财物的种类有认识错误，误以为里面是饼干，实际上是相机，表明没有意识到财物（相机）的现实存在。

范例2（调换价签案），甲在超市，偷偷将一部便宜相机（标价2000元）和一部贵重相机（标价2万元）的价签对调，贵重相机上贴着便宜相机的价签，拿着贵重相机到收银台。收银员看到相机本身（没有外包装盒、包装袋），看到机身上的价签，误以为该相机价值2000元，收了甲2000元，将相机交给甲。第一步，收银员对贵重相机有处分行为。第二步，收银员看到相机本身，将该相机交给甲，表明收银员意识到该相机的现实存在，意识到，自己将自己占有的相机转移给甲占有，表明收银员对相机有处分意识。因此，甲构成诈骗罪，而非盗窃罪。当收银员意识到财物的现实存在，就“处分意识”这个要件而言就够了，并不要求收银员在此基础上还意识到财物的数量、价格。成立诈骗罪，只要求被害人有处分意识，并不要求在方方面面都有完全正确的处分意识。如果要求在方方面面都有完全正确的处分意识，则表明被害人没有上当受骗。

[总结]“基于认识错误而处分财物”包括两个心理活动。一是“处分财物”的意识，这便要求意识到财物的现实存在。二是“基于认识错误”，这里的认识错误是指“处分意识”之外的认识错误，例如对财物的数量、价格的认识错误。这些认识错误是引发处分财物（处分行为、处分意识）的动因。

% Source PDF physical page 292; see audit/source_coverage.json.
\CJKunderwave{总结做题方法}：

1.种类错误，定盗窃罪。

被害人误以为是A种类财物，实际是B种类财物，表明没有意识到B种类财物的现实存在，表明没有处分B种类财物的意识，因此对B种类财物，犯罪人构成盗窃罪。

例如（邮票案），甲发现乙的书《错误论的新视角》中夹有一张清代邮票，让乙赠送这本书给自己。乙没有意识到书中夹有邮票，将书赠送给甲。第一步，客观上，乙对邮票有处分行为。第二步，主观上，乙误以为只是赠送一本书（A种类财物），实际上还有邮票（B种类财物），没有意识到邮票的现实存在，因此没有处分邮票的意识。所以，甲不构成诈骗罪，而构成盗窃罪。又如，上述例1“调包案”，性质也相同，定盗窃罪。

2.数量错误，分两种情况：

（1）被害人对财物是整体处分（论斤卖），被害人有数量错误，结论：犯罪人构成诈骗罪。

被害人整体处分A种类财物，被害人认识到A种类财物的现实存在，便有了处分A种类财物的意识。因此，犯罪人构成诈骗罪。被害人上当受骗体现在，对A种类财物的数量（重量、总价）有认识错误。也即，被害人是基于这些认识错误而处分了财物的占有。

例1（米店短斤少两案），甲来到米店买米，让老板称一袋米，甲事先暗中在老板的电子秤上作了手脚，这袋米真实重量是1斤，但电子秤显示是8两。老板收了甲8两的钱，将这袋米交给甲。老板论斤卖米，意识到这袋米（A种类财物）的现实存在，有处分这袋米的意识，因此甲构成诈骗罪。老板上当受骗体现在，对这袋米的重量有认识错误。反过来，如果老板在秤上做手脚，将1斤米显示成1.2斤，收甲1.2斤的钱，则老板也构成诈骗罪。\CJKunderwave{总结}：对同一种类财物，搞短斤少两，构成诈骗罪。

例2（多放一把米案），甲在超市里撕下一个塑料袋，装入散装的大米，交给过秤员乙，乙称了一下，贴上“1斤米”的标签，扎好口子，交给甲，让甲去收银台付款。甲悄悄撕开口子，又抓一把米放进去，扎好口子，交给收银员。收银员以为是1斤米，收取了1斤米的钱。甲偷偷放进一把米，并不是转移占有行为，这个环节不构成盗窃罪。收银员是关键。收银员意识到这袋米的现实存在，有处分这袋米的意识，因此甲构成诈骗罪。收银员上当受骗体现在，对这袋米的重量有认识错误。这也是对同一种类财物搞短斤少两。

例3，在例2中，假设，甲不是多抓一把米，而多塞进一颗枣，埋进大米里，交给收银员，收银员以为只是一袋大米，收了一袋米的钱，将这袋米交给甲。收银员存在种类错误，没有意识到枣的存在，没有处分枣的意识，因此甲对枣构成盗窃罪。

（2）被害人对财物是单个处分（论个卖）。

①被害人意识到所处分的财物的现实存在，对该财物有处分意识。结论：定诈骗罪。被害人上当受骗体现在对该财物的价格有认识错误。例如，上述“调换价签案”就是如此。

②被害人没有意识到所处分的财物的现实存在，对该财物没有处分意识。结论：定盗窃罪。在此，犯罪人加塞了某个财物，导致被害人没有意识到该财物的现实存在。在此，被害人无论是种类错误还是数量错误，都定盗窃罪。

例如（包装盒加塞案），某超市出售剃须刀是论个卖，一个包装盒里只放一个剃须刀。甲将一个剃须刀（价值2000元）的包装盒里的填充材料取出。

情形一（种类错误），悄悄塞进一个手机，将包装盒重新盖好，交给收银员。收银员看到外包装盒上写着“一个剃须刀，2000元”，以为里面只有一个剃须刀，别无他物，收

% Source PDF physical page 293; see audit/source_coverage.json.
了一个剃须刀的钱，将包装盒交给甲。收银员没有意识到手机的存在。甲对手机构成盗窃。或者，甲将剃须刀拿出来，放进一个手机，然后包装好，交给收银员。甲对手机仍是盗窃罪。这种调包行为与加塞行为没有本质区别，都属于种类错误，定盗窃罪。

情形二（数量错误），悄悄塞进一个相同型号的剃须刀（2000元），将包装盒重新盖好，交给收银员。收银员看到外包装盒上写着“一个剃须刀，2000元”，以为里面只有一个剃须刀，别无他物，收了一个剃须刀的2000元，将包装盒交给甲。收银员没有意识到另一个剃须刀的存在。甲对另一个剃须刀构成盗窃。

[总结］无论甲塞进的东西是不同种类的财物（情形一），还是相同种类的财物（情形二），无论收银员是种类错误（情形一），还是数量错误（情形二），收银员都没有意识到加塞的财物的存在，对其都没有处分意识，因此，犯罪人对加塞的财物一律构成盗窃罪。\footnote{提示：情形二“多塞一个剃须刀”与上述例2“多放一把米案”，相同点在于，收银员都有数量错误。区别在于，“多塞一个剃须刀”案中，财物处分规则是单个处分，论个卖。而“多放一把米”案中，财物处分规则是整体处分，论斤卖。日本部分学者没有看出上述两种案件的区别，认为只要是数量错误，一律定诈骗罪。实际上，应看前提条件，处分规则是论斤卖，还是论个卖。}

\InsertDiagram{figures/manual_p293_1.tex}

[应用］处分资金。资金钱款都是以某个单位数量（如1元）来计算，所以，处分资金属于单个处分（论个卖，而非论斤卖）。如果只意识到自已处分了1元钱，实际处分了1001元钱，存在数量错误，没有意识到更多资金的存在，没有处分更多资金的意识，那么犯罪人构成盗窃罪。\CJKunderwave{结论}：处分资金的处分意识，要求意识到具体金额。

例如（2017年第17题），甲冒充银行客服发送短信，称乙手机银行即将失效，需重新验证。乙信以为真，按短信提示输入银行卡号、密码等信息后，又将收到的编号为135423的“验证码”输入手机页面。后乙发现，其实是将135423元汇入了甲账户。乙有处分行为，但没有处分意识，没有意识到处分了这笔资金。甲构成盗窃，而非诈骗。具体而言，甲利用被害人乙的不知情，使乙自已转移占有，构成盗窃罪的间接正犯。

[巩固练习]题1（2016年第17题），臧某给金某发送一个交易金额标注为1元而实际植入了支付30万元的计算机程序的虚假链接，谎称金某点击该1元支付链接后，支付1元，其即可查看到付款成功的记录。金某在诱导下点击了该虚假链接，其银行账户中的30万+1元随即支付到臧某的账户中。金某有处分1元钱的意识，但没有处分30万元的意

% Source PDF physical page 294; see audit/source_coverage.json.
识，臧某构成盗窃罪。\footnote{最高人民法院指导案例27号：臧进泉等盗窃、诈骗案。}如果按照日本学者的看法，只要是数量错误，都定诈骗罪，那么对臧某应定诈骗罪。这种看法显然是错误的。

题2（多加个零案），甲知道自已在饭店消费了3000元，收银员乙在P0S机上故意输入30000元，甲没有注意到多了一个零，支付了30000元。甲虽然有处分3000元的意识，但是没有处分增多的27000元的意识，因此乙构成盗窃罪。（2018年主观题，观点展示）

\subsection*{三、诈骗罪与盗窃罪的区分II（三角关系的情形）}

\InsertDiagram{figures/manual_p294_1.tex}

诈骗罪的结构中，受骗人与处分人必须是同一人，但受骗人与受害人不要求是同一人。

两角诈骗：行为人一受骗人（处分人），同时是受害人。

三角诈骗：行为人—受骗人（处分人）—受害人。也即，受骗人与受害人不是同一人。

三角诈骗与盗窃罪（间接正犯）的区分标准：受骗人有无处分受害人财物的权利地位。若有，则受骗人是合格的处分人，便具备了诈骗罪的第三步，定诈骗罪。

\InsertDiagram{figures/manual_p294_2.tex}

例1（保姆案），甲来到乙家，乙不在家，甲对乙的保姆谎称：“我是干洗店的老板，你家主人让我来取他的衣服。”保姆不知情便将衣服给了甲。甲是行骗人，保姆是受骗人，乙是受害人。基于保姆的职业习惯，保姆具有处分主人衣服的权利地位，是合格的处分人，因此，保姆将衣服交给甲，可以评价为诈骗罪的第三步“处分财物”。所以甲属于三角诈骗，构成诈骗罪。

[提示】法考真题认为，保姆有处分主人旧电脑、旧冰箱的权利地位。但是钟点工没有这种权利地位。

例2（店员案），甲是干洗店老板，欺骗店员：“乙想将衣服拿到咱们店里干洗，你到他家院子，将他的衣服取来。”店员不知情而照办。甲是行骗人，店员是受骗人，乙是受害人。店员不具有处分乙的财物的权利地位，不是合格的处分人，其将衣服拿走给甲，不能评价为诈骗罪的第三步“处分财物”。由于缺少诈骗罪的第三步，所以甲不构成三角诈骗，而是盗窃罪的间接正犯，店员是甲实施盗窃的工具。

[巩固练习］题1，甲发现乙的电脑存放在商场一楼维修间。甲在早晨商场未开门时，对里面的保洁员谎称电脑是自已的，让保洁员递给自己。保洁员不知情而照办。保洁员没有处分的权利地位。甲构成盗窃罪的间接正犯。（2023年试题）

题2，乙进入地铁车厢，看到长椅上有个包，旁边坐着甲，对甲说：“是不是你的包？”甲谎称是的。乙便递给甲。甲拿走包下车。实际上该包是旁边睡着的丙的。甲构成盗窃罪的间接正犯。

% Source PDF physical page 295; see audit/source_coverage.json.
题3，甲伪造欠条，向法院提起虚假诉讼，要求乙偿还借款3万元。法官受骗，判决乙偿还借款。受骗人是法官，法官具有处分乙的财产的权利，具有处分人的资格地位，法官作出判决可以评价为诈骗罪的第三步“处分行为”，因此甲构成三角诈骗。由于该诈骗发生在诉讼领域，又称为诉讼诈骗。甲同时触犯虚假诉讼罪（第307条之一），想象竞合，择一重罪论处，定诈骗罪。

题4（2014年第19题），乙购物后，将购物小票随手扔在超市门口。甲捡到小票，立即拦住乙说：“你怎么把我购买的东西拿走？”乙莫名其妙，甲便向乙出示小票，两人发生争执。适逢交警丙路过，乙请丙判断是非，丙让乙将商品还给甲，有口难辩的乙只好照办。交警不是法官，没有权利处分乙的财物，所以甲不构成三角诈骗，而构成盗窃罪的间接正犯。如果要求盗窃罪具有秘密性，由于甲是公开进行的，不构成盗窃罪。甲也不构成抢夺罪。对甲便无法定罪。这显然是不合理的。命题组通过这道题想揭示“盗窃罪要求秘密性”的缺陷。

\subsection*{四、盗窃罪、诈骗罪、侵占罪的区分}

\subsubsection*{（一）总结打手机问题}

\InsertDiagram{figures/manual_table_p295_1.tex}

\subsubsection*{（二）总结调包问题}

1.盗窃罪与侵占罪

（1）行为人在自己占有前调包，将真货转移为自己占有，构成盗窃罪。

例如，甲在柜台要购买香烟，店员将1条中华香烟递给甲看。甲趁店员招呼其他顾客，将事先准备好的1条假的中华香烟与柜台上的中华香烟对调。结束后，甲对店员表示不买了，然后离开。甲在自已占有前，将店员的财物调包拿走，构成盗窃罪。（2006年第

% Source PDF physical page 296; see audit/source_coverage.json.
17题，2009年第59题）

（2）行为人在自己占有后调包，将真货据为己有，构成侵占罪。

例如，乙将新稻谷运送至甲的米业工厂，欲加工成大米，约定三日后来取货。甲收到新稻谷后，将陈年大米替代新稻谷加工的大米，交付给乙。甲将乙所有的、自己占有的新稻谷，据为已有，构成侵占罪。甲用陈年大米冒充新大米，交付给乙，属于通过欺骗方式免除返还新大米的义务，属于诈骗财产性利益，构成诈骗罪。由于甲的前后两个犯罪给乙只造成一份财产损失，甲也仅获得一份好处，因此对前后两个犯罪不需要并罚。根据重罪吸收轻罪的原理，以诈骗罪论处。（2022年试题）

2.盗窃罪与诈骗罪

（1）顾客调包

例1（调包案），甲在超市，偷偷将一箱饼干中的饼干取出，将一部照相机放进去，重新包装好，拿到收银台。收银员误以为里面是饼干，收了一箱饼干的钱，将箱子交给甲。收银员没有意识到照相机的现实存在，因此没有处分照相机的意识。也即收银员对照相机没有处分意识。因此，甲不构成诈骗罪，而构成盗窃罪。

例2（调换价签案），甲在超市，偷偷将一部便宜相机（标价2000元）和一部贵重相机（标价2万元）的价签对调，贵重相机上贴着便宜相机的价签，拿着贵重相机到收银台。收银员看到相机本身（没有外包装盒、包装袋），看到机身上的价签，误以为该相机价值2000元，收了甲2000元，将相机交给甲。收银员看到相机，对相机有处分意识。因此，甲构成诈骗罪，而非盗窃罪。

例3（调换主板案），甲买到一部手机，在家里拆下新主板，换上废旧主板，然后将手机交给商家，要求退货退款。商家不知情，予以全额退款。甲对自已买到的手机可以任意处置，所以调换主板不构成犯罪。甲用废旧主板冒充新主板，要求退款，构成诈骗罪。（2020年试题）

（2）商家调包

例如，甲在柜台要购买香烟，店员将1条中华香烟递给甲看。甲在看时，接到电话。店员趁甲接电话，将1条假的中华香烟与柜台上的中华香烟对调。甲接完电话，表示愿意购买，不知情而买走假的中华香烟。店员将假货冒充真货卖给甲，构成诈骗罪。

\subsubsection*{（三）总结送快递问题（不考虑数额）}

模型1：快递员将快递放到收件人乙的门口。邻居甲看到后，拿走，构成盗窃罪。

模型2：快递员错将快递放到甲家门口。甲发现后，知道收件人是乙，仍据为已有。快递员将一个快递放到甲家门口，属于放在甲实际控制的领域，甲在占有快递。甲将他人所有、自己占有的财物，变成自己所有，构成侵占罪。（2022年试题）

模型3：快递员敲甲家门，问甲：“你是收件人某某某吗？”甲明知自己不是收件人，隐瞒真相，予以冒领。甲构成诈骗罪。\footnote{本案与前文“默然接受多找钱案”有区别。本案中，快递员在确认收件人身份后，才会交付包裹。这时，甲“谎称是收件人”予以冒领，是一种积极的作为方式，不是单纯的利用维持快递员的认识错误。如果甲不谎称，快递员不会交付包裹。而“默然接受多找钱案”中，柜台员已经做出要给顾客找钱的动作，顾客只是没有指出来，顾客不用谎称，柜台员也会继续实施多找钱的动作。所以，顾客只是单纯利用柜台员的错误行为。}

% Source PDF physical page 297; see audit/source_coverage.json.
模型4：快递员甲在送快递包裹过程中，打开包裹，取走里面的财物。一般的快递包裹包装严密、标示隐私，属于封物。甲没有占有内容物。甲打开包裹，取走内容物，构成盗窃罪。

\section{职务侵占罪}

[问题]村长狗蛋利用职权，贪污村集体的财产。狗蛋构成贪污罪还是职务侵占罪？（村长贪污案）

职务侵占罪（第271条），是指公司、企业或者其他单位的工作人员，利用职务上的便利，将本单位财物非法占为己有，数额较大的行为。本罪实质上是非国家工作人员贪污罪。\footnote{提示：《刑法修正案（十一）》将本罪的法定最高刑从15年有期徒刑修订为无期徒刑。}

\subsection*{一、行为主体}

行为主体是公司、企业或者其他单位的工作人员，也即非国家工作人员。这是与贪污罪中的国家工作人员相对应的概念。

1.国有公司、企业或者其他国有单位中从事公务的人员和国有公司、企业或者其他国有单位委派到非国有公司、企业以及其他单位从事公务的人员，是贪污罪的行为主体。

2.国家机关、国有公司、企业、事业单位中并未从事公务的非国家工作人员，是职务侵占罪的行为主体。

3.在国有资本控股、参股的股份有限公司中从事管理工作的人员，除受国家机关、国有公司、企业、事业单位委派从事公务的以外，不属于国家工作人员，可以成为职务侵占罪的主体。（2008年第63题）

4.村民委员会等村基层组织人员，利用职务便利侵吞集体财产的，以职务侵占罪论处；但是如果在协助人民政府从事行政管理工作时，利用职务上的便利侵占公共财物的，则成立贪污罪。例如，本节开头“村长贪污案”中，村长构成职务侵占罪。（2008年第63题）

\subsection*{二、行为对象}

行为对象是单位的财产，不是国有财产，也不是某个个人的财产。人罪数额是3万元，与贪污罪入罪数额相同。

1.这里的财产包括财产性利益（债权）。例如，公司会计让债务人将欠款打入自己账户，构成职务侵占罪。

2.这里的财产包括利用职务产生的财物。例如，网游公司技术经理利用职权，私下生成游戏装备，卖给他人，贩卖所得归自己所有，构成职务侵占罪。

3.这里的财产包括公司的股份。

例1，A公司享有B公司20\%的股份，A公司委派甲参与B公司的管理。甲利用职务便利，将A公司的股份转移到自已名下，构成职务侵占罪。

% Source PDF physical page 298; see audit/source_coverage.json.
例2，A公司享有B公司20\%的股份，由A公司的经理甲代持，A公司委派甲参与B公司的管理。甲利用职务便利，将代持的股份擅自出售，出售所得归自已所有。甲属于侵占行为，构成职务侵占罪。

[巩固练习]题1，乙享有A公司20\%的股份。A公司的经理甲利用职务便利，在乙不知情的情况下，将乙的股份变更为自已所有。第一种观点认为，甲构成盗窃罪，不构成职务侵占罪，因为甲侵犯的不是单位的财产，而是他人的财产。第二种观点认为，甲构成职务侵占罪。\footnote{这是公安部经侦局的意见，只能认为是特殊规定。}

题2，乙享有A公司20\%的股份，由A公司的经理甲代持。甲利用职务便利，在乙不知情的情况下，将代持的股份擅自出售，出售所得归自己所有。甲侵占乙的财产，构成侵占罪。

\subsection*{三、实行行为}

本罪的实行行为有：侵占，盗窃，诈骗。理由是，刑法第183条第1款规定，保险公司的职员利用职务便利，骗取本公司保险金，定职务侵占罪。\footnote{少数说认为，本罪的实行行为只有侵占，也即，将单位所有、自己基于职务而占有的财物，变成自己所有。少数说是张明楷教授的观点。由于刑法第183条第1款的规定，学界和实务界多数认为，本罪的实行行为不限于侵占。}

侵占罪、盗窃罪、诈骗罪是A，职务侵占罪是A+B。B1是行为主体；B2是行为对象；B3是行为方式，要求利用职务便利。B1、B2、B3，缺少一个，职务侵占罪就退回侵占罪、盗窃罪或诈骗罪。

1.“利用职务便利”中的“职务”

这是指具有一定管理权限的职务。例如，主管、管理、经营等职务。依据这种职权，能够占有或支配单位财物。纯粹的体力性劳务，没有管理权限，只是机械执行的角色，不属于这里的“职务”，其对单位财物只是占有的辅助者，而没有独立占有。对此可联系盗窃罪中关于上下级占有的知识加以判断。

例1，单纯的没有管理权的装卸工、搬运工、分抹员，不是这里的“职务”。但仓库管理员、产品质检员，具有一定管理权限，属于这里的“职务”。（2009年第18题，2014年第17题）

例2，甲是某公司司机，按照公司安排独自一人将价值7万元的货物运往外地。甲在途中将货物变卖，携款潜逃。长途运输司机单独运输，没有押运员，表明公司对司机有充分授权，司机在占有公司财物。甲构成职务侵占罪。（2008年四川卷第18题）

2.“利用职务便利”中的“利用”

这是指实质的利用而非形式的利用，也即要求是利用职务上主管、管理单位财物的便利，而不是仅仅利用因工作关系而熟悉作案环境、易于接近作案目标、容易进入作案场所。

例如，某民营公司财务总监甲利用权限登录公司账户，盗划一笔资金。这是实质利用职权，构成职务侵占罪。甲又来到公司仓库，库管员认识甲，知道甲是本公司的财务总监，便没多问，甲进去后偷拿一个货物。这是形式利用职权，仅构成盗窃罪，不构成职务侵占罪。

% Source PDF physical page 299; see audit/source_coverage.json.
\section{普通罪名}

\subsection*{一、故意毁坏财物罪}

第275条故意毁坏公私财物，数额较大或者有其他严重情节的，处三年以下有期徒刑、拘役或者罚金；数额巨大或者有其他特别严重情节的，处三年以上七年以下有期徒刑。

1.毁坏既包括物理性毁坏，也包括功能性毁坏。功能包括客观功能，也包括一般人主观认可的功能。

2.即未使物理性毁坏和功能性毁坏，但使所有权人丧失占有，也构成毁坏。例如，打开他人鸟笼，放飞笼中小鸟。又如，打开他人鱼塘，让鱼回归大海。再如，将他人戒指扔到大海里。（2006年第60题）

\subsection*{二、挪用资金罪（第272条）}

这是指，公司、企业或者其他单位的工作人员，利用职务上的便利，挪用本单位资金归个人使用或者借贷给他人，\CJKunderwave{数额较大（5万元）、超过3个月未还的}，或者虽未超过3个月，但\CJKunderwave{数额较大（5万元）、进行营利活动的}，或者\CJKunderwave{进行非法活动（要求数额达到3万}\CJKunderwave{元）的}行为。

1.如果行为人挪用单位资金后，故意不归还，则行为性质转化为职务侵占罪。

2.对于挪用单位资金进行非法活动构成其他犯罪的，应当实行数罪并罚。

3.《刑法修正案（十一）》增设一款规定：成立本罪，在提起公诉前将挪用的资金退还的，可以从轻或者减轻处罚。其中，犯罪较轻的，可以减轻或者免除处罚。

\subsection*{三、挪用特定款物罪}

第273条挪用用于救灾、抢险、防汛、优抚、扶贫、移民、救济款物，情节严重，致使国家和人民群众利益遭受重大损害的，对直接责任人员，处三年以下有期徒刑或者拘役；情节特别严重的，处三年以上七年以下有期徒刑。

1.行为主体是自然人（直接责任人员），包括国家工作人员和非国家工作人员。\footnote{这是法考多数说。少数说认为，本罪的行为主体是单位，但只处罚直接责任人员。}

2.行为方式是指有关单位的直接责任人员改变特定款物的专用用途。如果国家工作人员利用职务便利，挪用特定款物归个人使用，则构成挪用公款罪，并从重处罚（第384条第2款）。

\subsection*{四、拒不支付劳动报酬罪}

第276条之一，以转移财产、逃匿等方法逃避支付劳动者的劳动报酬或者有能力支付而不支付劳动者的劳动报酬，数额较大，经政府有关部门责令支付仍不支付的，处三年以下有期徒刑或者拘役，并处或者单处罚金；造成严重后果的，处三年以上七年以下有期徒刑，并处罚金。

% Source PDF physical page 300; see audit/source_coverage.json.
单位犯前款罪的，对单位判处罚金，并对其直接负责的主管人员和其他直接责任人员，依照前款的规定处罚。

有前两款行为，尚未造成严重后果，在提起公诉前支付劳动者的劳动报酬，并依法承担相应赔偿责任的，可以减轻或者免除处罚。

1.本罪是不作为犯罪，即不履行支付劳动报酬的义务。

2.成立本罪有个前置条件：经政府有关部门责令支付仍不支付。（2012年第1题）

3.本罪的从宽处罚条件：尚未造成严重后果，在提起公诉前支付劳动者的劳动报酬，并依法承担相应赔偿责任。

\section{财产犯罪的既遂}

\subsection*{一、既遂标准}

取得型财产犯罪（盗窃罪、诈骗罪、敲诈勒索罪、抢劫罪、抢夺罪）要构成既遂，要求实现非法占有目的，也即要求建立自己的占有。因此，这类犯罪的既遂标准是取得控制财物（取得控制说）。\footnote{注意1：2008年以前的法考，盗窃罪的既遂标准是失控说。2008年后已经改为取得控制说。参见张明楷：《刑法学》（第六版），法律出版社2021年版，第1256页。注意2：司法解释的特殊规定。抢劫罪有个司法解释规定“即使行为人没抢到财物，只要造成被害人轻伤”也构成抢劫罪既遂。然而，第一，虽然抢劫罪的保护法益有人身和财产，但二者地位不是相等的，财产法益具有决定性。抢劫罪的属性是取得型财产犯罪，因此只有取得控制了财物才构成既遂。第二，甲抢劫乙，使用暴力先将乙打成轻伤，压制反抗，然后取得财物。按照该司法解释，甲将乙打成轻伤时就既遂了，那么取得财物就不能算既遂，因为一个犯罪只能既遂一次。然而，取得财物不算既遂，显然不妥。因此，只能将该司法解释的规定视为特殊规定。考试考出来了，就按该司法解释答题。}

犯罪人是否取得控制财物，主要考察犯罪人对财物的控制力大小。控制力越大，越倾向得出既遂结论。控制力的程度，至少要求排除被害人的支配。\footnote{注意：主人占有自己的财物，要求财物在自己的实际控制范围内，但不要求具有排除他人支配的可能性。例如，甲坐火车，将提包放在行李架上。虽然其他人能够拿走该提包，但是该提包属于甲在占有。}现以盗窃罪为例予以说明。

\subsection*{二、具体分析}

1.小件物品

[总结]装进衣服口袋的问题。

（1）小偷在大街上、公交车上、商场盗窃，将小件物品装进口袋，构成既遂。这是因为，普通公民没有搜身权，财物一旦进入小偷口袋，进入了个人专属领域，便处于小偷的控制之下。（2013年第54题）

（2）小偷在超市内盗窃，将小件商品装进口袋，没出收银台，构成既遂。不过，如果小偷的举动一直受到超市人员的监控，超市人员已经布控，则小偷只有将商品拿出收银台才能构成既遂。

（3）小偷入户盗窃，在主人家里将小件物品装进口袋，不算既遂。这是因为，主人对进入家里的小偷有搜身权和控制权。

（4）保姆、客人在主人家里盗窃，将小件物品装进口袋，均构成既遂。这是因为，主人不能对保姆、客人搜身。

% Source PDF physical page 301; see audit/source_coverage.json.
2.大件物品

例1，甲到家电商场盗窃电冰箱，将电冰箱拖出专卖柜台时不是既遂，拖出商场大门时才是既遂。

例2，甲进入某工厂仓库盗窃大件物品。如果工厂大门有严格门禁制度，则偷出大门才算既遂。如果工厂大门随便进出，形同虚设，则偷出仓库门就算既遂。

例3，乙发现甲下车没有拔钥匙，便悄悄上车，欲开走车。甲发现，马上制止并抓住了乙。乙构成盗窃罪未遂。只有开走汽车（移动了汽车）才算既遂。

3.在空间上，行为人建立自己的占有，既包括紧密占有，也包括松弛占有。也即，只要财物处于自己实际控制范围内即可，可以与之保持相当的空间距离。

[总结]放在秘密地点问题。一般而言，犯罪人将财物放在秘密地点，就既遂。犯罪人指示被害人将财物放在公共场所的秘密地点，不算既遂，因为被害人具有取回财物的能力，犯罪人未达到排除被害人支配的效果；犯罪人取得财物才算既遂。犯罪人指示被害人将财物放在自已控制的秘密地点，就算既遂。

例1，甲在火车上趁对面乘客睡觉，将其行李箱扔出车厢，扔到自已村头预定的地点草丛中，就属于既遂。如果扔到了长江里，只能是未遂。如果甲将财物扔到僻静处，已经既遂后，下车后回头捡时，发现被别人捡走，仍然是既遂，不能变回未遂。

例2（2008年第6题），深夜，甲入室盗窃，窃得主人的皮箱，扔到院墙外的僻静处。“深夜僻静处”，表达的意思是，此刻这个地方很少有人路过，因此，视为甲将财物置于自已的实际控制范围，甲构成既遂。即使此后意外地被别人捡走，仍然是既遂，不能变回未遂。既遂排斥未遂。

[巩固练习]题1，保姆将主人戒指藏在卫生间夹缝里，准备下周辞职时拿走，次日便被抓。保姆构成既遂。

题2，甲敲诈勒索乙，要求将一万元现金放进街边垃圾桶里。乙将现金放进垃圾桶里，甲因为害怕未前往取现金。后乙将现金取回。甲是否构成敲诈勒索罪既遂？（2021年主观题）

（1）少数说：甲构成犯罪既遂。理由：垃圾桶是个秘密地点，乙将现金放进垃圾桶，甲便取得控制了现金。

（2）多数说：甲构成犯罪未遂。理由：犯罪人甲指示被害人乙将财物放在公共场所的秘密地点，乙具有取回财物的能力，甲对财物的控制力没有达到排除乙支配的效果。甲只有拿到手才算既遂。\footnote{参见张明楷：《刑法学》（第六版），法律出版社2021年版，第485页。}

4.因果关系问题。行为人取得财物与被害人遭受财产损失的关系。行为人取得了财物，必然导致被害人遭受财产损失。但是，被害人遭受财产损失，并不意味着行为人取得财物。

例如（2015年第5题），甲诈骗乙，乙受骗，给甲汇钱，但是搞错卡号，将钱汇到了丙的账户。这时候，甲没有取得财物，构成诈骗罪未遂。不过，乙的财产损失与甲的诈骗行为仍有因果关系，甲应负赔偿责任。

% Source PDF physical page 302; see audit/source_coverage.json.
\section{财产犯罪的数额}

\subsection*{一、数额档次}

根据财物的价值数额大小，刑法保护财物的档次如下：

第一档：价值数额较大，值得刑法完整保护。所谓完整保护，是指所有的财产犯罪都要保护。数额较大一般指2000元左右。例如，商场发行的购物卡，不记名、不挂失。存有3000元的购物卡本身就是价值数额较大的财物。盗窃这样的购物卡，就等于盗窃了价值数额较大的财物，构成盗窃罪。

第二档：价值数额不大，但值得刑法有限保护。所谓有限保护，是指不是所有财产犯罪都保护，只有四个罪名保护：抢劫罪，特殊类型的盗窃罪（入户盗窃、携带凶器盗窃、扒窃、多次盗窃），多次抢夺，多次敲诈勒索。例如，一辆价值300元的自行车、一张银行卡、存折、欠条，属于价值数额不大，但是值得刑法有限保护的财物。普通盗窃一辆300元的自行车，不构成盗窃罪，按照《治安管理处罚法》处理。多次盗窃300元的自行车，则构成盗窃罪。抢劫300元的自行车，则构成抢劫罪。（2011年第16题）

例1，欠条只是证明债权的一种证据，不等于债权（财产性利益）本身，记载3000元债权的欠条不等于价值3000元的财物，因此，用普通盗窃方式盗窃一张欠条或诈骗一张欠条，不构成盗窃罪、诈骗罪。但欠条属于虽然价值数额不大，但是值得刑法有限保护的财物。因此，抢劫欠条或入户盗窃欠条，构成抢劫罪或盗窃罪。

例2，银行卡的特征是实名制、设有密码、可挂失。盗窃了银行卡，并不能排他的控制支配卡里的资金，并不能对卡里资金建立自己的占有。因此，盗窃了存有3000元的银行卡，不等于盗窃了价值3000元的财物。因此，用普通盗窃方式盗窃银行卡或诈骗银行卡，不构成盗窃罪或诈骗罪，只有此后使用卡，花了其中的钱，才可能构成犯罪。不过，银行卡虽然价值数额不大，但值得刑法有限保护，因此，抢劫或入户盗窃一张银行卡，构成抢劫罪、盗窃罪。

第三档：价值很小，不值得刑法保护。例如，一张餐巾纸、一张普通名片、一根葱等，由于价值极其低廉，不是财产犯罪的对象。即使抢劫这些财物，也不构成抢劫罪；如果致人轻伤，可定故意伤害罪。

\InsertDiagram{figures/manual_table_p302_1.tex}

% Source PDF physical page 303; see audit/source_coverage.json.
\ContinuedTable{figures/manual_table_p303_1.tex}
\SourceFootnote{\textbf{表注①}\quad 根据司法解释规定，多次盗窃、多次抢夺、多次敲诈勒索，是指两年内实施三次。刑法在诈骗罪中没有规定，“多次诈骗的，构成诈骗罪”。因此，两年内诈骗三次他人价值300元的自行车，不构成诈骗罪。但是，如果实施十次这样的诈骗（不要求两年内实施），骗取财物的数额累计达到3000元（诈骗罪的定罪数额“数额较大”），那么构成诈骗罪。假如每次仅诈骗一根葱，诈骗了三万次，累计数额达到3000元，仍不构成诈骗罪，因为一根葱不值得刑法保护，可以用治安管理处罚法处罚这种行为。}

\subsection*{二、主客观相一致问题}

盗窃罪（暂不考虑特殊盗窃）、抢夺罪、诈骗罪、敲诈勒索罪的定罪数额是“数额较大的财物”，适用基础法定刑，量刑数额有“数额巨大的财物”“数额特别巨大的财物”，适用升格法定刑。以盗窃罪为例，司法解释规定的“数额较大”是2000元左右，“数额巨大”是3万元以上，“数额特别巨大”是30万元以上。

由于这些财产犯罪都是故意犯罪，作为行为对象，\CJKunderwave{“数额较大（巨大、特别巨大）的}\CJKunderwave{财物”是个主客观相一致的要素}。常见误解是，将“数额较大（巨大、特别巨大）的财物”理解成纯粹客观要素，只看客观上是否具备。

\InsertDiagram{figures/manual_table_p303_2.tex}

% Source PDF physical page 304; see audit/source_coverage.json.
\ContinuedTable{figures/manual_table_p304_1.tex}
\SourceFootnote{\textbf{表注①}\quad 本案与共同犯罪一讲的渔网案的区别在于，本案中，甲对3万元的珠宝产生了紧迫危险，属于着手实行，所以定未遂。渔网案中，甲对预定目标尚未产生紧迫危险，不属于着手实行，所以定犯罪预备。}

[巩固练习］题1，孙某敲诈勒索赵某，要求赵某交付价值800万元的名画，赵某将价值8千元的一幅画谎称是800万元的画交给孙某。（1）多数说：孙某的一个行为一方面构成敲诈勒索数额特别巨大财物的未遂，另一方面构成敲诈勒索数额较大财物的既遂，想象竞合，择一重刑论处。（2）少数说：孙某只构成敲诈勒索数额较大财物的既遂。（2016年主观题）

题2，王新明使用伪造的户口本、身份证，冒充房主即王新明之父的身份，以出售一处房屋为由，与徐某签订房屋买卖合同，约定购房款为100万元，收取徐某支付的购房首付款30万元，并约定余款过户后给付。后王新明的虚假身份被建委工作人员发现，余款未取得。王新明骗到30万元（数额巨大），既遂数额是30万元（数额巨大）；剩下70万元未骗到，未遂数额是70万元（数额特别巨大）。多数说认为，王新明一个行为同时触犯诈骗数额巨大的既遂和诈骗数额特别巨大的未遂，想象竞合，择一重刑论处。最高人民法院的裁判要旨就是这个意思。\footnote{最高人民法院指导案例62号：王新明合同诈骗案。裁判要旨：“在数额犯中，犯罪既遂部分与未遂部分分别对应不同法定刑幅度的，应当先决定对未遂部分是否减轻处罚，确定未遂部分对应的法定刑幅度，再与既遂部分对应的法定刑幅度进行比较，选择适用处罚较重的法定刑幅度。”}

% Source PDF physical page 305; see audit/source_coverage.json.
题3（2017年第5题），王某欲通过合同诈骗的方式骗取徐某钱财。王某与徐某的房屋签约价是220万元，徐某首付100万元，之后王某被抓。王某的一个行为同时触犯诈骗“数额特别巨大”（100万元）的既遂和诈骗“数额特别巨大”（120万元）的未遂，想象竞合，择一重刑论处，以“数额特别巨大”（100万元）的既遂论处。

[提示］未遂数额是未获得的数额，是欲骗数额减去获得数额，而不是欲骗数额。在

题2中，未遂数额是100万元-30万元=70万元，不是100万元。题3中，未遂数额是220万元-100万元=120万元，不是220万元。\footnote{提示：不确切的期待数额，不能作为犯罪数额。例1（2013年第6题），甲偷到一个玉坠，认为是根本不值钱的仿制品，想骗乙的钱，谎称是文物，出价5万元。乙以为是文物，还价3万元。3万元成交。日常交易中，卖家有意抬高出价，买家再杀价，是正常现象。甲虽然出价5万元，但对5万元并不抱确切认真的希望，而是先把价抬高，能骗多少是多少，是一种不确切的心理期待。这种出价数额不能作为犯罪数额来考虑。对甲只需要根据实际获得数额加以认定即可。甲构成诈骗罪，既遂数额是3万元。在此不能认为，甲的未遂数额是2万元。本案中的5万元数额与王新明案中的签约数额不同，后者是被害人答应的确切数额，被害人有遭受这个数额损失的现实危险性。例2，狗蛋听说小芳家挺有钱的，便入室盗窃小芳家，抱着“能偷多少是多少，越多越好，能偷100万元的话，下辈子就不愁了”的心理，翻箱倒柜，偷到5000元。在此，不能把狗蛋心里的“100万元”当回事。对盗窃罪的数额按照实际获得数额认定即可，也即狗蛋构成盗窃罪，既遂数额是5000元。在此不能认为，未遂数额是995000元。}

\subsection*{三、违禁品的数额}

财物的价值包括积极价值和消极价值。消极价值，是指不能给人或社会带来积极价值，但落人不法分子之手会给他人或社会带来财产损失或其他危害。例如，违禁品（如毒品、假币、淫秽物品）有消极价值，落入不法分子手里会危害社会，在不法分子手里也能卖钱，所以属于有价值的财物。这些有消极价值的财物，能够成为财产犯罪的对象。

在计算违禁品的价值数额时，可参考其非法交易的价值数额（行情价）。根据违禁品的具体价值数额来确定其处在三个档次的哪个档次，依此判断犯罪成立条件和既遂条件。\footnote{有人可能认为，司法解释规定，对于违禁品按照情节来量刑。实际上，对于违禁品按照情节来量刑，不影响计算其消极价值的大小。何况，许多情况下必须计算其价值大小。例如，盗窃罪、抢夺罪、诈骗罪、敲诈勒索罪的第一档法定刑条件是“数额较大”，第二档是“数额巨大或者有其他严重情节”，第三档是“数额特别巨大或者有其他特别严重情节”。由于第一档没有规定“或者其他较重情节”，所以在衡量第一档时，需要计算违禁品的价值数额是否较大。}

例1（2012年第54题），乙抢夺邹某现金20万元，后发现全部是假币。20万元的假币具有消极价值，该价值可以根据非法交易的行情价来衡量，至少达到2000元左右，也即“数额较大”。因此，乙构成抢夺罪既遂。如果乙想抢夺数额较大的财物，抢夺到对方的提包，发现里面只有500元假币，由于500元假币的价值没有达到“数额较大”，因此乙构成抢夺罪未遂。

例2，甲想盗窃乙的100克毒品，并且盗窃到手。经评估，100克毒品价值1万元（非法交易的行情价）。甲构成盗窃罪既遂。如果甲盗窃到手的只有1克（价值100元），则构成盗窃罪未遂。

[巩固练习]题1，甲想抢劫或入户盗窃乙的5克毒品，并且取得到手。经评估，5克毒品价值500元（非法交易的行情价），属于数额没有达到较大、但也不是价值微薄的毒品，属于第二档次的财物。甲构成抢劫罪既遂或盗窃罪既遂。

题2，甲想盗窃乙的一箱淫秽光盘（非法交易行情价是3000元），盗窃到手的只有一张淫秽光盘，构成盗窃罪未遂。如果甲只想盗窃或抢劫乙的一张淫秽光盘，则属于盗窃或抢劫价值微薄的财物，不构成盗窃罪或抢劫罪。

% Source PDF physical page 306; see audit/source_coverage.json.
\subsection*{四、犯罪数额的认定}

1.犯罪所得数额（既遂数额）以犯罪时（行为时）为准。

例如（2008年第15题），甲想抢夺他人价值数额较大的财物，抢夺到手后误以为根本不值钱，实际价值数额较大。甲便丢弃。“抢夺数额较大财物”是个主客观相一致的要素。“主客观相一致”原则中的“主观认识”是指行为时的主观认识，而非结果时的主观认识。甲在行为时的主观认识是“他人的财物是数额较大的财物”，客观上的确也是数额较大的财物，主客观是相一致的。并且，甲将该财物抢夺到手，因此构成犯罪既遂。在结果时的主观认识，甲认为“这是根本不值钱的财物”，这种认识与客观事实不一致，这种主客观不一致不重要。不能因此认为甲构成犯罪未遂。

2.犯罪所得的收益的数额以犯罪后案发时为准。

例如，甲盗窃了乙的10万元，然后用来炒股，赚了1万元，后又赔了5000元，截至案发时经结算，共赚了2万元。犯罪所得数额是10万元。犯罪所得的收益是2万元。这一点在受贿罪中经常考查。例如，甲送官员价值100万元的股份，日后涨到150万元。受贿数额是100万元，多的50万元是收益数额。（2020年试题）

\section{财产性利益}

财产性利益是指有形财物以外的具有财产性价值的利益。这是考试的重点和难点。希望勇敢的你耐心学习，多看两遍。

\subsection*{一、总结债权问题}

\subsubsection*{（一）类型一：行为人是债务人，被害人是合法债权人}

1.\CJKunderwave{用诈骗、抢劫方式，使债权人放弃行使债权，构成诈骗罪、抢劫罪}。

例1（骗免过路费案），甲使用假军车牌，骗免过路费，收费员见甲的车是军车，便免收过路费。甲通过欺骗方式使收费员放弃行使债权，属于诈骗财产性利益，构成诈骗罪。

例2（抢劫车费案），王某乘坐出租车，到目的地后，司机要求支付车费。王某用刀威胁司机：“要车费，我就要你的命！”司机被迫说道：“算我倒霉，你下车吧！”王某下车离开。王某通过强制手段使司机放弃行使债权，属于抢劫财产性利益，构成抢劫罪。（2002年第12题，暂不考虑数额）\footnote{此处“财产性利益”的含义。合法债权人只是放弃行使债权，在民法上并未丧失债权，日后可以向法院起诉并主张债权。那么，所谓合法债权人损失了“财产性利益”，是指损失了什么？合法债权人放弃了当场行使债权。“债权人当场实现债权”就是一种财产性利益。当场收不到钱，就是一种经济损失。相应的，债务人获得的财产性利益，并非获得债权，或者在民法上免除债务，而是延迟履行债务。延迟履行债务对债务人而言，也是一种经济利益。拖欠不还，对债务人而言也是一种好处。同理，甲盗窃了乙的轿车，一般笼统地说，甲侵犯的法益是乙的所有权。然而，在民法上，乙没有丧失所有权。所以，准确地说，甲侵犯的法益是乙对所有权的正常行使，也即无法正常用车了。“所有权的正常行使”对乙而言就是一种经济利益。}

[注意】“转移占有”的问题。

少数说认为，上文例1、例2的行为人不构成诈骗罪、抢劫罪。理由是，诈骗罪、抢劫罪的成立，都要求将他人占有的财物转移为自己占有。而这些案件中，行为人并未将他人的债权转移为自己占有。

% Source PDF physical page 307; see audit/source_coverage.json.
多数说认为，“转移占有”是针对有形财物要求的。对于侵犯他人财产性利益，无法要求物理位移性质的转移占有，只要求：一方面他人丧失了财产性利益，另一方面自己获得了财产性利益，并且二者之间具有两种关系，一是条件关系，也即若没有行为人的行为，被害人就不会遭受损失；二是加减关系，也即行为人获得的利益就是被害人损失的利益。

[巩固练习]题1，甲欠乙50万元。乙讨债，甲杀死乙。多数说认为，甲使乙无法实现债权，属于抢劫财产性利益，构成抢劫罪（故意）致人死亡。少数说认为，甲不构成抢劫罪，只构成故意杀人罪。

题2（2016年主观题），赵某在家保管钱某的一幅名画，钱某讨要，赵某为了免除返还义务，杀死钱某。多数说认为，赵某构成抢劫罪。理由是，赵某使用暴力使钱某无法实现债权，属于抢劫财产性利益。赵某构成抢劫罪（故意）致人死亡。少数说认为，赵某不构成抢劫罪，只构成故意杀人罪。

2.\CJKunderwave{用盗窃方式，逃避债务}。

例如（餐厅逃单案），甲在餐馆开始有付费意思，吃完后，店家让买单时，甲不想付费，从后门偷偷溜走。甲的行为属于盗窃行为，盗窃财产性利益。对此存在观点展示。

多数说认为，甲不构成盗窃罪。理由有：第一，店家没有放弃行使债权。这一点与上文“骗免过路费案”“抢劫车费案”不同。后者中，债权人放弃了行使债权。第二，这种逃单行为按照民事纠纷处理即可。刑法在此应保持谦抑性。

少数说认为，甲构成盗窃罪。理由是，甲的行为导致店家无法当场行使债权。“当场行使债权”是店家的财产性利益。店家损失了这份利益。甲逃避债务，获得了延迟履行债务的利益。

[巩固练习］甲驾车到了收费站，不想付过路费，迅速跟着前车逃出收费站，或者冲撞栏杆，逃出收费站。甲属于盗窃财产性利益或者抢夺财产性利益。对此有多数说（不构成盗窃罪或抢夺罪）和少数说（构成盗窃罪或抢夺罪）。（2021年试题）

[总结]逃费行为

\InsertDiagram{figures/manual_table_p307_1.tex}

% Source PDF physical page 308; see audit/source_coverage.json.
\subsubsection*{（二）类型二：行为人是债务人，被害人是非法债权人}

问题：用非法方式，使非法债权人放弃行使非法债权，是否构成犯罪？

例如（骗免资案），甲与卖淫女乙约定资是3千元，事后，甲用假币支付资。甲构成使用假币罪。甲是否构成诈骗罪？

法律的财产说（少数说）：刑法只保护民法上的合法财产权。乙的债权不是民法上的合法财产权。甲不构成诈骗罪。

经济的财产说（少数说）：只要是具有经济价值的利益，刑法就应保护。乙的债权具有经济利益（“血汗钱”）。甲构成诈骗罪，与使用假币罪想象竞合，择一重罪论处。

法律与经济的财产说（多数说）：刑法保护的是法秩序认可的、具有经济价值的利益。法秩序不认可乙的非法债权。例如，乙不享有向法院起诉的诉权。法律程序不保护乙的非法债权。甲不构成诈骗罪。

[巩固练习］欺骗分赃案：苏某与王某共同抢劫了周某的一部手机。二人约定销赃款平分，由苏某负责销赃。苏某销赃获得赃款1万元。随后，苏某对王某谎称手机仅卖得4000元，分给王某2000元。苏某对王某是否构成诈骗罪？法律与经济的财产说（多数说）认为，苏某不构成诈骗罪，因为法秩序不认可王某对销赃款的请求权。这表明，这种类型的“黑吃黑”不构成犯罪。

\subsubsection*{（三）类\CJKunderwave{型三}：行为人与被害人没有债权债务关系}

\CJKunderwave{1.用非法方式，为自已生成非法债权，使被害人成为债务人。结论：该行为构成财}\CJKunderwave{产犯罪}。若给被害人造成实际损失，则构成既遂。

例1（强迫写欠条案），甲持刀逼迫乙写下“乙欠甲1万元”的欠条，甲拿走欠条。后甲未能实现债权。甲构成抢劫罪，属于未遂。\footnote{最高人民法院（2000）刑他字第9号批复：被告人以暴力、胁迫手段强行夺回欠款凭证（欠条），并让债权人在被告人已写好的收条上签字，以消灭其债务的行为，符合抢劫罪的特征，应以抢劫罪定罪处罚。第一个行为是抢劫欠条本身，成立抢劫罪。第二个行为是逼迫写收条，强迫债权人放弃行使债权，属于抢劫财产性利益，成立抢劫罪。两个抢劫罪是连续犯关系，只处罚一个抢劫罪即可。}

例2（非法充值案），甲在某电影院办了一张会员卡，采用技术手段非法充值1万元，然后使用，导致电影院损失1万元票房收入。甲非法充值，使自己享有了针对电影院的1万元非法债权，使电影院成为债务人。甲构成盗窃罪，并且既遂。\footnote{行为与行为对象同时存在原则。这是指，在实施行为时，行为对象应已然存在。例如，上文“类型一”“类型二”中，行为人侵犯被害人的债权时，被害人的债权已然存在。然而，上文“类型三”中，行为人不是侵犯被害人已然存在的债权，而是为自已生成非法债权。此时非法债权不属于行为对象，而属于行为滋生之物，所以在此不需要遵守“行为与行为对象同时存在原则”。例如，在贩卖毒品罪中，“毒品”是行为对象。在制造毒品罪中，“毒品”是行为滋生之物。债权既可以是行为对象，也可以是行为滋生之物，是因为，债权属于请求权，既可以已然存在，然后被侵犯，也可以通过单方面请求而产生，也即可以“无中生有”。}

2.\CJKunderwave{非法使用被害人的财产}。行为人既没有侵犯被害人已然存在的债权，也没有为自已生成债权，则不构成侵犯财产性利益的犯罪。

例如（2013年第61题），甲得知邻居出国，撬开门，对乙谎称自已是房东，将房屋出租给乙。甲对邻居不构成盗窃财产性利益。这是因为，第一，在甲撬门进去时，邻居并未已然拥有财产性利益（租金收入）。第二，甲并未生成债权，并未使邻居成为债务人。因此，甲对邻居不构成盗窃罪。当然，甲的侵入行为构成民事侵权，所得租金收入应按照

% Source PDF physical page 309; see audit/source_coverage.json.
不当得利返还给邻居。若拒不返还，则构成侵占罪。\footnote{（1）多数说认为，甲非法进入邻居家的行为不构成非法侵入住宅罪，因为该罪的保护法益是居住的安宁状态。屋内长期无人，甲没有侵犯居住的安宁状态。（2）甲对乙是否构成诈骗罪，关键看乙的损失情况。如果乙先住，甲后收钱，则乙不会有损失。甲对乙不构成诈骗罪。如果甲先收一年钱，那么乙存在损失的危险，例如房东随时回来，则甲构成诈骗罪；如果乙住够租期，没有损失，则甲构成诈骗罪未遂；如果房东提前回来，乙住不下去，损失了3千元支付出去的租金，则甲构成诈骗罪既遂。}

\CJKunderwave{3.冒名行使债权}。被害人对第三人享有债权。行为人在被害人不知情的情况下，代为行使债权，导致被害人丧失债权。结论：多数说认为构成三角诈骗，少数说认为构成盗窃罪。

例如（冒名收款案），狗蛋是甲公司的业务员，一直负责向乙客户收款。甲公司辞退了狗蛋，但未将该情况及时通知乙客户。狗蛋谎称代表公司来收款，乙客户不知情而将款项给了狗蛋。

做题方法：这种题有个三角关系，犯罪人一债权人一债务人。做题第一步，在债权人和债务人中找出被害人，也即遭受财产损失的人。对此判断，需要看债务人是否有效履行了债务。如果算是履行了债务，则债务人不是被害人，那么债权人就是被害人。本案中，由于狗蛋的收款行为具有权利外观，形成表见受领，乙客户算是履行了付款义务，没有财产损失，不是被害人。甲公司不能再向乙客户行使债权，但钱款又没收到，因此是被害人。

少数说：狗蛋对甲公司构成盗窃罪，在甲公司不知情的情况下代为行使了甲公司针对乙客户的债权，属于盗窃债权（财产性利益）。

多数说：狗蛋构成新型三角诈骗，受骗人是乙客户，被害人是甲公司，受骗人与被害人不是同一人，所以是三角诈骗。之所以是新型三角诈骗，是因为，传统三角诈骗中，受骗人处分的是被害人的财物，而新型三角诈骗中，受骗人处分的是自己的财物。也即，受骗人乙客户处分的是自己的财物。

[巩固练习］偷换二维码案：甲偷偷将某餐厅收银台上的二维码调换成自己的二维码。某顾客在某餐厅用餐后买单，餐厅收银员让顾客将餐费打入指定的二维码账户（实际是甲的二维码账户），顾客照办。甲的账户收到顾客的餐费。顾客将餐费打入该账户，就履行了付款义务，餐厅无权要求顾客再次付款，因此顾客没有遭受损失，不是被害人。餐厅不能再向顾客主张债权，钱又没收到，是被害人。少数说：甲在餐厅不知情的情况下，代为行使了餐厅对顾客的债权，盗窃了餐厅的债权，并实现了债权（获得一笔钱款），构成盗窃罪（既遂）。多数说：甲构成新型三角诈骗，受骗人是顾客，被害人是餐厅，受骗人与被害人不是同一人，所以是三角诈骗。之所以是新型三角诈骗，是因为，受骗人顾客处分的是自己的财物。

\subsection*{二、总结存款问题}

储户将现金存人银行，现金归银行所有、占有，储户换来针对银行的债权，称为存款债权。这种存款债权，被称为卡里“资金”，也即已经被实体化。这种实体化的债权与普通债权具有不同特征。普通债权，例如乙欠甲钱，甲对乙享有债权，甲对该债权无法自由地转移、转让。而存款债权，就像有体物一样，可以自由地转移、转让。例如，甲可以自

% Source PDF physical page 310; see audit/source_coverage.json.
由地将自己银行卡里的“资金”划到自己另一个卡里或者另一个人的卡里。

因此，在分析侵犯这种存款债权的案件时，可以将其理解为有形财物，使用“占有”及“转移占有”的概念进行分析。卡主（存款人）在占有存款债权（卡里资金），不占有存款指向的现金，现金由银行所有并占有。例如（非法转账），甲利用技术手段将乙存折里的资金（实际是指存款债权）划到自已存折里，对这笔资金实施了转移占有，构成盗窃罪，盗窃存款债权。\footnote{如果甲利用技术手段将乙银行卡里的资金划到自己账户，按照2018年司法解释，应定信用卡诈骗罪。这是特殊规定。对此，在后文“信用卡诈骗罪”中有专门讲解。}

\subsubsection*{（一）卖卡后，针对合法资金，掐卡、转账、取款}

模型：甲将自已的银行卡卖给（或借给）乙。日后，甲发现该卡中乙存入1万元资金。首先，乙是该资金的所有权人。争议焦点是，谁是该资金的占有人？

\CJKunderwave{事实占有说}（多数说）：乙在事实上占有该资金，甲没有占有该资金。这是实质标准。

\CJKunderwave{法律占有说}（少数说）：甲是法律上、名义上的卡主，甲在占有该资金。这是形式标准。

1.事实占有说的看法（多数说）

（1）甲掐卡，也即凭身份证到银行柜台挂失该卡。多数说认为，该掐卡行为（挂失卡）构成盗窃罪。理由是，甲的掐卡行为导致乙无法支配卡内资金，使乙对卡内资金失去占有，并且甲对卡内资金建立了排他性占有。\footnote{参见张明楷：《供卡人掐卡、取款的行为性质》，载《法学评论》2024年第1期，第1页。}（2025年试题）

（2）甲掐卡后转账，也即将卡内1万元资金转入自己另一账户。多数说认为，该行为不再构成犯罪。理由是，掐卡行为便构成盗窃罪既遂，也即已经占有了卡内资金，然后再转账，属于将自己左口袋财物放进自己右口袋。

（3）甲掐卡后取款，也即让银行职员取出1万元现金。多数说（事实占有说）认为，甲对银行构成诈骗罪，诈骗的行为对象是银行的现金，因为甲不是资金的所有权人，没有取款权限。\footnote{参见张明楷：《领取无正当原因汇款的行为性质》，载《法学》2020年第11期，第18页。此外，甲不构成信用卡诈骗罪，因为甲不属于“冒用他人的信用卡”。提示：银行只是直接的受害人。银行会根据合同规则，扣掉储户的存款债权，由此挽回损失。银行不是最终的受害人。但是，刑法主要关注的是直接受害人及直接损失。}少数说（法律占有说中的一种观点）认为，甲对银行不构成诈骗罪，因为甲是资金的法律占有人，有取款权限。若在ATM机取现金，多数说认为，甲对银行构成盗窃罪，盗窃银行的现金；少数说认为，甲对银行不构成盗窃罪。\footnote{甲的取款行为也导致用卡人乙对卡里“资金”彻底丧失，但是由于甲的掐卡行为已经导致乙对卡里“资金”失去占有，构成盗窃罪，所以，甲的取款行为针对乙不需要再定罪。}

2.法律占有说认为（少数说），对掐卡行为与之后的转账、取款行为应一体看待，甲的这些行为属于将乙所有、自己占有的资金变成自己所有，构成侵占罪。

\subsubsection*{（二）卖卡后，针对赃款，掐卡、转账、取款}

模型：甲向乙卖卡，知道乙用于电信网络诈骗。某日，甲得知卡里多了1万元，是乙诈骗所得赃款。甲先掐卡（挂失卡），然后转账、取款。

1.甲出售卡的行为构成诈骗罪的帮助犯、帮助信息网络犯罪活动罪，想象竞合，择一重罪论处。

% Source PDF physical page 311; see audit/source_coverage.json.
2.甲的掐卡、转账、取款行为。

（1）事实占有说的看法（多数说）

将掐卡与转账、取款一体化看待，构成盗窃罪。\footnote{2023年4月11日两高《关于印发办理涉“两卡”犯罪案件典型案例的通知》（案例三）（高检发办字

[2023】43号），“吕某钊、宁某帮助信息网络犯罪活动盗窃，吕某钊掩饰、隐瞒犯罪所得案”。由于这是两高发布的典型案例，所以是多数说。此外，甲不构成信用卡诈骗罪，是因为甲是卡的名义人，不是在冒用他人的银行卡。}这是因为，第一，刑法保护乙对卡里“资金”的非法占有事实。这种非法占有事实能够对抗第三人。第二，甲只是诈骗罪的帮助犯，没有事实占有赃款，乙在事实占有赃款。\footnote{乙构成诈骗罪既遂，甲构成乙的诈骗罪的帮助犯，因此也构成既遂。甲的这个既遂的含义是指，甲对乙的诈骗罪的既遂结果要负责，亦即也要承担既遂的刑事责任。这并不意味着，甲也在现实的占有犯罪所得。共犯人中，谁在现实占有犯罪所得，需要根据事实进行具体判断。}刑法不允许甲对乙的“黑吃黑”。（2024年试题，2023年主观题，2025年主观题）

(2）法律占有说的看法（少数说）

将掐卡与转账、取款一体化看待，构成侵占罪。这是因为，甲作为卡主在法律上占有卡里资金，甲将他人所有、自己占有的资金变成自己所有，构成侵占罪。\footnote{还有第二种少数说：具体判断说。其观点是，第一，掐卡行为不构成盗窃罪。理由是，由于卡内资金是乙的赃款，甲的掐卡行为导致乙无法转移赃款，客观上帮助了司法机关的追赃活动，不是危害行为。第二，掐卡后转账，不构成犯罪。理由是，掐卡行为已经导致甲占有控制了卡内资金。然后再转账，是进一步占有控制该资金。第三，掐卡后取款，例如，在银行ATM机或银行柜台取现金。对此，第一种意见认为（多数意见），甲对银行构成盗窃罪或诈骗罪，犯罪对象是银行的现金，甲没有取款权限。第二种意见认为（少数意见），甲对银行不构成犯罪，因为甲有取款权限。这种具体判断说是张明楷教授的观点，参见张明楷：《供卡人掐卡、取款的行为性质》，载《法学评论》2024年第1期，第10页。该观点未得到学界和司法文件的认可，因此是少数说。}

[练习1］甲向乙卖卡，不知道乙用于电信网络诈骗。某日，甲得知卡里多了1万元，才得知是乙诈骗所得赃款。甲先掐卡（挂失卡），然后转账、取款。甲出售卡的行为不构成诈骗罪的帮助犯、帮助信息网络犯罪活动罪。甲的掐卡、转账、取款行为，分析结论与上文相同。\footnote{甲的掐卡、转账、取款行为不构成掩饰、隐瞒犯罪所得罪（多数说）。理由是，成立掩饰、隐瞒犯罪所得罪，要求不违背上游犯罪人的意思。如果恶意侵害上游犯罪人的犯罪所得，则有可能构成财产犯罪。例如，甲窃得乙的财物，丙从甲处窃得该财物，丙构成盗窃罪。}

[练习2］乙诈骗丙，使丙将资金汇入乙的银行卡账户。乙向甲告知真相，让甲到银行柜台取走该现金（或者让甲在ATM机取走该现金）。乙对丙构成诈骗罪既遂。第一，多数说认为，甲对银行构成诈骗罪（或盗窃罪），同时构成掩饰、隐瞒犯罪所得罪，想象竞合，择一重罪论处。第二，少数说认为，甲不构成诈骗罪，仅构成掩饰、隐瞒犯罪所得罪。（2019年试题）

\InsertDiagram{figures/manual_table_p311_1.tex}

% Source PDF physical page 312; see audit/source_coverage.json.
[总结］“黑吃黑”

（1）“黑吃黑”，有罪的情形：对甲非法占有的财物，乙非法转移为自己占有，乙有罪。刑法保护甲对财物的非法占有。刑法禁止这种“黑吃黑”。\CJKunderwave{非法占有能够对抗第三人}。（多数说）

第一种类型：甲非法占有的赃物。例如，甲偷到丙的摩托车，乙从甲家偷走。乙构成盗窃罪。刑法保护甲对赃物的非法占有。

第二种类型：甲非法占有的赃款。例如，甲诈骗到丙的钱款，存在存折里。乙非法转移为自己占有。乙构成盗窃罪。刑法保护甲对赃款的非法占有。\footnote{还有一种原理相同的类型：共同犯罪中的赃物、赃款。例如，甲乙共同盗窃丙的摩托车，甲是实行犯，乙是帮助犯。由甲负责保管占有，乙没有事实占有该车。乙从甲家偷走该车。乙构成盗窃罪。刑法保护甲对赃物的非法占有。可能有人问，乙构成盗窃罪帮助犯既遂，车也是乙的赃物，乙为何还构成盗窃罪？这是因为，乙虽然构成盗窃罪既遂，但这只是指，乙要对甲的既遂结果负责，并不意味着乙在现实占有犯罪所得。盗窃罪看的是，将他人占有的财物非法转移为自己占有。}

第三种类型：甲非法占有违禁品（毒品、淫秽物品、假币、枪支等）。例如，甲非法持有毒品。乙从甲家偷走。乙构成盗窃罪。刑法保护甲对违禁品的非法占有。

（2）“黑吃黑”，无罪的情形：甲没有非法占有财物，甲对财物仅有非法请求权，乙侵犯该非法请求权，乙无罪。刑法不禁止这种“黑吃黑”。（多数说）

第一种类型：甲仅有非法债权。

例1（骗免资案），卖淫女为客服务完后，客用假币支付资。客不构成诈骗罪，因为法秩序不认可卖淫女的非法债权。客构成使用假币罪。

例2（欺骗分赃案），甲乙共同盗窃到丙的手机，二人约定销赃款平分，由乙负责销赃。乙销赃获得赃款1万元。乙对甲谎称手机仅卖得4000元，分给甲2000元。乙对甲不构成诈骗罪，因为法秩序不认可甲对销赃款的请求权。

第二种类型：基于不法原因的委托物。（侵占罪一节有详细讲解）

例1，委托给付的财物。甲欲向官员行贿，将财物委托给乙，让乙转交给官员，乙却将财物据为已有，拒不退还。乙不构成侵占罪，因为此时财物属于贿赂款，甲的返还请求权是非法的，不受法秩序认可。

例2，委托保管的财物。甲将其盗窃的摩托车委托给乙保管，告知是偷来的。乙窝藏后，并且拒不返还。乙构成掩饰、隐瞒犯罪所得罪，但不构成侵占罪，因为甲对赃物的返还请求权是非法的，不受法秩序认可。

\subsubsection*{（三）错误汇款}

模型：乙本想给丙汇款1万元，不慎打到了甲的卡里。甲占有卡里的这笔资金，属于不当得利，有返还义务。

1.甲针对乙。由于资金是种类物，甲花掉这笔钱（如购物），不构成侵占罪。但是，甲拒不返还，则构成侵占罪。

2.甲针对银行。甲从银行ATM机上取走这1万元现金。

（1）多数说认为，\footnote{参见张明楷：《领取无正当原因汇款的行为性质》，载《法学》2020年第11期，第19页。}甲构成盗窃罪，行为对象是银行的现金。理由是，甲不是该资金

% Source PDF physical page 313; see audit/source_coverage.json.
的所有权人，没有取款权限。\footnote{罪数上，甲如果先是拒不返还构成侵占罪，后取现金构成盗窃罪，则对先后两罪，可以重罪吸收轻罪，不并罚，因为甲最终只获得一份好处。如果甲尚无拒不返还行为，直接在ATM机取现金，则这个行为只构成盗窃罪。这个行为尚不构成侵占罪，因为对于资金这种种类物，只有拒不返还才构成侵占罪。}

（2）少数说认为，甲不构成盗窃罪。理由是，甲是该资金的法律占有人，甲有取款权限。

引申情形：甲在银行欺骗柜台员取走这1万元现金。多数说认为，甲构成诈骗罪。少数说认为，甲不构成诈骗罪。

\subsection*{三、电力、电话、网络问题}

\subsubsection*{（一）电力与电费}

例1，甲在邻居的电线上偷接一根电线，供自已使用，属于盗窃电力，构成盗窃罪。

例2，甲家的用电方式是，先用电，月底根据电表计量缴纳电费。甲在使用过程中，在电表上安装某种非法设备，导致电表运转很慢，实际使用了5000度电，只显示1000度电。电力公司收取电费时，甲向收费员展示电表计数，收费员以为甲只使用了1000度电，收取了1000度电费。收费员基于认识错误而处分财产性利益，其处分行为表现为不作为，亦即不行使权利。甲构成诈骗罪，属于骗免电费，诈骗财产性利益。（2021年试题）

例3，甲家的用电方式是，需要先给电卡里充值，然后使用电力，如果电卡里没余额了，就会断电。甲的电卡里已经没有余额，但是甲采取技术手段，能够继续使用电力。甲的行为属于盗窃电力，构成盗窃罪。

\subsubsection*{（二）电话、网络（司法解释总结）}

1.盗打电话、盗用网络

第一，以牟利为目的，盗接他人通信线路、复制他人电信码号或者明知是盗接、复制的电信设备、设施而使用的，定盗窃罪（第265条）。盗窃罪的数额根据给他人造成的资费损失数额来计算。

第二，盗用他人上网账号、密码上网，造成他人资费损失数额较大的，定盗窃罪。盗窃罪的数额根据给他人造成的资费损失数额来计算。\footnote{提示：以上两种盗窃（盗用）的对象是通讯信号、网络信号，而非电话费、上网费等财产性利益（债权）。类似的，盗接他人电线，拉到自己家里使用，盗窃的是电力，而非电费（债权）。只是“通讯信号”“网络信号”“电力”等无形财物的价值在计算时根据资费来计算。}

2.非法充值

第一，将电信卡非法充值后使用，造成电信资费损失数额较大的，定盗窃罪。盗窃罪的数额根据给他人造成的资费损失数额来计算。

第二，明知是非法充值的电信卡而使用或者购买而使用，造成资费损失数额较大的，定盗窃罪。盗窃罪的数额根据给他人造成的资费损失数额来计算。

3.诈骗电信工作人员

以虚假、冒用的身份证件，带着不想交费的目的，欺骗电信公司工作人员，办理人网手续（先使用后交费），使用电话或网络，造成资费损失数额较大的，定诈骗罪。

% Source PDF physical page 314; see audit/source_coverage.json.
\section{财产犯罪的罪数联系}

\subsection*{一、销赃行为}

1.盗窃与销赃。例如（卖树案），甲见外地商贩乙来本村购买木材，便指着邻居丙（丙不在家）的树木说：“这些树是我的，3万元卖给你。”乙信以为真，交给甲3万元。甲便说：“这些树归你了，你自已伐后运走。”乙便伐倒运走。甲的一个卖树行为，一方面对邻居构成盗窃罪，另一方面对乙构成诈骗罪（多数说），想象竞合，择一重罪论处。（2014年第60题，2015年第18题，2016年第17题，2019年试题）

2.诈骗与销赃（两头骗）。例如，甲诈骗到乙的摩托车，然后向丙谎称是自己的摩托车，卖给了丙。甲先对乙构成诈骗罪，后对丙构成诈骗罪（多数说）。两个诈骗罪并罚。

3.侵占与销赃。

例1，甲正常租用乙的汽车，到期乙向甲索还时，甲拒不归还，然后又伪造虚假材料，谎称汽车是自已所有，以正常行情价卖给丙。甲拒不归还，对乙构成侵占罪。甲的卖车行为，对丙构成诈骗罪，两罪并罚。

例2，二重买卖（2017年第18题）。甲销售汽车给乙，乙分期付款购买，约定车款付清前汽车由乙使用，但所有权归甲，待付清车款，所有权再转移给乙。乙在车款付清前，在使用期间，对丙谎称是自已所有的车，将车卖给丙。乙只有一个处分行为，一方面对甲构成侵占罪，一方面对丙构成诈骗罪，想象竞合，择一重罪论处。\footnote{引申练习：甲先将自己的房子卖给乙，乙交了钱，但尚未办理过户。甲又高价卖给丙（丙不知甲乙事宜），并过户给丙。甲对乙构成违约，不构成犯罪。}

\subsection*{二、免除返还义务}

1.先侵占后诈骗。例如，乙将财物交给甲，让甲拿回家保管。后甲将财物变卖。当乙要求返还时，甲欺骗乙：“被盗了（或被抢了）”，乙免除了其返还义务。甲的第一个行为构成侵占罪，第二个行为构成诈骗罪，诈骗财产性利益。首先，对两罪不需要并罚，因为乙最终遭受的财产损失只有一份，甲最终获得的好处只有一份。其次，重罪吸收轻罪，定诈骗罪。\footnote{少数说认为，只定前行为的罪名，也即侵占罪，将后行为视为不可罚的事后行为。}（2022年试题）

2.先侵占后杀害。例如（2016年主观题），甲接受乙的委托，保管乙的财物。甲将财物变卖。当乙要求归还时，甲为了免除返还义务，杀害了乙。甲的第一个行为构成侵占罪。甲的第二个行为属于抢劫财产性利益，也即通过暴力手段免除了返还义务，构成抢劫罪（故意）致人死亡，吸收前面的侵占罪。\footnote{少数说认为，第二个行为不构成抢劫罪，只构成故意杀人罪，然后与侵占罪并罚。}

3.先诈骗后抢劫。例如（2017年第8题），甲骗到乙的名画，乙日后发现，要求返还，甲暴力威胁，迫使乙免除返还义务。甲的第一个行为构成诈骗罪，第二个行为构成抢劫罪，抢劫财产性利益。对两罪不需要数罪并罚，可以重罪吸收轻罪，因为乙最终遭受的财产损失只有一份，甲最终获得的好处只有一份。\footnote{少数说认为，第二个行为不构成抢劫罪和其他犯罪。}

% Source PDF physical page 315; see audit/source_coverage.json.
\subsection*{三、先偷后骗}

例1，甲将摩托车借给朋友乙使用，后又从乙家偷回来。次日，乙告知甲摩托车丢失，并提出赔偿4000元。甲隐瞒真相，接受赔偿。（2003年第10题）

例2，甲从超市盗窃一部手机，又要求“退货”，换取到现金。

例3，甲从游戏厅的游戏机中盗窃游戏币，又“兑换”现金。

上述案例，甲的第一个行为构成盗窃罪，第二个行为构成诈骗罪。\footnote{对于例2、3中甲的第二个行为，少数说认为不构成诈骗罪，理由是，该行为属于等价交换，相对方没有财产损失。}由于甲最终只获得一份好处，被害人最终只遭受一份财产损失，所以没必要数罪并罚，可重罪吸收轻罪。\footnote{少数说认为，两个行为是独立行为，且两个行为的对象不具有同一性，前者是财物，后者是现金或另一财物，应数罪并罚。}

